\documentclass[10pt]{article}
% %%%%%%%%%%%%%%%%%%%%%%%%%% Table set up

\usepackage{tabu}
%%%%%%%%%%%%%%%%%%%%%%%%%%%%%%%%%%%%% FORMAT
%%%%%%%%%%%%%%%%%%%%%%%%%%% footer and header
\usepackage{lastpage}

\usepackage{fancyhdr}
\usepackage{xparse}
\pagestyle{fancy}
%%%%%%%%%%%%%%%%%%%%%%%%%%% Headers and footers

\fancyhead{} % clear all header fields
\fancyfoot{} % clear all footer fields


\rfoot{Page \thepage\ of \pageref{LastPage}}
% Define the footer separation line
\renewcommand{\footrulewidth}{0.4pt}
\addtolength{\headwidth}{4.8cm}
\usepackage[headheight=80pt,tmargin=85pt,headsep=25pt, footskip=25pt]{geometry}

\geometry{a4paper, margin=2cm}

\usepackage[numbers]{natbib}

%%%%%%%%%%%%%%%%%%%%%%%%%%% Modify the description list
\usepackage{nameref}

\makeatletter
\let\orgdescriptionlabel\descriptionlabel
\renewcommand*{\descriptionlabel}[1]{%
  \let\orglabel\label
  \let\label\@gobble
  \phantomsection
  \edef\@currentlabel{#1\unskip}%
  \let\label\orglabel
  \orgdescriptionlabel{#1}%
}
\makeatother
%%%%%%%%%%%%%%%%%%%%%%%%%%% Font colour

\usepackage[utf8]{inputenc}
\usepackage[english]{babel}
\usepackage[dvipsnames]{xcolor}

\usepackage{listings}
\usepackage{color}
% %%%%%%%%%%%%%%%%%%%%%%%%% Tikz packages:
\usepackage{tikz}
\usepackage{tikz-layers}
\usepackage{array}
\newcolumntype{C}[1]{>{\centering\arraybackslash}p{#1}}
%
% the positioning-plus library (is on TeX.sx),
% the node-families library (is on TeX.sx),
% the fit library (loaded by positioning-plus),
% the backgrounds library to draw stuff behind other stuff,
% the calc library for some funny coordinate calculations and the
% shapes.geometric library for the ellipse shape

\usetikzlibrary{shapes.geometric,automata,positioning,arrows,shadows,backgrounds,calc}

\definecolor{frameColor}{RGB}{128,128,128}
\definecolor{initialStateColor}{RGB}{10,100,10}
\definecolor{impactStateColor}{RGB}{0,255,0}
\definecolor{admittanceStateColor}{RGB}{0,0,255}

\definecolor{resetStateColor}{RGB}{10,100,10}

\definecolor{stateColor}{RGB}{51,204,204}

\definecolor{controllerColor}{RGB}{10,10,255}
\definecolor{plannerColor}{RGB}{10,150,150}
\definecolor{estimatorColor}{RGB}{10,200,10}
\definecolor{robotColor}{RGB}{255,0,0}
\definecolor{modelColor}{RGB}{180,100,10}

%%% Enumeration
\usepackage{enumitem}

%%%%%%%%%%%%%%%%%%%%%%%%%%% Table
\usepackage{multirow}


%%%%%%%%%%%%%%%%%%%%%%%%%%% Adjust font size
\usepackage{setspace}

%%%%%%%%%%%%%%%%%%%%%%%%%%%
%%%%%%%%%%%%%%%%%%%%%%%%%%%%%%%%%%%%%% MATH:
%%%%%%%%%%%%%%%%%%%%%%%%% AMS packages:
\usepackage{mathtools}
\usepackage{amsmath, amsthm, amsfonts}
\usepackage{amssymb}
\usepackage{graphicx}
\usepackage[font=normal,labelfont=bf]{caption}
\usepackage{subcaption}
\usepackage{url}
\usepackage[colorlinks=true,hyperindex=true,bookmarksnumbered]{hyperref} % Turn on "bookmarks" and "hyperindex" to generate outlines.

%%%%%%%%%%%%%%%%%%%%%%%%% Answer and reply
\usepackage{thmtools}
\declaretheoremstyle[headfont=\bfseries,
bodyfont=\normalfont,
numbered=no]{normalhead}
\declaretheorem[style=normalhead]{Reply}

\theoremstyle{plain}
\newtheorem*{Reply*}{Reply}

\newtheorem*{AuthorReply*}{AuthorReply}

\declaretheoremstyle[
  headfont=\color{blue}\normalfont\bfseries,
  bodyfont=\color{red}\normalfont\itshape,
  numbered=no
]{quoteColor}

\declaretheorem[style=quoteColor]{Quote}
\newtheorem*{Quote*}{Quote}

\declaretheoremstyle[
  headfont=\color{black}\normalfont\bfseries,
  bodyfont=\color{blue}\normalfont,
  numbered=no
]{reviewColor}
\declaretheorem[style=reviewColor]{review}
\newtheorem*{review*}{review}

\theoremstyle{definition}

\theoremstyle{remark}
% -----------------------------------------------------------------


%%%%%%%%%%%%%%%%%%%%%%%%%%%%%%%%%%%%%%%%%%%%%%%%%%%%%%%%%%%%%%%%%%%%%%%%%%%%%%%%
%% 1. PACKAGES AND ENVIRONMENT SETUP
%%%%%%%%%%%%%%%%%%%%%%%%%%%%%%%%%%%%%%%%%%%%%%%%%%%%%%%%%%%%%%%%%%%%%%%%%%%%%%%%

\usepackage{xparse} % For enhanced command definitions, coupled with transposeRow
\usepackage{xstring} % For string parsing, e.g., \IfStrEq{#1}{dot}{}{}
\usepackage{ifthen} % For semantic selectors used by frame and notation macros.
\usepackage{mathtools} % Core display math and \DeclareMathOperator for legacy direct inputs.
\usepackage{amsfonts,amssymb}
\usepackage{enumitem} % Structured enumerate labels used throughout the reports.

\ExplSyntaxOn
\NewDocumentCommand{\yuquanIfBlankTF}{mmm}{%
  \tl_if_blank:nTF {#1} {#2} {#3}%
}
\ExplSyntaxOff

\usepackage{mdframed}
\usepackage{natbib} % For \citet command

\newmdenv[
shadow=true,
shadowsize=6pt,
linewidth=0pt,
backgroundcolor=gray!20,
shadowcolor=black!40,
roundcorner=10pt
]{shadowed}

\usepackage{makecell}

%%%%%%%%%%%%%%%%%%%%%%%%%%% For  the cross mark
\usepackage[utf8]{inputenc}
\usepackage[english]{babel}

\usepackage{color}
% %%%%%%%%%%%%%%%%%%%%%%%%% Tikz packages:
\usepackage{tikz}
\usepackage{tikz-layers}
\usepackage{array}
%
% the positioning-plus library (is on TeX.sx),
% the node-families library (is on TeX.sx),
% the fit library (loaded by positioning-plus),
% the backgrounds library to draw stuff behind other stuff,
% the calc library for some funny coordinate calculations and the
% shapes.geometric library for the ellipse shape

\usetikzlibrary{shapes.geometric,automata,positioning,arrows,shadows,backgrounds,calc,fit}

\usepackage{setspace}




%%%%%%%%%%%%%%%%%%%%%%%%%% MATLAB setup

\usepackage{listings}
\lstset{
  language=Matlab,                		% choose the language of the code
  %	basicstyle=10pt,       				% the size of the fonts that are used for the code
  numbers=left,                  			% where to put the line-numbers
  numberstyle=\footnotesize,      		% the size of the fonts that are used for the line-numbers
  stepnumber=1,                   			% the step between two line-numbers. If it's 1 each line will be numbered
  numbersep=5pt,                  		% how far the line-numbers are from the code
  %	backgroundcolor=\color{white},  	% choose the background color. You must add \usepackage{color}
  showspaces=false,               		% show spaces adding particular underscores
  showstringspaces=false,         		% underline spaces within strings
  showtabs=false,                 			% show tabs within strings adding particular underscores
  %	frame=single,	                			% adds a frame around the code
  %	tabsize=2,                				% sets default tabsize to 2 spaces
  %	captionpos=b,                   			% sets the caption-position to bottom
  breaklines=true,                			% sets automatic line breaking
  breakatwhitespace=false,        		% sets if automatic breaks should only happen at whitespace
  escapeinside={\%*}{*)}          		% if you want to add a comment within your code
}


%%%%%%%%%%%%%%%%%%%%%%%%%% Self-defined commands:


\usepackage{xspace}
\lstset{frame=tb,
  language=C++,
  backgroundcolor=\color{black!5}, % set backgroundcolor
  basicstyle=\footnotesize,% basic font setting
  aboveskip=3mm,
  belowskip=3mm,
  showstringspaces=false,
  columns=flexible,
  numbers=none,
  numberstyle=\tiny\color{gray},
  keywordstyle=\color{blue},
  commentstyle=\color{dkgreen},
  stringstyle=\color{mauve},
  breaklines=true,
  breakatwhitespace=true,
  tabsize=3
}


\usepackage{soul}

\usepackage{pifont}% http://ctan.org/pkg/pifont

\usepackage{stackengine}
\newcommand\barbelow[1]{\stackunder[1.2pt]{$#1$}{\rule{.8ex}{.075ex}}}

% Algorithms.  Use the standard algorithmicx packages when available.  A
% small compatibility implementation keeps the reports buildable on minimal
% TeX installations; it supports the constructs used in this repository.
\IfFileExists{algorithm.sty}{%
  \usepackage{algorithm}% http://ctan.org/pkg/algorithms
  \usepackage{algorithmicx}
  \usepackage{algpseudocode}% http://ctan.org/pkg/algorithmicx
}{%
  \usepackage{float}
  \floatstyle{ruled}
  \newfloat{algorithm}{tbp}{loa}
  \floatname{algorithm}{Algorithm}
  \newlength{\algorithmicindent}
  \setlength{\algorithmicindent}{1.5em}
  \NewDocumentEnvironment{algorithmic}{O{}}{%
    \begin{list}{}{%
      \setlength{\leftmargin}{1.5em}%
      \setlength{\itemsep}{0.15em}%
      \setlength{\parsep}{0pt}%
    }%
  }{\end{list}}
  \newcommand{\State}{\item}
  \newcommand{\Procedure}[2]{\item\textbf{procedure} ##1\yuquanIfBlankTF{##2}{}{\ (##2)}}
  \newcommand{\EndProcedure}{\item\textbf{end procedure}}
  \newcommand{\For}[1]{\item\textbf{for} ##1 \textbf{do}}
  \newcommand{\EndFor}{\item\textbf{end for}}
  \newcommand{\If}[1]{\item\textbf{if} ##1 \textbf{then}}
  \newcommand{\Else}[1]{\item\textbf{else}\yuquanIfBlankTF{##1}{}{ ##1}}
  \newcommand{\EndIf}{\item\textbf{end if}}
  \newcommand{\While}[1]{\item\textbf{while} ##1 \textbf{do}}
  \newcommand{\EndWhile}{\item\textbf{end while}}
  \newcommand{\Comment}[1]{\hfill$\triangleright$~##1}
}

%%% predecessor:
\usepackage{dsfont}

%%%%%%%%%%%%%%%%%%%%%%%%%%%%%%%%%%%%%%%%%%%%%%%%%%%%%%%%%%%%%%%%%%%%%%%%%%%%%%%%
%% 2. COLOR DEFINITIONS
%%%%%%%%%%%%%%%%%%%%%%%%%%%%%%%%%%%%%%%%%%%%%%%%%%%%%%%%%%%%%%%%%%%%%%%%%%%%%%%%

\definecolor{frameColor}{RGB}{128,128,128}
\definecolor{initialStateColor}{RGB}{10,100,10}
\definecolor{impactStateColor}{RGB}{0,255,0}
\definecolor{admittanceStateColor}{RGB}{0,0,255}

\definecolor{resetStateColor}{RGB}{10,100,10}

\definecolor{stateColor}{RGB}{51,204,204}

\definecolor{controllerColor}{RGB}{10,10,255}
\definecolor{plannerColor}{RGB}{10,150,150}
\definecolor{estimatorColor}{RGB}{10,200,10}
\definecolor{robotColor}{RGB}{255,0,0}
\definecolor{modelColor}{RGB}{180,100,10}

%%% Enumeration

%%%%%%%%%%%%%%%%%%%%%%%%%%% Adjust font size
\definecolor{dkgreen}{rgb}{0,0.6,0}
\definecolor{gray}{rgb}{0.5,0.5,0.5}
\definecolor{mauve}{rgb}{0.58,0,0.82}

\definecolor{svaColor}{RGB}{7, 160, 2}
\definecolor{uncertainColor}{RGB}{255, 0, 0}

\definecolor{spatialVelColor}{RGB}{189, 38, 96}
\definecolor{geometricColor}{RGB}{66, 126, 245}
\definecolor{bodyVelColor}{RGB}{13, 191, 191}

\definecolor{velColor}{RGB}{50, 205, 50}

\definecolor{textColor}{RGB}{0, 11, 178}


\colorlet{svaColorVariable}{svaColor}
\colorlet{geometricColorVariable}{geometricColor}

% %%%%%%%%%%%%%%%%%%%%%%%%% comment
\definecolor{csrStateColor}{RGB}{255,255,153}
\definecolor{scrStateColor}{RGB}{153,204,0}
\definecolor{crStateColor}{RGB}{153,204,155}

%%%%%%%%%%%%%%%%%%%%%%%%%% DDP commands:

% Open-loop gain for computing the optimal control action

%%%%%%%%%%%%%%%%%%%%%%%%%%%%%%%%%%%%%%%%%%%%%%%%%%%%%%%%%%%%%%%%%%%%%%%%%%%%%%%%
%% 3. FORMATTING AND PRESENTATION
%%%%%%%%%%%%%%%%%%%%%%%%%%%%%%%%%%%%%%%%%%%%%%%%%%%%%%%%%%%%%%%%%%%%%%%%%%%%%%%%

\newcommand{\shadowRemark}[1]{
\begin{shadowed}
\begin{remark}
{#1}
\end{remark}
\end{shadowed}
}

\newcommand{\shadowProblem}[2]{
\begin{shadowed}
\infoBlock{Problem:} \hlBlock{#1}
{#2}
        \end{shadowed}
}

%%%%%%%%%%%%%%%%%%%%%%%%%%% Splitting lines within a table cell
\newcommand{\crossmark}{
  \ding{55}
}


%%%%%%%%%%%%%%%%%%%%%%%%%%% Font colour
\newcommand{\mcrtc}{mc$\_$rtc\xspace}
\newcommand{\sota}{state-of-the-art\xspace}

\newcommand{\poincare}{Poincar\'e\xspace}

%%%%%%%%%%%%%%%%%%%%%%%%%% Units:
\newcommand{\ith}{$i$-th\xspace}
%%%%%%%%%%%%%%%%%%%%%%%%%% Units:
\newcommand{\highlight}[1]{%
  \colorbox{gray!20}{$\displaystyle#1$}}

\newcommand{\highlightblue}[1]{%
  \colorbox{blue!20}{$\displaystyle#1$}}
\newcommand{\highlightgreen}[1]{%
  \colorbox{green!20}{$\displaystyle#1$}}


%%%%%%%%%%%%%%%%%%%%%%%%%% Bounds on tasks
\newcommand{\threshold}[1]{\Theta_{#1}}

%%%%%%%%%%%%%%%%%%%%%%%%%% Control Lyapunov function

%%%%%%%%%%%%%%%%%%%%%%%%%% Control Barrier function


%%%%%%%%%%%%%%%%%%%%%%%%%% Optimal Control
% Estimated quantity
\NewDocumentCommand{\cpp}{v}{%
\texttt{\textcolor{blue}{#1}}%
}

%%%%%%%%%%%%%%%%%%%%%%%%%%%%% New reference command:
\NewDocumentCommand{\secRef}{m}{%
Sec.~\ref{#1}}

\NewDocumentCommand{\remarkRef}{m}{%
Remark.~\ref{#1}}

\NewDocumentCommand{\tableRef}{m}{%
Table.~\ref{#1}}

\NewDocumentCommand{\algRef}{m}{%
Algorithm.~\ref{#1}}

\NewDocumentCommand{\lemmaRef}{m}{%
Lemma.~\ref{#1}}

\NewDocumentCommand{\theoremRef}{m}{%
Theorem.~\ref{#1}}

\NewDocumentCommand{\figRef}{m}{%
  Fig.~\ref{#1}}

\NewDocumentCommand{\appRef}{m}{%
  Appendix~\ref{#1}}

\NewDocumentCommand{\lineRef}{m}{%
line.~\ref{#1}}

\NewDocumentCommand\orderedTwoS{mm}%
  {$<$#1,#2$>$}
\NewDocumentCommand\orderedThreeS{mmm}%
  {$<$#1,#2,#3$>$}

\newif\ifBlockComment
\BlockCommentfalse %\BlockCommenttrue
% %%%%%%%%%%%%%%%%%%%%%%%%%


\newcommand{\comment}[1]{{\bf \colorbox{teal}{\textcolor{white}{#1}}}}
\newcommand{\fix}[1]{{\bf \colorbox{purple}{\textcolor{white}{#1}}}}
% %%%%%%%%%%%%%%%%%%%%%%%%% Highlight with colors
\sethlcolor{teal}
\newcommand{\hlBlock}[1]{{\sethlcolor{teal}\color{white}\hl{#1}}}
\newcommand{\errorBlock}[1]{{\sethlcolor{violet}\color{white}\hl{#1}}}
% Question block
\newcommand{\qtBlock}[1]{{\sethlcolor{violet}\color{white}\hl{#1}}}
\newcommand{\infoBlock}[1]{{\sethlcolor{olive}\color{white}\hl{#1}}}
\newcommand{\bashBlock}[1]{{\sethlcolor{black}\color{white}\hl{#1}}}

\newcommand{\tBlock}[1]{{\bf \colorbox{teal}{\textcolor{white}{#1}}}}
\newcommand{\oBlock}[1]{{\bf \colorbox{olive}{\textcolor{white}{#1}}}}

\newcommand{\bBlock}[1]{{\sethlcolor{brown}\color{white}\hl{#1}}}


\newcommand{\fixedwidth}[1]{
{\parbox{\dimexpr\linewidth-20\fboxsep}{#1}}
}
\newcommand{\wrongBlock}[1]{{\color{brown}\ding{55}}~~{\bf \colorbox{brown}{\textcolor{white}{#1}}}}
\newcommand{\correctBlock}[1]{
{{\color{olive}\ding{51}}~{\bf \colorbox{teal}{\textcolor{white}{
	\fixedwidth{#1}
	}}}}
}
% Delay the lightweight fallback until every package has been loaded.  This
% lets reports opt into todonotes after inputting this legacy configuration
% without triggering a duplicate \todo definition.
\AtBeginDocument{%
  \providecommand{\todo}[1]{To do:~
    \textcolor{bodyVelColor}{#1}%
  }%
}


% The critical frequency
\newcommand{\cmark}{\ding{51}}%
\newcommand{\xmark}{\ding{55}}%


% The motion constraint Jacobian

%%%%%%%%%%%%%%%%%%%%%%%%%%%%%%%%%%%%%%%%%%%%%%%%%%%%%%%%%%%%%%%%%%%%%%%%%%%%%%%%
%% 4. UNITS
%%%%%%%%%%%%%%%%%%%%%%%%%%%%%%%%%%%%%%%%%%%%%%%%%%%%%%%%%%%%%%%%%%%%%%%%%%%%%%%%

\newcommand{\unitHz}{~Hz\xspace}
\newcommand{\unitMs}{~ms\xspace}
\newcommand{\unitTime}{~s\xspace}
\newcommand{\unitDegree}{~deg.\xspace}
\newcommand{\unitMass}{~kg\xspace}
\newcommand{\unitMeter}{~m\xspace}
\newcommand{\unitHertz}{~N/m$^{3/2}$\xspace}
\newcommand{\unitGPa}{~GPa\xspace}


\newcommand{\unitPosTS}{~m\xspace}
\newcommand{\unitPosJS}{~rad\xspace}
\newcommand{\unitVelTS}{~m/s\xspace}
\newcommand{\unitVelJS}{~rad/s\xspace}
\newcommand{\unitAccTS}{~m/s$^2$\xspace}
\newcommand{\unitAccVel}{~rad/s$^2$\xspace}
\newcommand{\unitForce}{~N\xspace}

%%%%%%%%%%%%%%%%%%%%%%%%%%%
\newcommand{\unitImpulse}{~N$\cdot$s\xspace}

%%%%%%%%%%%%%%%%%%%%%%%%%% LQR commands:

%%%%%%%%%%%%%%%%%%%%%%%%%%%%%%%%%%%%%%%%%%%%%%%%%%%%%%%%%%%%%%%%%%%%%%%%%%%%%%%%
%% 5. BASIC MATH NOTATION
%%%%%%%%%%%%%%%%%%%%%%%%%%%%%%%%%%%%%%%%%%%%%%%%%%%%%%%%%%%%%%%%%%%%%%%%%%%%%%%%

\newcommand{\setDef}[2]{
  \defeq \{#1 \mid #2 \}
}

\newcommand{\quickEq}[2]{
  \begin{equation}
    \label{#1}
    {#2}
  \end{equation}
}

\newcommand{\quickBM}[1]{
    \begin{bmatrix}
    {#1}
    \end{bmatrix}
}

\newcommand{\quickAd}[1]{
  \begin{aligned}
    {#1}
  \end{aligned}
}

\newcommand{\de}[1]{
  \frac{d}{dt}(#1)
}
%%%%%%%%%%%%%%%%%%%%%%%%%% Collision observer


\newcommand{\setSymbol}{\mathcal{X}}

\newcommand{\feasibleSymbol}{\mathcal{F}}
\newcommand{\viableSymbol}{\nu}

\newcommand{\feasibleSet}[1]{\feasibleSymbol_{#1}}
\newcommand{\viableSet}[1]{\viableSymbol_{#1}}

\newcommand{\set}[1]{\setSymbol_{#1}}
% Interval
\newcommand{\interval}[2]{[\ #1,~#2 ]\ }



% %%%%%%%%%%%%%%%%%%%%%%%%% COLOUR:

\newcommand{\quadraticPN}[4]{
#1\pTwo{#4}  #2{#4}   #3 = 0
}
%%%%%%%%%%%%%%%%%%%%%%%%% Momenta
\newcommand{\weightMatrix}[1]{\mathbb{#1}}
\newcommand{\weightmatrixH}{\mathbf{H}}
\newcommand{\weightmatrixW}{\mathbf{W}}
\newcommand{\weightmatrixQ}{\mathbf{Q}}


\newcommand{\matrixA}{\mathbf{A}}

\IfFileExists{accents.sty}{%
  \usepackage{accents}
  \newcommand{\ubar}[1]{\underaccent{\bar}{##1}}
}{%
  % A conventional underline is preferable to making the entire package fail
  % when the optional accents package is not installed.
  \newcommand{\ubar}[1]{\underline{##1}}
}
\newcommand{\boundSymbol}{l}

\newcommand{\applyBoundSuperscript}[1]{%
    \IfStrEqCase{#1}{%
        {min}{\text{min}}%
        {max}{\text{max}}%
    }[] % "default" string
}

\newcommand{\applyBoundSubscript}[1]{%
        #1
}


\NewDocumentCommand{\bound}{O{}O{}}{%
  \ensuremath{%
    \boundSymbol
    \yuquanIfBlankTF{#2}{}{^{\applyBoundSubscript{#2}}}%
    \yuquanIfBlankTF{#1}{}{_{\applyBoundSuperscript{#1}}}%
  }
}


\newcommand{\lowerBound}[1]{{#1}_{\text{min}}}
\newcommand{\upperBound}[1]{{#1}_{\text{max}}}
%%%%%%%%%%%%%%%%%%%%%%%%%%%%%%%%%%%%%%%%%%%%%%%%%%%%%%%%%%%%%%%%% Bound

\newcommand{\stateBound}[1]{
        \lowerBound{#1} \leq #1 \leq \upperBound{#1}
}
\newcommand{\stateAbsBound}[1]{
        | #1 |  \leq \upperBound{#1}
}
\newcommand{\stateBoundConstraint}[2]{
        \lowerBound{#1} \leq #2 \leq \upperBound{#1}
}
\newcommand{\doubleBound}[1]{\ubar{\bar{#1}}}


\newcommand{\defeq}{\vcentcolon=}
\newcommand{\eqdef}{=\vcentcolon}


% Operators
% New operators must defined as such to have them typeset
% correctly. As an example we define the Jacobian:
% -----------------------------------------------------------------
\DeclareMathOperator{\Jac}{Jac}


% Shortcuts.
% One can define new commands to shorten frequently used
% constructions. As an example, this defines the R and Z used
% for the real and integer numbers.
% -----------------------------------------------------------------
\def\RR{\mathbb{R}}
\def\ZZ{\mathbb{Z}}

% Integrate over time:  #1 starting time, #2 finishing time, #3 quantity
\newcommand{\timeIntegral}[3]{
\int^{#2}_{#1} #3 dt
}

\newcommand{\RRv}[1]{\mathbb{R}^{#1}}
\newcommand{\RRm}[2]{\mathbb{R}^{#1 \times #2}}

%%%%%%%%%%%%%%%%%%%%%%%%%%% Impact-handling

\DeclareMathOperator*{\argmin}{arg\,min} % thin space, limits underneath in displays
\DeclareMathOperator*{\lexmin}{lex\,min} % Lexicographic-order minimization

%%%%%%%%%%%%%%%%%%%%%%%%%% Probabilistic robotics


%%%%%%%%%%%%%%%%%%%%%%%%%% Collision avoidance notations
% A link
\newcommand{\fpd}[2]{{#1}_{#2}}
% Second-order partial derivative #1 quantity, #2 variable, #3 second variable
\newcommand{\spd}[3]{{#1}_{#2#3}}



% The argument denotes the initial state
\newcommand{\Vee}[1]{\left(#1\right)^\vee}


%%%%%%%%% Represent the derivative operator
\newcommand{\derivative}[1]{\frac{d}{dt}(#1)}
%%% #1:function, #2:var1, #3:var2
\newcommand{\pdv}[3]{
  \frac{\partial^2 #1}{\partial #2 \partial #3}
}
%%% #1:function, #2:var1
\newcommand{\pd}[2]{
  \frac{\partial #1}{\partial #2 }
}

\newcommand{\up}[1]{\skewMatrix{#1}}

\newcommand{\liebra}[2]{
  [#1,#2]
}

% #1 Notation of the twist coordinate
\newcommand{\liebraResult}[2]{
  \Vee{\liebraDefDef{#1}{#2}}
}
%%%% #1:linear velocity #2: angular velocity #3:linear velocity #4: angular velocity
\newcommand{\seLiebraResult}[4]{
(\cross{#2}{#3} - \cross{#4}{#1}, \skewMatrix{#2}#4 - \skewMatrix{#4}#2 )
}



\newcommand{\liebraDef}[2]{
  \skewMatrixTwo{\cross{#1}{#2}}
}

\newcommand{\liebraDefDef}[2]{
  #1#2 - #2#1
}

%%%%%%%%%%% ----------------------------------------

%% atan2
% #1: y #2: x
\newcommand{\twoNorm}[1]{
  {{\left\lVert#1\right\rVert}^2}
}

\ExplSyntaxOn
\NewDocumentCommand{\applyNormSymbols}{mm}{%
  \str_case:nnF {#1}
    {
      {two}{\left\lVert#2\right\rVert^2_2}% Squared Euclidean norm.
      {square}{\left\lVert#2\right\rVert^2}% Squared unspecified norm.
      {one}{\left\lVert#2\right\rVert_1}% One norm.
      {abs}{\left\lvert#2\right\rvert}% Scalar absolute value.
      {inf}{\left\lVert#2\right\rVert_{\infty}}% Infinity norm.
    }
    {\left\lVert#2\right\rVert}% Default norm.
}
\ExplSyntaxOff

\NewDocumentCommand{\norm}{O{}m}{%
  \applyNormSymbols{#1}{#2}%
}

\newcommand{\abs}[1]{\left\vert#1\right\vert}
\newcommand{\cross}[2]{ #1 \times #2 }
\newcommand{\crossDef}[1]{ % Define the 3-dimensional cross operator for a vector
\begin{bmatrix}
0 & -{#1}_3 & {#1}_2 \\
{#1}_3 & 0 &  -{#1}_1 \\
-{#1}_2 & {#1}_1 & 0
\end{bmatrix}
}
\newcommand{\crossDefTwo}[1]{ % Define the 3-dimensional cross operator for a vector
\begin{bmatrix}
0 & -\z{#1} & \y{#1} \\
\z{#1} & 0 &  -\x{#1} \\
-\y{#1} & \x{#1} & 0
\end{bmatrix}
}



\newcommand{\dualCross}[2]{ #1 \times^* #2 }
\newcommand{\dualCrossDef}[2]{ -#1 \times^\top #2 }
\newcommand{\innerP}[2]{
  \transpose{#1}#2
}

% Computes the outer product
\newcommand{\outerP}[2]{
  #1\transpose{#2}
}



% Calculates the reciprocal product between two twists
% #1:v_1,  #2:w_1,  #3:v_2, #4:w_2.
\newcommand{\agg}[3]{
\sum^{#2}_{#1=1}{#3}
}
%%% #1:index, #2:first index, #3:last index, #3 component
\newcommand{\aggTwo}[4]{
\sum^{#3}_{#1=#2}{#4}
}

%%% #1:index, #2:last index, #3 component
\newcommand{\union}[3]{
\bigcup^{#2}_{#1=1}{#3}
}


%%%%%%%%%%% #1:row index, #2:column index, #3 index (between 1 and dof)
\newcommand{\columnTwo}[2]{
  \begin{bmatrix}
    #1&  #2
\end{bmatrix}
}

\newcommand{\vectorTwo}[2]{
  \begin{bmatrix}
    #1\\
    #2
\end{bmatrix}
}

\newcommand{\vectorThree}[3]{
  \begin{bmatrix}
    #1\\
    #2\\
    #3
\end{bmatrix}
}

\newcommand{\vectorFour}[4]{
  \begin{bmatrix}
    #1\\
    #2\\
    #3\\
    #4
\end{bmatrix}
}
\newcommand{\vectorTwoRow}[2]{
  [#1^\top, #2^\top]^\top
}

\newcommand{\vectorThreeRow}[3]{
  [{#1}^\top, {#2}^\top, {#3}^\top]^\top
}
\newcommand{\matrixTwo}[4]{
  \begin{bmatrix}
    #1 & #2\\
    #3 & #4
\end{bmatrix}
}

\newcommand{\matrixThree}[9]{
  \begin{bmatrix}
    #1 & #2 & #3\\
    #4 & #5 & #6\\
    #7 & #8 & #9
\end{bmatrix}
}

\newcommand{\order}[1]{\mathcal{O}(#1)}


\newcommand{\SE}[1]{SE{(#1)}}
\newcommand{\se}[1]{se{(#1)}}

\newcommand{\SO}[1]{SO{(#1)}}
\renewcommand{\so}[1]{so{(#1)}}
\newcommand{\soPowerTwoDef}[1]{
  \projectionDef{#1} - \twoNorm{#1} \identityMatrix
}
\newcommand{\soPowerThreeDef}[1]{
  - \twoNorm{#1} \skewMatrix{#1}
}


\newcommand{\mapDef}[2]{
  : #1 \rightarrow #2
}


% transform a twist coordinate from coordinate #1 to coordinate #2
\newcommand{\quadraticObj}[2]{\transpose{#1}#2#1}

\newcommand{\complexNum}[2]{
  (#1 ~ #2*i)
}

%% #1 upper bound, #2 middle #3 lower bound
\newcommand{\saturation}[3]{
\left\{
  \begin{array}{lr}
    #1 \\
    #2 \\
    #3
  \end{array}
\right.
}
% #1 velocity #2 mass
\newcommand{\quadratic}[2]{\frac{1}{2} #1^\top #2 #1}
\newcommand{\squareTwo}[1]{ #1^\top #1}
\newcommand{\squared}[1]{  {#1}^2 }


\ExplSyntaxOn

 % Define a sequence variable to store processed items
 \seq_new:N \trSeq

% Define the command to process and display the list
\NewDocumentCommand{\transposeRow}{m}
 {

  % Clear the sequence at the start
  \seq_clear:N \trSeq

  % Process each item in the list and store in the sequence
  \clist_map_inline:nn { #1 }
   {
    \seq_put_left:Nn \trSeq { \transpose{##1} }
   }

  % Output all items in the sequence, separated by spaces (or any other delimiter)
  \seq_use:Nn \trSeq {}
 }
\ExplSyntaxOff


\ExplSyntaxOn

%% % Define a sequence variable to store processed items
\seq_new:N \imSeq

% Define the command to process and display the list
\NewDocumentCommand{\innerMany}{m}
 {
  % Clear the sequence at the start
  \seq_clear:N \imSeq

  % Process each item in the list and store in the sequence
  \clist_map_inline:nn { #1 }
   {
    \seq_put_left:Nn \imSeq { \transpose{##1} }
    \seq_put_right:Nn \imSeq { {##1} }
   }

  % Output all items in the sequence, separated by spaces (or any other delimiter)
  \seq_use:Nn \imSeq {}
 }
\ExplSyntaxOff


% # 1 vel #2 mass # 3 partial d of vel
\newcommand{\quadraticd}[3]{ #3^\top #2 #1 +  #1^\top #2 #3}
\newcommand{\pTwo}[1]{#1^{2}}


%%%%%%%%%%%%%%%%%%%%%%%%%% Cross product
\newcommand{\paral}[2]{#1 \parallel #2}
\newcommand{\perpen}[2]{#1 \perp #2}

\newcommand{\bs}{\boldsymbol}
\newcommand{\backfill}{\hfill\(\blacksquare\)}

\newcommand{\skewMatrix}[1]{\widehat{#1}}

\newcommand{\skewMatrixTwo}[1]{({#1})^{\widehat{~}}}


\newcommand{\vectriple}[3]{#1\times(#2 \times #3)}
\newcommand{\vectripleexpan}[3]{
({#1}\bullet{#3}){#2} - ({#1}\bullet{#2}){#3}
}
\newcommand{\vectripleexpannegate}[3]{
-({#1}\bullet{#3}){#2} + ({#1}\bullet{#2}){#3}
}
\newcommand{\scalartriple}[3]{{#1}\bullet({#2} \times {#3})}
\newcommand{\vcd}[3]{({#1}\bullet{#3}){#2} - ({#1}\bullet{#2}) {#3}}  % Vector cross product definition

\newcommand{\inner}[2]{#1\bullet#2} % Inner product
\newcommand{\innerDef}[2]{\transpose{{#1}}#2} % Inner product
\newcommand{\normal}[1]{\bs{n}_{{#1}}}

\newcommand{\nc}[1]{\frac{#1}{| #1| }}% normal calc

\newcommand{\transpose}[2][default]{
  \IfStrEq{#1}{inv}
  {
    % Inverse transpose:
    {#2}^{-\top}
  }
  {
    % Default:
    {#2}^\top
  }
}

\newcommand{\transInverse}[1]{#1^{-\top}}
\newcommand{\pseudoInverse}[1]{{#1}^{+}}
% When  #row >  #column
\newcommand{\pseudoInverseColumnDef}[1]{\inverse{( \transpose{#1} #1 )}\transpose{#1}}
% When  #column >  #row
\newcommand{\pseudoInverseRowDef}[1]{\transpose{#1}\inverse{( #1 \transpose{#1} )}}

\newcommand{\weightedPseudoInverse}[1]{#1^{\dagger}}
% % #1 Jacobian #2 The weight, e.g., the inertia matrix.
\newcommand{\weightedPseudoInverseRowDef}[2]{\inverse{#2}\transpose{#1}\inverse{( #1 \inverse{#2} \transpose{#1} )}}
\newcommand{\weightedPseudoInverseColumnDef}[2]{\inverse{( \transpose{#1} \inverse{#2} #1 )}\transpose{#1}\inverse{#2}}

\newcommand{\inverse}[2][default]{
\IfStrEq{#1}{trans}
{
{#2}^{-\top}
}{
{#2}^{-1}
}
}


%%%%%%%%%%%%%%%%%%%%%%%%%% Resultant torque
% The argument are: force, torque, application point of the input wrench, application point of the resultant wrench.
\newcommand{\resultantTorque}[4]{#2 + (#3 - #4)\times #1}
\newcommand{\resultantTorqueTwo}[4]{#2 + \pointer{#4}{#3}\times #1}

\newcommand{\torqueFrame}[2]{{\torque_{#1#2}}}

\newcommand{\frameSymbol}{\mathcal{F}}


\newcommand{\applyFrameSubscript}[1]{%
    \ifthenelse{\equal{#1}{inertial}}{\inertialFrame}{%
    \ifthenelse{\equal{#1}{cmm}}{\cmmFrame}{%
    \ifthenelse{\equal{#1}{com}}{\com}{%
    \ifthenelse{\equal{#1}{fb}}{\fbFrame}{%
    \ifthenelse{\equal{#1}{torso}}{\torsoFrame}{%
    \ifthenelse{\equal{#1}{ee}}{\text{e}}{%
    \ifthenelse{\equal{#1}{cp}}{\cpFrame}{%
    \text{$#1$} % Default case, just print the argument as text
    }}}}}}}%
}

\NewDocumentCommand{\cframe}{O{}}{{\frameSymbol_{\applyFrameSubscript{#1}}}}



\newcommand{\origin}{\point[][][\inertialFrame]}


% The Gravito-inertial wrench
\newcommand{\giwrench}{\bs{W}^{gi}}
% The contact wrench
\newcommand{\cwrench}{\bs{W}^c}
\newcommand{\numEE}{m}

\newcommand{\forceJacobian}[2]{\jacobian_{#1 #2}}

\newcommand{\configurationFrame}[2]{g_{#1 #2}}



\newcommand{\spatialJacobian}[2]{
  \textcolor{spatialVelColor}{
    \spatial{\jacobian}_{#1 #2}
  }
}

\newcommand{\spatialJacobiand}[2]{
  \textcolor{spatialVelColor}{
    \dot{\spatial{\jacobian}}_{#1 #2}
  }
}

\newcommand{\forceJacobianDef}[2]{
\begin{bmatrix}
  \rotation{#1}{#2} \\
  \pointer{#1}{#2} \times \rotation{#1}{#2}
\end{bmatrix}}

% Transform from #1 to #2
\newcommand{\svaForceJacobianDef}[2]{
  \textcolor{svaColorVariable}{
    \vectorTwo{
      \translationSkew{#2}{#1}\rotationInv{#1}{#2}
    }{
      {\rotationInv{#1}{#2}}
    }
  }
}
% Transform from #1 to #2
\newcommand{\svaForceJacobianDefTwo}[2]{
  \textcolor{svaColorVariable}{
    \vectorTwo{
      -\rotation{#2}{#1}\translationSkew{#1}{#2}
    }{
      {\rotation{#2}{#1}}
    }
  }
}


\newcommand{\mass}{\text{m}}

\newcommand{\applyNumSymbol}[1]{%
       \IfStrEq{#1}{upper}{\text{N}}{%
         \text{n}
    }%
}

\newcommand{\applyNumSuperscript}[1]{%
    \IfStrEqCase{#1}{%
    {float}{
    \text{fb}
    }%
    {actuated}{
    \text{at}
    }%
    {ieq}{
    \text{ieq}
    }%
    {eq}{
    \text{eq}
    }%
    }[] % "default" string
}
\newcommand{\applyNumSubscript}[1]{%
    \IfStrEqCase{#1}{%
    {obj}{
    \text{obj}
    }%
    {cons}{
    \text{cons}
    }%
    {joints}{
    \text{j}
    }%
    {contact}{
    \text{sc}
    }%
    {impact}{
    \impulse
    }%
    }[] % "default" string
}

\NewDocumentCommand{\num}{O{default}O{default}O{default}}{%
    \ensuremath{%
        % #3 Shall we use the lower case symbol n?
        \applyNumSymbol{#3}_{\applyNumSubscript{#1}}% Check the first optional argument for obj, cons, joints, contact, impact,
        % Check the second optional argument for actuated, floating-base
        ^{\applyNumSuperscript{#2}}
    }
}


\newcommand{\numContacts}{{\text{m}_{\text{sc}}}} % The number of sustained contacts

\newcommand{\numImpacts}{{\text{m}_{\text{im}}}}


\newcommand{\numFree}{\numEE - \numContacts -\numImpacts }
\newcommand{\contactPoint}{\text{c}}

% End-effector
\newcommand{\ee}{\text{e}}

%%%%%%%%%%%%%%%%%%%%%%%%%%%%%%%% wrench, force, and torque
\newcommand{\applyWrenchPrescript}[1]{%
#1
}

\newcommand{\applyWrenchSuperscript}[1]{%
    \IfStrEqCase{#1}{%
        {ref}{\text{ref}}%
        {des}{\text{des}}%
        {mea}{\circ}%
        {min}{\text{min}}%
        {max}{\text{max}}%
    }[#1] % "default" string
}

\newcommand{\applyRelativeFrameSubscript}[2]{%
    {#1
    \ifx\hfuzz#1\hfuzz
         % #1 is empty
    \else % #1 is not empty
        \ifx\hfuzz#2\hfuzz
        % #2 is empty
        \else
          ,
          \fi
    \fi
    #2}
}

\newcommand{\wrenchSymbol}{{\textcolor{textColor}{\text{\bf w}}}\xspace}

\NewDocumentCommand{\test}{O{}O{}}{%
    {\ensuremath{
        {\wrenchSymbol_{#1}}_{#2}
    }
    }
}


\NewDocumentCommand{\wrench}{O{}O{}O{}O{}}{%
    {\ensuremath{
        % #1 reference frame #2 application point frame
        ~^{\applyWrenchPrescript{#4}}
        \wrenchSymbol_{\applyRelativeFrameSubscript{#1}{#2}}
        % #3 ref, mea, or des
        ^{\applyWrenchSuperscript{#3}}
    }
    }
}

\newcommand{\applyForceScalarSuperscript}[1]{%
    \IfStrEqCase{#1}{%
        {x}{\text{x}}%
        {y}{\text{y}}%
        {z}{\text{z}}%
    }[] % "default" string
}


\newcommand{\forceSymbol}{{{\bf f}}\xspace}
\NewDocumentCommand{\force}{O{}O{}O{}O{}}{%
{    \ensuremath{
      %% ^{\applyWrenchPrescript{#4}}
        % #1 reference frame #2 application point frame
       \bs{f}_{\applyRelativeFrameSubscript{#1}{#2}}
        ^{
          \applyForceScalarSuperscript{#4}% #4 x,y, or z
          \ifx\hfuzz#3\hfuzz
          % #3 is empty
          \else % #3 is not empty
          \ifx\hfuzz#4\hfuzz
          % #4 is empty
          \else
          ,
          \fi
          \fi
          %% \ifx\hfuzz#3\hfuzz\else
          %% \ifx\hfuzz#4\hfuzz\else
          %%   ,\applyWrenchSuperscript{#3}%
          %% \fi
          %% \fi
          \applyWrenchSuperscript{#3}
          }
        }
}
}

\NewDocumentCommand{\torque}{O{}O{}O{}O{}}{%
{    \ensuremath{
        % #1 reference frame #2 application point frame
       \bs{\tau}_{\applyRelativeFrameSubscript{#1}{#2}}
        ^{
          \applyForceScalarSuperscript{#4}% #4 x,y, or z
          \ifx\hfuzz#3\hfuzz\else
          \ifx\hfuzz#4\hfuzz\else
            ,\applyWrenchSuperscript{#3}%
          \fi
          \fi
          \applyWrenchSuperscript{#3}
          }
        }
}
}

\newcommand{\applyDurationSubscript}[2]{%
    \IfStrEqCase{#1}{%
    {stance}{
      {#2}_{\text{st.}}
    }%
    {flight}{
      {#2}_{\text{flt.}}
    }%
    {nominal}{
      {#2}_{\text{nom.}}
    }%
    {horizon}{
      {#2}_{\text{hor.}}
    }%
    }[
    #2
    ] % "default" string
}


%%%%%%%%%%%%%%%%%%%%%%%%%%%%%%%% ZMP, COM, and DCM

\newcommand{\applySuperscript}[1]{%
    \IfStrEqCase{#1}{%
    {ref}{
    \text{ref}
    }%
    {des}{
    \text{des}
    }%
    {mea}{
    \circ
    }%
    {est}{
    \wedge
    }%
    {min}{
    \text{min}
    }%
    {max}{
    \text{max}
    }%
    }[#1] % "default" string
}

\newcommand{\applyOverheadSymbols}[2]{%
    \IfStrEqCase{#1}{%
    {dddot}{
    \dddot{#2}
    }%
    {ddot}{
    \ddot{#2}
    }%
    {dot}{
    \dot{#2}
    }%
    }[
    #2
    ] % "default" string
}

\newcommand{\applyPostscript}[1]{%
    \IfStrEqCase{#1}{%
    {time}{
    (t)
    }%
    {current}{
    (k)
    }%
    {next}{
    (k+1)
    }%
    {previous}{
    (k-1)
    }%
    }[#1] % "default" string
}

\newcommand{\applySubscript}[1]{%
    \IfStrEqCase{#1}{%
    {plane}{
    \text{x,y}
    }%
    {x}{
    \text{x}
    }%
    {y}{
    \text{y}
    }%
    {z}{
    \text{z}
    }%
    }[#1] % "default" string
}




\newcommand{\zmpSymbol}{\text{\bf z}}
\NewDocumentCommand{\zmp}{O{}O{}O{}O{}}{%
{    \ensuremath{%
        % Check the first optional argument for dot or double dot
        %% \applyOverheadSymbols{#1}\text{\bf z}
        \applyOverheadSymbols{#1}{\zmpSymbol}
        % Check the second optional argument for 'ref'
        ^{\applySuperscript{#2}}%
        _{\applySubscript{#3}}
        \applyPostscript{#4}
    }
    }
}



\newcommand{\comSymbol}{{\textcolor{textColor}{\text{\bf c}}}\xspace}
\NewDocumentCommand{\com}{O{}O{}O{}O{}}{%
{    \ensuremath{%
        % Check the first optional argument for dot or double dot
        \applyOverheadSymbols{#1}{\comSymbol}
        % Check the second optional argument for 'ref'
        ^{\applySuperscript{#2}}
        _{\applySubscript{#3}}
        \applyPostscript{#4}
    }
    }
}

\newcommand{\dcmSymbol}{\textcolor{textColor}{\bs{\xi}}}
\NewDocumentCommand{\dcm}{O{}O{}O{}O{}}{%
{    \ensuremath{%
        % Check the first optional argument for dot or double dot
        \applyOverheadSymbols{#1}{\dcmSymbol}
        % Check the second optional argument for 'ref'
        ^{\applySuperscript{#2}}%
        _{\applySubscript{#3}}
        \applyPostscript{#4}
    }
    }
}

\newcommand{\pointSymbol}{\boldsymbol{q}}
\NewDocumentCommand{\point}{O{}O{}O{}O{}}{%
  {\ensuremath{%
        %% % Check the first optional argument for dot or double dot
        %% \IfStrEq{#1}{ddot}{\ddot{\boldsymbol{q}}}{%
        %%     \IfStrEq{#1}{dot}{\dot{\boldsymbol{q}}}{\boldsymbol{q}}%
        %% }%
        \applyOverheadSymbols{#1}{\pointSymbol}
        % Check the second optional argument for 'ref'
        ^{\applySuperscript{#2}}%
        % The mandatory subscript
        \ifx&#3&%
        % #3 is empty
        \else
        _{#3}%
        \fi
        % Post script
        \applyPostscript{#4}
    }
  }
}


\newcommand{\applySigmoidSymbols}[2]{%
    \IfStrEqCase{#1}{%
    {def}{
      \frac{\exp{#2}}{ 1 + \exp{#2}}
    }%
    %
    }[
      \sigmoidSymbol(#2)
    ] % "default" string
}


\newcommand{\sigmoidSymbol}{\sigma}
\NewDocumentCommand{\sigmoid}{O{default}O{default}}{%
  {\ensuremath{%
    \applySigmoidSymbols{#2}{#1} % #1, argument #2, def or not
    }
  }
}


\newcommand{\durationSymbol}{\text{T}}
\NewDocumentCommand{\duration}{O{default}}{%
  {\ensuremath{%
    \applyDurationSubscript{#1}{\durationSymbol}
    }
  }
}

\NewDocumentCommand{\pointFrame}{O{default}O{default}O{default}O{default}}{%
{    \ensuremath{%
       \point[#3][#4]_{\applyRelativeFrameSubscript{#1}{#2}} % #1 Reference frame, #2 Notation
    }
}
}


\newcommand{\forceScalar}{f}
\newcommand{\forceFrame}[2]{{\force^{#1}_{#2}}}


% The arguments are: (1) frame (2) Application point.
\newcommand{\impulseFrame}[2]{{\impulse^{#1}_{#2}}}
\newcommand{\applyWrenchFrameSuperscript}[1]{%
    \IfStrEq{#1}{ref}{, \text{ref}}{%
        \IfStrEq{#1}{des}{, \text{des}}{%
            \IfStrEq{#1}{mea}{, \circ}{}%
        }%
    }%
}
\NewDocumentCommand{\wrenchFrame}{mmO{default}}{%
  % #1 = representation frame
  % #2 = application point
  % #3 = optional argument, specifies ref, des, or mea
   \wrench^{{#1}\applyWrenchFrameSuperscript{#3}}_{#2}
}

% The foundamental subspaces
% #1 The matrix
\newcommand{\nullSpace}[1]{{N(#1)}}
% #1 The range (column) space
\newcommand{\colSpace}[1]{C(#1)}
\newcommand{\csProjector}[1]{{P}_{#1}}
% #1 Jacobian with: #col > #row
\newcommand{\csProjectorDef}[1]{\pseudoInverse{#1}#1}
% #1 Twist of a joint with: #col == 1
\newcommand{\jcsProjectorDef}[1]{#1\transpose{#1}}

\newcommand{\nsProjector}[1]{{N}_{#1}}
% #1 Jacobian with: #col > #row
\newcommand{\nsProjectorDef}[1]{(\identityMatrix - \csProjector{#1})}
\newcommand{\nsProjectorDefDef}[1]{(\identityMatrix - \csProjectorDef{#1})}
% #1 Twist of a joint with: #col == 1
\newcommand{\jnsProjectorDef}[1]{(\identityMatrix - \jcsProjectorDef{#1})}


% #1 The row  space
\newcommand{\rowSpace}[1]{C(\transpose{#1})}
% #1 The null space of the rows
\newcommand{\rowNullSpace}[1]{N(\transpose{#1})}
\newcommand{\rnsProjector}[1]{{N}_{\transpose{#1}}}
% #1 Jacobian with: #col > #row
\newcommand{\rnsProjectorDef}[1]{\transpose{(\identityMatrix - \csProjectorDef{#1})}}
% #1 Twist of a joint with: #col == 1
\newcommand{\jrnsProjectorDef}[1]{\transpose{(\identityMatrix - \jcsProjectorDef{#1})}}

\newcommand{\rsProjector}[1]{{P}_{\transpose{#1}}}
\newcommand{\rsProjectorDef}[1]{\transpose{(\csProjector{#1})}}
% #1 Twist of a joint with: #col == 1
\newcommand{\jrsProjectorDef}[1]{\transpose{(\jcsProjectorDef{#1})}}

% Joint row null space projector (to obtain the constraint wrench P*LagrangeMultiplier)
\newcommand{\jcwProjector}{T}

% The joint torque null-space projector of a redundant robot.
% #1 The Jacobian: v = J \dot{q}
\newcommand{\jtnsProjector}[1]{N_{\transpose{#1}}}
\newcommand{\jtnsProjectorDef}[1]{\transpose{(\identityMatrix - \gJacobianInv{#1} #1)}}

% The generalized Jacobain inverse (inertia-weighted pseudoinverse):
% #1 The Jacobian: v = J \dot{q},
% \gJacobianInv{J} = J^{-1}, when the manipulator is non-redundant.
\newcommand{\gJacobianInv}[1]{\overline{#1}}
\newcommand{\gJacobianInvDef}[1]{\inverse{\inertiaMatrix} \transpose{#1}\osInertia }

% The arguments are: (1) Equivalent frame name (2) Application point of the force
\newcommand{\forceGraspMatrix}[2]{
\begin{bmatrix}
    \rotation{\bs{#1}}{#2}^\top \\
    \pointer{#1}{#2} \times \rotation{#1}{#2}^\top
  \end{bmatrix}
}
\newcommand{\graspMatrix}[2]{
  {G_{#1#2}}
}

\newcommand{\graspMatrixDef}[2]{
\begin{bmatrix}
  \rotation{#1}{#2} & 0\\
  \pointer{#1}{#2} \times \rotation{#1}{#2} & \rotation{#1}{#2}
\end{bmatrix}
}


\newcommand{\pendulumC}{w}
\newcommand{\const}{\text{C}}
\renewcommand{\exp}[1]{e^{#1}}

\newcommand{\pendulumCDef}{\sqrt{\frac{\gSymbol}{\com[][][z]}}}


% Hadamard product
\newcommand{\hProduct}[2]{{#1}\bigodot {#2}}


%%%%%%%%%%%%%%%%%%%%%%%%%%%%%%%%%%%%%%%%%%%%%%%%%%%%%%%%%%%%%%%%%%%%%%%%%%%%%%%%
%% 6. COORDINATE FRAMES
%%%%%%%%%%%%%%%%%%%%%%%%%%%%%%%%%%%%%%%%%%%%%%%%%%%%%%%%%%%%%%%%%%%%%%%%%%%%%%%%

\newcommand{\dFrame}{\text{des}}

% The reference frame
\newcommand{\rFrame}{\text{ref}}

% Operational space coriolios and centrifugal forces

\newcommand{\coFrame}{\mathcal{C}}



\newcommand{\hybridFrame}{H}
\newcommand{\cmmFrame}{\text{G}}
\newcommand{\fbFrame}{\text{B}}
\newcommand{\torsoFrame}{\text{t}}
\newcommand{\inertialFrame}{\text{O}}
\newcommand{\anchor}{\text{A}} % The anchor frame of a legged robot
% The inertialFrame
\newcommand{\iFrame}{\text{O}}

\newcommand{\bodyFrame}{\text{b}}
\newcommand{\spatialFrame}{\text{s}}


\newcommand{\nz}{\text{\bf n}}
\newcommand{\nzDef}{
\begin{bmatrix}
0\\
0\\
1
\end{bmatrix}
}

\newcommand{\linkFrame}{L}
\newcommand{\jointFrame}{J}

% Gravity force

%%%%%%%%%%%%%%%%%%%%%%%%%%%%%%%%%%%%%%%%%%%%%%%%%%%%%%%%%%%%%%%%%%%%%%%%%%%%%%%%
%% 7. RIGID BODY GEOMETRY
%%%%%%%%%%%%%%%%%%%%%%%%%%%%%%%%%%%%%%%%%%%%%%%%%%%%%%%%%%%%%%%%%%%%%%%%%%%%%%%%

\newcommand{\epMapDef}[1]{
\identityMatrix + #1  + \frac{#1^2}{2!}  + \frac{#1^3}{3!} + \ldots
}
%%% #1:unit skew matrix #2 realnumber
\newcommand{\unitSkewExpDef}[2]{
\identityMatrix + \skewMatrix{#1}sin(#2) + \skewMatrix{#1}^2(1-cos(#2))
}

%%%%%%%%%%% -------------------------------------- Jacobian:
%%%% #1 reference frame #2 body frame #3 joint angle
\newcommand{\transformConfig}[3]{
  \textcolor{geometricColorVariable}{
    \transform{#1}{#2}(#3)
  }}
\newcommand{\transformConfigDef}[3]{
  \textcolor{geometricColorVariable}{
    \epMap{\twist{#1}{#2}#3} \transform{#1}{#2}(0)
  }}

\newcommand{\transformdConfigDef}[3]{
  \textcolor{geometricColorVariable}{
    \twist{#1}{#2} \dot{#3}\epMap{\twist{#1}{#2}#3} \transform{#1}{#2}(0)
  }}
\newcommand{\transformdConfigDefTwo}[3]{
  \textcolor{geometricColorVariable}{
    \epMap{\twist{#1}{#2}#3}\twist{#1}{#2} \dot{#3} \transform{#1}{#2}(0)
  }}

\newcommand{\transformddConfigDef}[3]{
  \textcolor{geometricColorVariable}{
    % \twist \ddot{#3}\epMap{\twist#3} \transform{#1}{#2}(0)
    % +\twist \dot{#3}\twist \dot{#3}\epMap{\twist#3} \transform{#1}{#2}(0)
    (\twist{#1}{#2} \ddot{#3}%\epMap{\twist#3} \transform{#1}{#2}(0)
    +\twist{#1}{#2} \twist{#1}{#2} \dot{#3}^2)\epMap{\twist{#1}{#2}#3} \transform{#1}{#2}(0)
  }}
\newcommand{\transformddConfigDefTwo}[3]{
  \textcolor{geometricColorVariable}{
    % \twist \ddot{#3}\epMap{\twist#3} \transform{#1}{#2}(0)
    % +\twist \dot{#3}\twist \dot{#3}\epMap{\twist#3} \transform{#1}{#2}(0)
    \epMap{\twist{#1}{#2}#3} (\twist{#1}{#2} \ddot{#3}%\epMap{\twist#3} \transform{#1}{#2}(0)
    +\twist{#1}{#2} \twist{#1}{#2} \dot{#3}^2)\transform{#1}{#2}(0)
  }}

% #1:reference frame, #:2 end-effector frame
\newcommand{\poeDef}[2]{
  \textcolor{geometricColorVariable}{
    \twistEp{#1}{1}{1}\twistEp{#1}{2}{2}\ldots\twistEp{#1}{#2}{#2}\transformConfig{#1}{#2}{\zeroVector}
  }}
\newcommand{\invPoeDef}[2]{
  \textcolor{geometricColorVariable}{
    \transformInvConfig{#1}{#2}{\zeroVector}\invTwistEp{#1}{#2}{#2}\ldots\invTwistEp{#1}{2}{2}\invTwistEp{#1}{1}{1}
  }}

%%%% #1 reference frame #2 body frame #3 joint angle
\newcommand{\transformInvConfig}[3]{
  \textcolor{geometricColorVariable}{
    \transformInv{#1}{#2}(#3)
  }}
\newcommand{\transformdConfig}[3]{
  \textcolor{geometricColorVariable}{
    \transformd{#1}{#2}(#3)
  }}
\newcommand{\transformddConfig}[3]{
 \textcolor{geometricColorVariable}{
   \transformdd{#1}{#2}(#3)
 }}
\newcommand{\transformInvConfigDef}[3]{
  \textcolor{geometricColorVariable}{
  \transformInv{#1}{#2}(0) \epMap{-\twist{#1}{#2}#3}
  }}



%%% Lie bracket
\newcommand{\epMap}[1]{e^{#1}}

\newcommand{\rotation}[2]{R_{#1 #2}}




\newcommand{\rotationInv}[2]{R^{\top}_{#1 #2}}
\newcommand{\rotationInvD}[2]{\dot{R}^{\top}_{#1 #2}}
\newcommand{\rotationInvDD}[2]{\ddot{R}^{\top}_{#1 #2}}
\newcommand{\rotationd}[2]{\dot{R}_{#1 #2}}
\newcommand{\rotationdd}[2]{\ddot{R}_{#1 #2}}

% #1 reference frame (point) #2 point
\newcommand{\translationDef}[2]{
  \textcolor{geometricColorVariable}{
    \pointer{#1}{#2} =  #2 - #1
}}


\newcommand{\translation}[2]{
  \vc{p}_{#1 #2}
}

\newcommand{\translationInv}[2]{
   \textcolor{geometricColorVariable}{
     \translation{#2}{#1}
     }
}
\newcommand{\translationInvSkew}[2]{
   \textcolor{geometricColorVariable}{
     \translationSkew{#2}{#1}
     }
}
\newcommand{\translationInvDef}[2]{
   \textcolor{geometricColorVariable}{
     -\rotationInv{#1}{#2}\translation{#1}{#2}
     }
}


% Complexity of an algorithm
\newcommand{\translationSkew}[2]{\skewMatrix{\vc{p}}_{#1 #2}}
\newcommand{\translationdSkew}[2]{\skewMatrix{\dot{\vc{p}}}_{#1 #2}}
\newcommand{\translationd}[2]{\dot{\vc{p}}_{#1 #2}}
\newcommand{\translationdd}[2]{{\ddot{\vc{p}}_{#1 #2}}}

\newcommand{\identityMatrix}{\mathds{1}} %matrix % alternativ \mathds{1} or {\mathbb{1}}
\newcommand{\idMatrix}[1]{{\identityMatrix_{#1}}} %matrix % alternativ \mathds{1} or {\mathbb{1}}

\newcommand{\zeroVector}{\mathbf{0}} %matrix % alternativ \mathds{0} or {\mathbb{0}} or \mathcal{O}
\newcommand{\zeroMatrix}{0} %matrix % alternativ \mathds{0} or {\mathbb{0}} or \mathcal{O}

\newcommand{\transform}[2]{
  \textcolor{geometricColorVariable}{
    g_{#1#2}
  }}

\newcommand{\transformd}[2]{
 \textcolor{geometricColorVariable}{
   \dot{g}_{#1#2}
 }}
\newcommand{\transformdDef}[2]{
 \textcolor{geometricColorVariable}{
   \matrixTwo{
     \rotationd{#1}{#2}
   }{
     \translationd{#1}{#2}
   }{
     \zeroVector
   }{
     0
   }
 }}
\newcommand{\transformdd}[2]{
 \textcolor{geometricColorVariable}{
   \ddot{g}_{#1#2}
 }}
\newcommand{\transformddDef}[2]{
 \textcolor{geometricColorVariable}{
   \matrixTwo{
     \rotationdd{#1}{#2}
   }{
     \translationdd{#1}{#2}
   }{
     \zeroVector
   }{
     0
   }
 }}
\newcommand{\transformInv}[2]{
  \textcolor{geometricColorVariable}{
    g^{-1}_{#1#2}
  }}
\newcommand{\transformInvD}[2]{
  \textcolor{geometricColorVariable}{
    % \dot{g}^{-1}_{#1#2}
    \derivative{g^{-1}_{#1#2}}
  }}
\newcommand{\transformInvDConfig}[3]{
  \textcolor{geometricColorVariable}{
    % \dot{g}^{-1}_{#1#2}
    \derivative{g^{-1}_{#1#2}(#3)}
  }}
\newcommand{\transformInvDDef}[2]{
  \textcolor{geometricColorVariable}{
    \matrixTwo{
      \rotationInvD{#1}{#2}
    }{
      -\rotationInvD{#1}{#2} \translation{#1}{#2}  - \rotationInv{#1}{#2} \translationd{#1}{#2}
    }{
      \zeroVector
    }{
      0
    }
  }}

\newcommand{\transformDef}[2]{
  \textcolor{geometricColorVariable}{
    \matrixTwo{
      \rotation{#1}{#2}
    }{
      \translation{#1}{#2}
    }{
      \zeroVector
    }{
      1
    }
  }
}

\newcommand{\transformInvDef}[2]{
  \textcolor{geometricColorVariable}{
    \matrixTwo{
      \rotationInv{#1}{#2}
    }{
      -\rotationInv{#1}{#2} \translation{#1}{#2}
    }{
      \zeroVector
    }{
      1
    }
  }
}
\newcommand{\adg}[2]{
  \textcolor{geometricColorVariable}{
    Ad_{g_{#1#2}}
  }
}

\newcommand{\adgDef}[2]{
  \textcolor{geometricColorVariable}{
  \begin{bmatrix}
    \rotation{#1}{#2} & \translationSkew{#1}{#2}\rotation{#1}{#2} \\
    \zeroMatrix & \rotation{#1}{#2}
  \end{bmatrix}
}}

\newcommand{\adgInv}[2]{
  \textcolor{geometricColorVariable}{
  Ad^{-1}_{g_{#1#2}}
}}
\newcommand{\adgTwoInv}[2]{
  \textcolor{geometricColorVariable}{
  Ad_{g^{-1}_{#1#2}}
}}

\newcommand{\adgTrans}[2]{
  \textcolor{geometricColorVariable}{
  Ad^{\top}_{g_{#1#2}}
}}
\newcommand{\adgTransDef}[2]{
  \textcolor{geometricColorVariable}{
  \begin{bmatrix}
    \rotationInv{#1}{#2} & \zeroMatrix \\
     -\rotationInv{#1}{#2}\translationSkew{#1}{#2} & \rotationInv{#1}{#2}
\end{bmatrix}
}}


\newcommand{\adgInvDefTwo}[2]{
  \textcolor{geometricColorVariable}{
  \begin{bmatrix}
    \rotationInv{#1}{#2} & -(\rotationInv{#1}{#2}\translation{#1}{#2})^{\skewMatrix{}}\rotationInv{#1}{#2} \\
    \zeroMatrix & \rotationInv{#1}{#2}
\end{bmatrix}
}}
\newcommand{\adgInvDef}[2]{
  \textcolor{geometricColorVariable}{
  \begin{bmatrix}
    \rotationInv{#1}{#2} & -\rotationInv{#1}{#2}\translationSkew{#1}{#2} \\
    \zeroMatrix & \rotationInv{#1}{#2}
\end{bmatrix}
}}

% From frame #1 to frame #2
\newcommand{\adgInvTrans}[2]{
  \textcolor{geometricColorVariable}{
  Ad^{\top}_{g^{-1}_{#1#2}}
}}

\newcommand{\adgInvTransDef}[2]{
  \textcolor{geometricColorVariable}{
  \begin{bmatrix}
    \rotation{#1}{#2} & \zeroMatrix \\
    \translationSkew{#1}{#2}\rotation{#1}{#2}  & \rotation{#1}{#2}
  \end{bmatrix}
  }
}



\newcommand{\adgTransform}[3]{
  \textcolor{geometricColorVariable}{
  \transform{#1}{#2} \skewMatrix{#3} \transformInv{#1}{#2}
}}


%%%%%%%%%%%%%%%%%%%%%%%%%%%%%%%%%%%%%%%%%%%%%%%%%%%%%%%%%%%%%%%%%%%%%%%%%%%%%%%%
%% 8. SCREW THEORY
%%%%%%%%%%%%%%%%%%%%%%%%%%%%%%%%%%%%%%%%%%%%%%%%%%%%%%%%%%%%%%%%%%%%%%%%%%%%%%%%

\newcommand{\atanTwo}[2]{
atan2({#1}, {#2})
}


%%%%%%%%%%% ----------------------------------------Scew


%%%%%%%%%%%% Transform defined by a screw
%%% #1:axis, #2:theta, #3:point, #4:pitch(finite)
\newcommand{\screwTransform}[4]{
\matrixTwo{
  \epMap{\skewMatrix{#1}#2}
}{
  (\identityMatrix - \epMap{\skewMatrix{#1}#2})#3
  + #4 #1 #2
}{
\zeroVector
}{
1
}
}

%%%%%%%%%%%% Wrench defined by a screw
%%% #1:magnitude, #2:point, #3:direction, #4:pitch(finite)
%%% #1: reference frame #2: screw name
\newcommand{\geoWrenchDef}[2]{
  \screwM_{#2} \vectorTwo{
    \angularVel_{#2}
  }{
    -\cross{\angularVel_{#2}}{\translation{\origin{#1}}{#2}} + \pitch_{#2}\angularVel_{#2}
  }
}

%%%%%% When the pitch is infinite:
%%% #1: screw name
\newcommand{\geoWrenchDefTwo}[1]{
  \screwM_{#1} \vectorTwo{
    \zeroVector
  }{
    \angularVel_{#1}
  }
}

\newcommand{\distance}{d}
%%% #1:screw 1, #2:screw 2
\newcommand{\screwAngle}[2]{
  \alpha_{#1#2}
}
\newcommand{\screwAngleDef}[2]{
atan2(\inner{(\cross{\angularVel_{#1}}{\angularVel_{#2}})}{\unitVec_{#1#2}}, \inner{\angularVel_{#1}}{\angularVel_{#2}} ).
}
\newcommand{\unitVec}{\bs{n}}

% Calculates the reciprocal product
\newcommand{\recip}[2]{
  #1 \odot{}   #2
}

\newcommand{\recipScrewDef}[2]{
  \screwM_{#1}\screwM_{#2}\left(
    (\pitch_{#1} + \pitch_{#2})\cos\screwAngle{#1}{#2} - \distance\sin\screwAngle{1}{2}
  \right)
}

%%%%%%%%%%%% Pitch:
\newcommand{\pitch}{
  h
}



% Twist2Screw:  #1 linear vel #2 angular vel
\newcommand{\pitchTwistDef}[2]{
  \frac{\innerP{#2}{#1}}{\twoNorm{#2}}
}

% Wrench2Screw: #1 force #2 moment
\newcommand{\pitchWrenchDef}[2]{
  \frac{\innerP{#1}{#2}}{\twoNorm{#1}}
}


\newcommand{\axis}{
  l
}

%%% #1:point, #2:angularVel direction
\newcommand{\screwAxisDef}[2]{
  \{ #1 + \lambda#2 : \lambda \in \RRv{} \}
}

% When angularVel is Not zero: #1:linearVel, #2:angularVel, #3:lambda (scalar)
\newcommand{\axisTwistDefThree}[3]{
\frac{\cross{#2}{#1}}{\twoNorm{#2}} + \axisTwistDefTwo{#1}{#3}
}

% When angularVel is zero: #1:linearVel, #2:lambda (scalar)
\newcommand{\axisTwistDefTwo}[2]{
   #1 #2
}

% When force is NOT zero: #1:force, #2:torque, #3:lambda (scalar)
\newcommand{\axisWrenchDefThree}[3]{
\frac{\cross{#1}{#2}}{\twoNorm{#1}} + \axisWrenchDefTwo{#2}{#3}
}

% When force is zero: #1 torque #2: lambda (scalar)
\newcommand{\axisWrenchDefTwo}[2]{
   #1 #2
}

%%%%%%%%%%%% Magnitude:
\newcommand{\screwM}{
  M
}

% Wrench2Screw: #1: force,  if force is nonzero. Otherwise #1 is the torque
% Twist2Screw: #1: angularVel,  if angularVel is nonzero. Otherwise #1 is the linearVel
\newcommand{\screwMDef}[1]{
  \norm{#1}
}
%%%%%%%%%%%% Twist:

% #1 axis #2 theta
\newcommand{\twistE}[2]{
  e^{\skewMatrix{#1}#2}
}


% Computes the iner product
\newcommand{\recipVelDef}[4]{
  \innerP{#1}{#4} + \innerP{#3}{#2}
}

% Calculates the reciprocal product between two wrenches
% #1:force,  #2:torque,  #3:force, #4:torque.
\newcommand{\recipWrenchDef}[4]{
  \innerP{#1}{#4} + \innerP{#3}{#2}
}

% Screw
\newcommand{\screw}{S}
%%% #1 screw notation
\newcommand{\screwDef}[1]{
  (\screwM_{#1}, \pitch_{#1}, \point[][][#1], \angularVel_{#1})
}


% Detail the wrench along a screw
\newcommand{\wrenchAS}[2]{S}

% Detail the twist about a screw
%%% #1:axis, #2:\point, #3 pitch
\newcommand{\linearTwistAS}[3]{
-\cross{#1}{#2} + #3#1
}

\newcommand{\linearTwistASTwo}[3]{
  \cross{#2}{#1} + #3#1
}

% Detail the twist about a screw, when the pitch is zero:
%%% #1:axis, #2:\point,
\newcommand{\linearTwistASRevolute}[2]{
-\cross{#1}{#2}
}
\newcommand{\linearTwistASRevoluteTwo}[2]{
{#2}\cross{#1}
}



%%%%%%%%%%% Composite rigid body (CRB) algorithm:
\newcommand{\twistDef}[2]{
  \matrixTwo{
    #2
  }{
    #1
  }{\zeroVector
  }{0
  }
}
%%%%%%%%%%%%%%% The exponential map of a twist:
%% #1 linear part #2 angularPart   #3 theta
\newcommand{\twistEpDef}[3]{
  \matrixTwo{
    \epMap{\skewMatrix{#2} #3}
  }{
    (\identityMatrix - \epMap{\skewMatrix{#2} #3})
    (\cross{#2}{#1})
    +
    \projectionDef{#2}#1#3
  }{
    \zeroVector
  }{
    1
  }
}
%% #1: reference frame, #2 twist notation #3 theta notation
\newcommand{\twistEp}[3]{
\epMap{\twist{#1}{#2} \theta_{#3}}
}
%% #1: reference frame, #2 twist notation #3 theta notation
\newcommand{\invTwistEp}[3]{
\epMap{-\twist{#1}{#2} \theta_{#3}}
}


%%%%%%%%%%%%%%%%%%%%%% Dynamics
%%% #1:index, #2:last index, #3 component
\newcommand{\twistFromScrew}[1]{
  \screwM_{#1}\vectorTwo{
    \linearTwistAS{\angularVel_{#1}}{\point[][][#1]}{\pitch_{#1}}
  }{
    \angularVel_{#1}
  }
}

\newcommand{\twistFromScrewRevolute}[1]{
  \screwM_{#1}\vectorTwo{
    \linearTwistASRevolute{\angularVel_{#1}}{\point[][][#1]}
  }{
    \angularVel_{#1}
  }
}


%%%%%%%%%%%%%%%%%%%%%%%%%%%%%%%%%%%%%%%%%%%%%%%%%%%%%%%%%%%%%%%%%%%%%%%%%%%%%%%%
%% 9. TWIST COORDINATES AND VELOCITIES
%%%%%%%%%%%%%%%%%%%%%%%%%%%%%%%%%%%%%%%%%%%%%%%%%%%%%%%%%%%%%%%%%%%%%%%%%%%%%%%%

\newcommand{\twistVel}{\bs{V}}


\newcommand{\twoCases}[4]{
\begin{dcases}
  {#1} & {#2} \\
  {#3} & {#4}
\end{dcases}
}
\newcommand{\crossSixD}[1]{
  \begin{bmatrix}
    \cross{\angularVel_{#1}}{} & \cross{\linearVel_{#1}}{}\\
    0 &\cross{\angularVel_{#1}}{}
  \end{bmatrix}
}

\newcommand{\dualCrossSixD}[1]{
  \begin{bmatrix}
    \cross{\angularVel_{#1}}{} & 0 \\
    \cross{\linearVel_{#1}}{} &\cross{\angularVel_{#1}}{}
  \end{bmatrix}
}


\newcommand{\twist}[2]{
  \textcolor{spatialVelColor}{
    \skewMatrix{\xi}_{#1#2}
  }
}

%%% #1: infinitesimal translation generator #2:so{3}
\newcommand{\twistC}[2]{
  \textcolor{spatialVelColor}{
    \xi_{#1#2}
  }
}

% #1:reference frame, #2:notation of the twist
\newcommand{\twistCd}[2]{
  \textcolor{spatialVelColor}{
    \dot{\xi}_{#1#2}
  }
}

% #1:reference frame, #2:notation of the twist
\newcommand{\spatialTwist}[2]{
  \textcolor{spatialVelColor}{
    \hat{\xi}'_{#1#2}
  }}

% #1:reference frame, #2:notation of the twist
\newcommand{\spatialTwistC}[2]{
  \textcolor{spatialVelColor}{
    \xi'_{#1#2}
  }}

% #1:reference frame, #2:notation of the twist
\newcommand{\spatialTwistCd}[2]{
  \textcolor{spatialVelColor}{
    \dot{\xi}'_{#1#2}
  }}

% #1:reference frame, #2:notation of the twist
\newcommand{\bodyTwistC}[2]{
  \textcolor{bodyVelColor}{
    \xi^{\dagger}_{#1#2}
  }}
% #1:reference frame, #2:index of the twist
\newcommand{\bodyTwistCd}[2]{
  \textcolor{bodyVelColor}{
    \dot{\xi}^{\dagger}_{#1#2}
  }}

\newcommand{\bodyTwist}[2]{
  \textcolor{bodyVelColor}{
    \hat{\xi}^{\dagger}_{#1#2}
  }}

% #1:reference frame, #2:index  of the twist, #3:index of the starting link
%%% t_0 means that starting time.
\newcommand{\spatialTwistCDef}[3]{
  \textcolor{spatialVelColor}{
    Ad_{\epMap{\twist{#1}{#3(t_0)} \theta_{#3}}\ldots \epMap{\twist{#1}{#2-1(t_0)} \theta_{#2-1}} }\twistC{#1}{#2(t_0)}
    }
}

% #1:reference frame, #2:index of the twist, #3:index of the last link
\newcommand{\bodyTwistCDef}[3]{
  \textcolor{bodyVelColor}{
    Ad^{-1}_{\epMap{\twist{#1}{#2(t_0)} \theta_{#2}}\ldots \epMap{\twist{s}{#3(t_0)} \theta_{#3}}\transformConfig{s}{#3}{0} }\twistC{#1}{#2(t_0)}
    }
}


%%% The twist coordinate
%%% #1: notation of the joint
\newcommand{\twistCRevolute}[1]{
\vectorTwo{
  \linearTwistASRevolute{\angularVel_{#1}}{\point[][][#1]}
  }{
    \angularVel_{#1}
  }
}

%%% #1: notation of the joint
\newcommand{\twistCPrismatic}[1]{
\vectorTwo{
  \linearVel_{#1}
  }{
    \zeroVector
  }
}

%%% #1:linearVel, #2: angularVel
\newcommand{\twistCDef}[2]{(#1, #2)}

%%% #1:notation of screw
\newcommand{\vc}[1]{\bs{#1}}


\newcommand{\caseTwo}[4]{
  \begin{cases}
    #1 & #2 \\
    #3 & #4
  \end{cases}
}
\newcommand{\caseThree}[6]{
  \begin{cases}
    #1 & #2 \\
    #3 & #4 \\
    #5 & #6 \\
  \end{cases}
}

\newcommand{\translationVel}{
  \bs{v}
}

\newcommand{\linearVel}{
\translationVel
}
\newcommand{\linearVelScalar}{
  v
}

\newcommand{\angularVel}{
  \bs{w}
}

\newcommand{\angularVelScalar}{
  w
}



%%% #1 reference frame #2 body frame
\newcommand{\spatialVel}[2]{
  \textcolor{spatialVelColor}{
    \vc{V}^{\text{s}}_{#1#2}
  }}

%%% #1 reference frame #2 body frame
\newcommand{\spatialVelDef}[2]{
  \textcolor{spatialVelColor}{
\vectorTwo{
  \spatialTVDef{#1}{#2}
  }{
    % (\rotationd{#1}{#2}\rotationInv{#1}{#2})^{\vee}
    \spatialRVDef{#1}{#2}
  }
}}
\newcommand{\spatialVelTwistC}[2]{
  \textcolor{spatialVelColor}{
\vectorTwo{
  \spatialTV{#1}{#2}
  }{
    % (\rotationd{#1}{#2}\rotationInv{#1}{#2})^{\vee}
    \spatialRV{#1}{#2}
  }
}}
\newcommand{\bodyVelTwistC}[2]{
  \textcolor{bodyVelColor}{
\vectorTwo{
  \bodyTV{#1}{#2}
  }{
    % (\rotationd{#1}{#2}\rotationInv{#1}{#2})^{\vee}
    \bodyRV{#1}{#2}
  }
}}

\newcommand{\spatialVelTwistDefTwo}[2]{
  \textcolor{spatialVelColor}{
    \transformdConfig{#1}{#2}{\theta}\transformInvConfig{#1}{#2}{\theta}
  }}

\newcommand{\bodyVelTwistDefTwo}[2]{
  \textcolor{spatialVelColor}{
    \transformInvConfig{#1}{#2}{\theta}\transformdConfig{#1}{#2}{\theta}
}}


% Defines the reference frame and body frame of the angular velocity
%%% #1 reference frame #2 body frame
\newcommand{\hybridVel}[2]{
  \textcolor{spatialVelColor}{
    \vc{V}^{h}_{#1#2}
  }}


\newcommand{\hybridVelDef}[2]{
  \textcolor{spatialVelColor}{
    \vectorTwo{
      \translationd{#1}{#2}
    }{
      \spatialRV{#1}{#2}
    }
  }}
\newcommand{\hybridVelDefDef}[2]{
  \textcolor{spatialVelColor}{
    \vectorTwo{
      \translationd{#1}{#2}
    }{
      \spatialRVDef{#1}{#2}
    }
  }}
%%%%%%%%%%%%%%%% Convert hybrid vel to body vel
%%% #1 Reference frame  #2 body frame
\newcommand{\hTwoB}[2]{
  \begin{bmatrix}
    \rotationInv{#1}{#2} & \zeroMatrix \\
    \zeroMatrix & \rotationInv{#1}{#2}
  \end{bmatrix}
}

%%% #1 Reference frame  #2 body frame
\newcommand{\bTwoH}[2]{
  \begin{bmatrix}
    \rotation{#1}{#2} & \zeroMatrix \\
    \zeroMatrix & \rotation{#1}{#2}
  \end{bmatrix}
}

%%%%%%%%%%%%%%%% Convert hybrid vel to body vel
\newcommand{\spatialTV}[2]{
  \textcolor{spatialVelColor}{
  \translationVel^{\spatialFrame}_{#1#2}
}}
\newcommand{\spatialTVDef}[2]{
  \textcolor{spatialVelColor}{
    -\rotationd{#1}{#2}\rotationInv{#1}{#2}\translation{#1}{#2} + \translationd{#1}{#2}
}}
\newcommand{\spatialTA}[2]{
  \textcolor{spatialVelColor}{
    \dot{\translationVel}^{\spatialFrame}_{#1#2}
  }}
\newcommand{\spatialTADef}[2]{
  \textcolor{spatialVelColor}{
    \translationdd{#1}{#2} - \spatialRATwist{#1}{#2}\translation{#1}{#2} - \spatialRVTwist{#1}{#2}\translationd{#1}{#2}
  }}
\newcommand{\spatialTADefTwo}[2]{
  \textcolor{spatialVelColor}{
    \translationdd{#1}{#2} + \translationSkew{#1}{#2}\spatialRA{#1}{#2} +\translationdSkew{#1}{#2} \spatialRV{#1}{#2}
  }}
\newcommand{\spatialTADefDef}[2]{
  \textcolor{spatialVelColor}{
\translationdd{#1}{#2} -\rotationdd{#1}{#2}\rotationInv{#1}{#2} \translation{#1}{#2}
-\rotationd{#1}{#2} \rotationInvD{#1}{#2} \translation{#1}{#2} - \rotationd{#1}{#2}\rotationInv{#1}{#2}\translationd{#1}{#2}
  }}
\newcommand{\spatialRADef}[2]{
  \textcolor{spatialVelColor}{
    \rotationdd{#1}{#2}\rotationInv{#1}{#2} + \rotationd{#1}{#2}\rotationInvD{#1}{#2}
  }}
\newcommand{\bodyRADef}[2]{
  \textcolor{bodyVelColor}{
    \rotationInv{#1}{#2}\rotationdd{#1}{#2} + \rotationInvD{#1}{#2}\rotationd{#1}{#2}
  }}
\newcommand{\spatialAcc}[2]{
  \textcolor{spatialVelColor}{
    \dot{\vc{V}}^{\spatialFrame}_{#1#2}
  }
}
\newcommand{\spatialAccTwist}[2]{
  \textcolor{spatialVelColor}{
    \skewMatrix{\dot{\vc{V}}}^{\spatialFrame}_{#1#2}
  }
}
\newcommand{\bodyAcc}[2]{
  \textcolor{bodyVelColor}{
    \dot{\vc{V}}^{\bodyFrame}_{#1#2}
  }
}
\newcommand{\bodyAccTwist}[2]{
  \textcolor{bodyVelColor}{
    \skewMatrix{\dot{\vc{V}}}^{\bodyFrame}_{#1#2}
  }
}
\newcommand{\bodyAccTwistDef}[2]{
  \textcolor{bodyVelColor}{
    \matrixTwo{
      \bodyRATwist{#1}{#2}
    }{
      \bodyTA{#1}{#2}
    }{
      \zeroVector
    }{
      0
    }
  }
}
\newcommand{\spatialAccTwistDef}[2]{
  \textcolor{spatialVelColor}{
    \matrixTwo{
      \spatialRATwist{#1}{#2}
    }{
      \spatialTA{#1}{#2}
    }{
      \zeroVector
    }{
      0
    }
  }
}
\newcommand{\spatialAccDef}[2]{
  \textcolor{spatialVelColor}{
    \matrixTwo{
      \spatialRATwist{#1}{#2}
    }{
      \translationdd{#1}{#2} - \spatialRATwist{#1}{#2}\translation{#1}{#2} - \spatialRVTwist{#1}{#2}\translationd{#1}{#2}
    }{
      \zeroVector
    }{
      0
    }
  }}

\newcommand{\spatialAccTwistC}[2]{
  \textcolor{spatialVelColor}{
    \vectorTwo{
      \translationdd{#1}{#2} - \spatialRATwist{#1}{#2}\translation{#1}{#2} - \spatialRVTwist{#1}{#2}\translationd{#1}{#2}
    }{
      \spatialRA{#1}{#2}
    }
  }}

\newcommand{\spatialAccDefDef}[2]{
  \textcolor{spatialVelColor}{
    \matrixTwo{
      \rotationdd{#1}{#2}\rotationInv{#1}{#2} + \rotationd{#1}{#2}\rotationInvD{#1}{#2}
    }{
      \translationdd{#1}{#2} -\rotationdd{#1}{#2}\rotationInv{#1}{#2} \translation{#1}{#2}
      -\rotationd{#1}{#2} \rotationInvD{#1}{#2} \translation{#1}{#2} - \rotationd{#1}{#2}\rotationInv{#1}{#2}\translationd{#1}{#2}
    }{
      \zeroVector
    }{
      0
    }
  }}
\newcommand{\bodyAccDefDef}[2]{
  \textcolor{bodyVelColor}{
    \matrixTwo{
      \rotationInv{#1}{#2}\rotationdd{#1}{#2} + \rotationInvD{#1}{#2}\rotationd{#1}{#2}
    }{
      \rotationInvD{#1}{#2}\translationd{#1}{#2} + \rotationInv{#1}{#2}\translationdd{#1}{#2}
    }{
      \zeroVector
    }{
      0
    }
  }}
\newcommand{\bodyAccDef}[2]{
  \textcolor{bodyVelColor}{
    \matrixTwo{
      \bodyRATwist{#1}{#2}
    }{
      % \bodyTV{#1}{#2} + \rotationInv{#1}{#2}\translationdd{#1}{#2}
      \bodyTA{#1}{#2}
    }{
      \zeroVector
    }{
      0
    }
  }}


\newcommand{\spatialRA}[2]{
  \textcolor{spatialVelColor}{
    \dot{\angularVel}^{\spatialFrame}_{#1#2}
  }}
\newcommand{\spatialRATwist}[2]{
  \textcolor{spatialVelColor}{
    \skewMatrix{\dot{\angularVel}}^{\spatialFrame}_{#1#2}
  }}
\newcommand{\bodyRATwist}[2]{
  \textcolor{bodyVelColor}{
    \skewMatrix{\dot{\angularVel}}^{\bodyFrame}_{#1#2}
  }}

\newcommand{\spatialRV}[2]{
  \textcolor{spatialVelColor}{
    \angularVel^{\spatialFrame}_{#1#2}
  }}

\newcommand{\spatialRVDef}[2]{
  \textcolor{spatialVelColor}{
    (\rotationd{#1}{#2}\rotationInv{#1}{#2})^{\vee}
}}
\newcommand{\spatialRVDefTwo}[2]{
  \textcolor{spatialVelColor}{
    \rotationd{#1}{#2}\rotationInv{#1}{#2}
}}

\newcommand{\spatialRVd}[2]{
  \textcolor{spatialVelColor}{
    \dot{\angularVel}^{\spatialFrame}_{#1#2}
}}

\newcommand{\spatialTVd}[2]{
  \textcolor{spatialVelColor}{
    \dot{\translationVel}^{\spatialFrame}_{#1#2}
}}

\newcommand{\spatialVelTwist}[2]{
  \textcolor{spatialVelColor}{
    \skewMatrix{\vc{V}}^{\spatialFrame}_{#1#2}
  }}

\newcommand{\spatialVelTwistDef}[2]{
  \textcolor{spatialVelColor}{
    \matrixTwo
    {
      \spatialRVTwist{#1}{#2}
    }{
      \spatialTV{#1}{#2}
    }{
      \zeroVector
    }{
      0
    }
}}
\newcommand{\spatialTVTwist}[2]{
  \textcolor{spatialVelColor}{
    \skewMatrix{\vc{v}}^{\spatialFrame}_{#1#2}
}}
\newcommand{\spatialRVTwist}[2]{
  \textcolor{spatialVelColor}{
    \skewMatrix{\vc{w}}^{\spatialFrame}_{#1#2}
}}
\newcommand{\spatialRVdTwist}[2]{
  \textcolor{spatialVelColor}{
    \skewMatrix{\dot{\vc{w}}}^{\spatialFrame}_{#1#2}
  }}

%%% #1 Representation frame, #2 Base frame, and #3 Body frame
\newcommand{\vel}[3]{
  \textcolor{velColor}{
    \vc{V}^{#1}_{#2#3}
  }}
\newcommand{\velJacobian}[3]{
  \textcolor{velColor}{
    \jacobian^{#1}_{#2#3}
  }}
\newcommand{\tV}[3]{
  \textcolor{velColor}{
    \vc{v}^{#1}_{#2#3}
  }}

\newcommand{\rV}[3]{
  \textcolor{velColor}{
    \vc{w}^{#1}_{#2#3}
  }}

%%% #1 Reference frame #2 Body frame
\newcommand{\bodyVel}[2]{
  \textcolor{bodyVelColor}{
    \vc{V}^{\bodyFrame}_{#1#2}
  }}
%%% #1 Reference frame #2 Body frame
\newcommand{\bodyTVDef}[2]{
  \textcolor{bodyVelColor}{
    \rotationInv{#1}{#2}\translationd{#1}{#2}
  }}
%%% #1 Reference frame #2 Body frame
\newcommand{\bodyRVDef}[2]{
  \textcolor{bodyVelColor}{
      (\rotationInv{#1}{#2}\rotationd{#1}{#2})^{\vee}
  }}


\newcommand{\bodyRVDefTwo}[2]{
  \textcolor{bodyVelColor}{
    \rotationInv{#1}{#2}\rotationd{#1}{#2}
  }}



%%% #1 Reference frame #2 Body frame
\newcommand{\bodyVelDef}[2]{
  \textcolor{bodyVelColor}{
    \vectorTwo{
      \bodyTVDef{#1}{#2}
    }{
      \bodyRVDef{#1}{#2}
    }
  }}

%%% #1:Reference frame of the screw, #2:Body frame attached to the screw motion, #3: angle
\newcommand{\screwBodyAcc}[3]{
  \textcolor{bodyVelColor}{
    \Vee{\twistTransform{#1}{#2(0)}  \twistC{#1}{#2}}\ddot{#3}
  }}
%%% #1:Reference frame of the screw, #2:Body frame attached to the screw motion, #3: angle
\newcommand{\screwBodyVel}[3]{
  \textcolor{bodyVelColor}{
    \Vee{\twistTransform{#1}{#2(0)}  \twistC{#1}{#2}}\dot{#3}
  }}
%%% #1:Reference frame of the screw, #2:Body frame attached to the screw motion, #3: angle
\newcommand{\screwSpatialAcc}[3]{
  \textcolor{spatialVelColor}{
    \twistC{#1}{#2}\ddot{#3}
  }}

%%% #1:Reference frame of the screw, #2:Body frame attached to the screw motion, #3: angle
\newcommand{\screwSpatialVel}[3]{
  \textcolor{spatialVelColor}{
    \twistC{#1}{#2}\dot{#3}
  }}



\newcommand{\bodyVelTwist}[2]{
  \textcolor{bodyVelColor}{
    \skewMatrix{\vc{V}}^{\bodyFrame}_{#1#2}
  }}
\newcommand{\bodyVelTwistDef}[2]{
  \textcolor{bodyVelColor}{
    \matrixTwo
    {
      \bodyRVTwist{#1}{#2}
    }{
      \bodyTV{#1}{#2}
    }{
      \zeroVector
    }{
      0
    }
  }}
\newcommand{\bodyVelTwistDefDef}[2]{
  \textcolor{bodyVelColor}{
    \matrixTwo
    {
      \rotationInv{#1}{#2}\rotationd{#1}{#2}
    }{
      \bodyTVDef{#1}{#2}
    }{
      \zeroVector
    }{
      0
    }
}}
\newcommand{\bodyTV}[2]{
  \textcolor{bodyVelColor}{
    \vc{v}^{\bodyFrame}_{#1#2}
  }}
\newcommand{\bodyTA}[2]{
  \textcolor{bodyVelColor}{
    \dot{\vc{v}}^{\bodyFrame}_{#1#2}
  }}
\newcommand{\bodyTADef}[2]{
  \textcolor{bodyVelColor}{
    % \rotationInvD{#1}{#2}\translationd{#1}{#2} + \rotationInv{#1}{#2}\translationdd{#1}{#2}
    \rotationInv{#1}{#2}\translationdd{#1}{#2} - \bodyRVTwist{#1}{#2}\bodyTV{#1}{#2}
  }}
\newcommand{\bodyTADefDef}[2]{
  \textcolor{bodyVelColor}{
    \rotationInvD{#1}{#2}\translationd{#1}{#2} + \rotationInv{#1}{#2}\translationdd{#1}{#2}
  }}


\newcommand{\bodyRA}[2]{
  \textcolor{bodyVelColor}{
    \dot{\angularVel}^{\bodyFrame}_{#1#2}
  }}
\newcommand{\bodyTVd}[2]{
  \textcolor{bodyVelColor}{
    \dot{\vc{v}}^{\bodyFrame}_{#1#2}
}}
\newcommand{\bodyTVTwist}[2]{
  \textcolor{bodyVelColor}{
    \skewMatrix{\vc{v}}^{\bodyFrame}_{#1#2}
}}

\newcommand{\bodyRV}[2]{
  \textcolor{bodyVelColor}{
    \vc{w}^{\bodyFrame}_{#1#2}
}}
\newcommand{\bodyRVd}[2]{
  \textcolor{geometricColorVariable}{
    \dot{\vc{w}}^{\bodyFrame}_{#1#2}
}}

\newcommand{\bodyForce}{
\force^{\bodyFrame}
}
\newcommand{\bodyTorque}{
\torque^{\bodyFrame}
}

\newcommand{\bodyRVTwist}[2]{
  \textcolor{bodyVelColor}{
    \skewMatrix{\vc{w}}^{\bodyFrame}_{#1#2}
}}



\newcommand{\twistCTransform}[2]{
  \textcolor{geometricColorVariable}{
    \twistTransform{#1}{#2}
  }
}
\newcommand{\twistCTransformDef}[2]{
  \textcolor{geometricColorVariable}{
    \twistTransformDef{#1}{#2}
  }
}
% transform a point from coordinate #1 to coordinate #2
\newcommand{\pointTransform}[2]{
  \textcolor{geometricColorVariable}{
    \transformInv{#1}{#2}
  }
}
\newcommand{\pointTransformDef}[2]{
  \textcolor{geometricColorVariable}{
    \transformInvDef{#1}{#2}
  }
}
\newcommand{\pointTransformTwo}[2]{
  \textcolor{geometricColorVariable}{
    \transform{#2}{#1}
  }
}
\newcommand{\pointTransformTwoDef}[2]{
  \textcolor{geometricColorVariable}{
    \transformDef{#2}{#1}
  }
}
% transform from coordinate #1 to coordinate #2
\newcommand{\twistTransformTwo}[2]{
  \textcolor{geometricColorVariable}{
    \adg{#2}{#1}
  }
}

\newcommand{\twistTransformdTwo}[2]{
  \textcolor{geometricColorVariable}{
    \derivative{\adg{#2}{#1}}
  }
}
\newcommand{\twistTransformdTwoDef}[2]{
  \textcolor{geometricColorVariable}{
    -\adg{#2}{#1}(\spatialVel{#1}{#2}\times)
  }
}

\newcommand{\twistTransformTwoDef}[2]{
  \textcolor{geometricColorVariable}{
    \adgDef{#2}{#1}
  }
}

\newcommand{\twistTransform}[2]{
  \textcolor{geometricColorVariable}{
    \adgInv{#1}{#2}
  }
}


% The direct transform: re-write the coordinates.
% #1 current frame, #2 target frame.
\newcommand{\dTransform}[2]{
  \textcolor{geometricColorVariable}{
    \matrixTwo{\tdTransform{#1}{#2}}{0}{0}{\tdTransform{#1}{#2}}
  }
}

% #1 current frame, #2 target frame.
\newcommand{\tdTransform}[2]{
  \textcolor{geometricColorVariable}{
    \rotationInv{#1}{#2}
  }
}


\newcommand{\twistTransformDef}[2]{
  \adgInvDef{#1}{#2}
}
\newcommand{\twistTransformd}[2]{
  \textcolor{geometricColorVariable}{
    \derivative{\twistTransform{#1}{#2}}
    }
}
\newcommand{\twistTransformdDef}[2]{
  \textcolor{geometricColorVariable}{
    -\twistTransform{#1}{#2}(\spatialVel{#1}{#2}\times)
    }
  }

\newcommand{\twistTransformdDefTwo}[2]{
  \textcolor{geometricColorVariable}{
    \twistTransform{#1}{#2}(\bodyVel{#2}{#1}\times)
    }
}





% #1: reference frame , #2 link(body) frame, and #3: New reference frame.
\newcommand{\spatialVelTransform}[3]{
  \textcolor{spatialVelColor}{
    % \spatialVel{#3}{#1} + \spatialVelStaticTransform{#3}{#1}\spatialVel{#1}{#2}
    \spatialVel{#3}{#1} + \twistTransform{#1}{#3}\spatialVel{#1}{#2}
  }
}
% #1: reference frame , #2 link(body) frame, and #3: New reference frame.
\newcommand{\spatialAccTransform}[3]{
  \textcolor{spatialVelColor}{
    % \spatialVel{#3}{#1} + \spatialVelStaticTransform{#3}{#1}\spatialVel{#1}{#2}
    \spatialAcc{#3}{#1}
    +\twistTransform{#1}{#3} (\spatialVelTwist{#1}{#2}\spatialVelTwist{#1}{#3} - \spatialVelTwist{#1}{#3}\spatialVelTwist{#1}{#2})^\vee
    % +\twistTransform{#1}{#3} (\bodyVelTwist{#3}{#1}\bodyVelTwist{#2}{#1})^\vee
    + \twistTransform{#1}{#3}\spatialAcc{#1}{#2}
  }
}
% #1: reference frame , #2 link(body) frame, and #3: New reference frame.
\newcommand{\spatialAccTransformTwo}[3]{
  \textcolor{spatialVelColor}{
    \spatialAcc{#3}{#1}
    +\twistTransform{#1}{#3}(\cross{\spatialVel{#1}{#2}}{\spatialVel{#1}{#3}})
    + \twistTransform{#1}{#3}\spatialAcc{#1}{#2}
  }
}

% #1: reference frame , #2 link(body) frame, and #3: New reference frame.
\newcommand{\bodyAccTransform}[3]{
  \textcolor{bodyVelColor}{
    \bodyAcc{#1}{#2}
    + \twistTransform{#1}{#2}(
    \bodyVelTwist{#2}{#1}\bodyVelTwist{#3}{#1}
    -\bodyVelTwist{#3}{#1}\bodyVelTwist{#2}{#1}
    )^{\vee}
    +\twistTransform{#1}{#2}\bodyAcc{#3}{#1}
  }
}
% #1: reference frame , #2 link(body) frame, and #3: New reference frame.
\newcommand{\bodyAccTransformTwo}[3]{
  \textcolor{bodyVelColor}{
    \bodyAcc{#1}{#2}
    % \twistTransformdDef{#1}{#2}
    +\twistTransform{#1}{#2}
    (\cross{\bodyVel{#2}{#1}}
    {\bodyVel{#3}{#1}})
    +\twistTransform{#1}{#2}\bodyAcc{#3}{#1}
  }
}

\newcommand{\spatialVelTransformTwo}[3]{
  \textcolor{spatialVelColor}{
    % \spatialVel{#3}{#1} + \spatialVelStaticTransform{#3}{#1}\spatialVel{#1}{#2}
    \spatialVel{#3}{#1} +
    \twistTransformTwo{#1}{#3}\spatialVel{#1}{#2}
  }
}

% Expressed in #2
% #1: reference frame , #2 link(body) frame, and #3: New reference frame.
\newcommand{\bodyVelTransform}[3]{
  \textcolor{bodyVelColor}{
    \twistTransform{#1}{#2}\bodyVel{#3}{#1} + \bodyVel{#1}{#2}
  }
}
\newcommand{\bodyVelTransformTwo}[3]{
  \textcolor{bodyVelColor}{
    \twistTransformTwo{#1}{#2}\bodyVel{#3}{#1} + \bodyVel{#1}{#2}
  }
}


% Transform from #1 to #2
\newcommand{\bodyWrench}{\bs{W}^{\bodyFrame}}



%%%%%%%%%%%%%%%%%%%%%%%%% Inertia and momentum:
%   #1 reference frame name, #2 body frame name
%  locally, we can specify: #1 body frame: F_i  #2 the local centroidal frame: F_com_i

%%%%%%%%%%%%%%%%%%%%%%%%%%%%%%%%%%%%%%%%%%%%%%%%%%%%%%%%%%%%%%%%%%%%%%%%%%%%%%%%
%% 10. SPATIAL VECTOR ALGEBRA (Featherstone)
%%%%%%%%%%%%%%%%%%%%%%%%%%%%%%%%%%%%%%%%%%%%%%%%%%%%%%%%%%%%%%%%%%%%%%%%%%%%%%%%

\newcommand{\svaMT}[2]{
  \svaFT{#1}{#2}
}


\newcommand{\svaTransform}[2]{
  \textcolor{svaColorVariable}{
    {^{#2}T_{#1}}
  }
}
% Transform from #1 to #2
\newcommand{\svaTransformDef}[2]{
  \textcolor{svaColorVariable}{
  \matrixTwo{
    \rotation{#2}{#1}
  }{
    -\rotation{#2}{#1}\translation{#1}{#2}
  }{
    \zeroVector
  }{
    1
  }
  }
}
\newcommand{\svaTransformDefTwo}[2]{
  \textcolor{svaColorVariable}{
  \matrixTwo{
    \rotationInv{#1}{#2}
  }{
    -\rotationInv{#1}{#2}\translation{#1}{#2}
  }{
    \zeroVector
  }{
    1
  }
  }
}

\newcommand{\svaT}{
  \textcolor{svaColorVariable}{
    \bs{X}
  }
}

% Transform from #1 to #2 (#2 is represented in #1)
\newcommand{\svaVT}[2]{
  \textcolor{svaColorVariable}{
  {^{#2}\svaT_{#1}}
  % \svaT{#1}{#2}
  }
}

% Transform Vel from #1 to #2
\newcommand{\svaVTDef}[2]{
  % eq.~2.24 of FeatherStone
  % \matrixTwo{
  %   \rotationInv{#1}{#2}
  % }{\zeroMatrix}{
  %   -\rotationInv{#1}{#2}\translationSkew{#1}{#2}
  % }
  % {
  %   \rotationInv{#1}{#2}
  % }
  \textcolor{svaColorVariable}{
  \matrixTwo{
    \rotation{#2}{#1}
  }{\zeroMatrix}{
    -\rotation{#2}{#1}\translationSkew{#1}{#2}
  }
  {
    \rotation{#2}{#1}
  }
  }
}

% Transform Vel from #1 to #2
\newcommand{\svaVTDefTwo}[2]{
  % eq.~2.26 of FeatherStone
  \textcolor{svaColorVariable}{
  \matrixTwo{
    \rotationInv{#1}{#2}
  }{\zeroMatrix
  }{
    \translationSkew{#2}{#1}\rotationInv{#1}{#2}
  }
  {
    \rotationInv{#1}{#2}
  }
  }
}

% Transform from #1 to #2
\newcommand{\svaFT}[2]{
  \textcolor{svaColorVariable}{
    {^{#2}\svaT^{*}_{#1}}
  }
}

% Transform from #1 to #2
\newcommand{\svaFTDef}[2]{
  % Following eq.2.27 of Featherstone
  % \matrixTwo
  % {\rotationInv{#2}{#1}}
  % {
  %   \translationSkew{#1}{#2}\rotationInv{#2}{#1}
  % }{
  %   \zeroMatrix
  % }
  %   {\rotationInv{#2}{#1}}
  \textcolor{svaColorVariable}{
  \matrixTwo
  {\rotationInv{#1}{#2}}
  {
    \translationSkew{#2}{#1}\rotationInv{#1}{#2}
  }{
    \zeroMatrix
  }
  {\rotationInv{#1}{#2}}
  }
}

% Transform from #1 to #2
\newcommand{\svaFTDefTwo}[2]{
  % Following eq.2.25 of Featherstone
  % \matrixTwo{
  %   \rotationInv{#1}{#2}
  % }{
  %   -\rotationInv{#1}{#2}\translationSkew{#1}{#2}
  % }{
  %   \zeroMatrix
  % }
  %   {\rotationInv{#1}{#2}}
  \textcolor{svaColorVariable}{
  \matrixTwo{
    \rotation{#2}{#1}
  }{
    -\rotation{#2}{#1}\translationSkew{#1}{#2}
  }{
    \zeroMatrix
  }
  {\rotation{#2}{#1}}
  }
}

% Reference frame: #1, link frame: #2
\newcommand{\svaBodyGInertiaDef}[2]{
  \textcolor{svaColorVariable}{
    \matrixTwo{\bodyMetaInertia{#1}{#2} }{\zeroMatrix}{\zeroMatrix}{\mass_{#2} \identityMatrix}
  }
}


% #1 Reference frame #2 body frame
\newcommand{\svaMomentaTransform}[2]{
  \transpose{\svaVT{#2}{#1}}
}
\newcommand{\svaMomentaTransformDef}[2]{
  \transpose{\svaVTDef{#2}{#1}}
}


% #1 and #2 defines the spatial generalized inertia to be transformed.
% #1 Reference frame #2 body frame
% #3 defines the new reference frame
\newcommand{\svaSpatialGInertiaTransform}[3]{
  \transpose{\svaVT{#3}{#1}}\spatialGInertia{#1}{#2}\svaVT{#3}{#1}
}

\newcommand{\svaSpatialCrbInertiaTransform}[3]{
  \transpose{\svaVT{#3}{#1}}\spatialCrbInertia{#1}{#2}\svaVT{#3}{#1}
}

%%% Generalized inertia transform
% #1:new reference frame
% #2:current reference frame
% #3:current Generalized inertia
\newcommand{\svaGInertiaTransform}[3]{
  \transpose{\svaVT{#1}{#2}}
  #3
  \svaVT{#1}{#2}
}


% #1 Reference frame #2 body frame
\newcommand{\svaSpatialGInertia}[2]{
  \textcolor{svaColorVariable}{
    \transpose{\svaVT{#1}{#2}} \bodyGInertia{#1}{#2}\svaVT{#1}{#2}
  }
}

% #1 Reference frame #2 body frame
\newcommand{\svaSpatialGInertiaDef}[2]{
  \textcolor{svaColorVariable}{
    \transpose{\svaVTDef{#1}{#2}} \svaBodyGInertiaDef{#1}{#2}\svaVTDef{#1}{#2}
  }
}



%%%%%%%%%%%%%%%%%%%%%%%%%%%%%%%%%%%%%%%%%%%%%%%%%%%%%%%%%%%%%%%%%%%%%%%%%%%%%%%%
%% 11. JOINT KINEMATICS
%%%%%%%%%%%%%%%%%%%%%%%%%%%%%%%%%%%%%%%%%%%%%%%%%%%%%%%%%%%%%%%%%%%%%%%%%%%%%%%%

\newcommand{\svaJointJac}[1]{\bf{\Phi}_{#1}}


\newcommand{\applyJangles}[1]{%
    \IfStrEqCase{#1}{%
    {actuated}{
    \bs{\theta}
    }%
    {fb}{
    \bs{q}_{\text{fb}}
    }%
    {generalized}{
    \bs{\mathcal{X}}
    }%
    }[\bs{q}] % "default" string
}


\newcommand{\applyJanglesPostscript}[1]{%
    \IfStrEqCase{#1}{%
    {current}{
    {t_k}
    }%
    {next}{
    {t_{k+1}}
    }%
    {previous}{
    {t_{k-1}}
    }%
    }[] % "default" string
}


\NewDocumentCommand{\jangles}{O{}O{}O{}O{}}{%
{    \ensuremath{%
        % Check the first optional argument for dot or double dot
        \applyOverheadSymbols{#1}{\applyJangles{#3}} % Third argument for actuated, fb joints
        % 'ref', 'mea', 'des'
        ^{\applySuperscript{#2}}%
        _{\applyJanglesPostscript{#4}}
    }
    }
}

\newcommand{\jvelocities}{\mathbf{\dot{q}}}
\newcommand{\jaccelerations}{\mathbf{\ddot{q}}}
\newcommand{\jtorques}{\boldsymbol{\tau}}
\newcommand{\omegaVector}{\boldsymbol{\omega}}
\newcommand{\tsInertia}{\boldsymbol{\Lambda}}




\newcommand{\janglesJump}{\jump \jangles}
\newcommand{\jvelocitiesJump}{\jump \jvelocities}
\newcommand{\jangle}{q}

%%%%%%%%%%%%%%%%%%%%%%%%%%%%%%%%%%%%%%%%%%%%%%%%%%%%%%%%%%%%%%%%% Bound
\newcommand{\jvelocitiesUpperBound}{\mathbf{\dot{\bar{\jangle}}}}
\newcommand{\jvelocitiesLowerBound}{\mathbf{\dot{\ubar{\jangle}}}}
\newcommand{\jvelocitiesJumpUpperBound}{\jump \jvelocitiesUpperBound}
\newcommand{\jvelocitiesJumpLowerBound}{\jump \jvelocitiesLowerBound}

%%%%%%%%%%%%%%%%%%%%%%%%%%%%%%%%%%%%%%%%%%%%%%%%%%%%%%%%%%%%%%%%%%%



%%%%%%%%%%%%%%%%%%%%%%%%%% Def equal
\newcommand{\jointVel}[2]{\twistVel^{#1}_{J#2}}%
\newcommand{\jointVelDef}[2]{\spatialVel{#1}{#2}}%
\newcommand{\jointAcc}[2]{\dot{\twistVel}^{#1}_{J#2}}%
\newcommand{\jointAccDef}[2]{\spatialAcc{#1}{#2}}%


% In Body coordinates: #1 Joint index
% *predecessor idx: #1 -1  *successor #1
\newcommand{\jointVelBody}[1]{
  \textcolor{bodyVelColor}{
    \twistVel_{J#1}
  }
}
\newcommand{\jointVelBodynew}[2]{
  \textcolor{bodyVelColor}{
    \twistVel^{#1}_{J#2}
  }
}
% Defines the Joint frame number. The joint locates on the link #1-1 and does not depend on the #1th joint position
% #1 Joint index
\newcommand{\jointFrameDef}[1]{
  (#1 -1, #1)
}
% In Body coordinates: #1 Joint index
% *predecessor idx: #1 -1  *successor #1
\newcommand{\jointVelBodyDef}[1]{
  \bodyVel{\jointFrameDef{#1}}{#1}
  % \bodyVel{(#1-1, #1)}{#1}
}
% In Body coordinates: #1 Reference Frame  #2 Joint index
% *predecessor idx: #2 -1  *successor #2
% Due to the rigid link assumption, this formula assumes \bodyVel{(#2-1)}{(#2-1, #2)} = 0.
\newcommand{\jointVelBodyDefDef}[2]{
  \bodyVel{#1}{#2} - \twistTransform{(#2-1)}{#2}\bodyVel{#1}{(#2-1)}
}

% In Body coordinates: #1 Joint index
% *predecessor idx: #1 -1  *successor #1
\newcommand{\jointAccBody}[1]{
  \textcolor{bodyVelColor}{
    \dot{\twistVel}_{J#1}
  }
}
% In Body coordinates: #1 Joint index
% *predecessor idx: #1 -1  *successor #1
\newcommand{\jointAccBodyDef}[2]{
  \bodyAcc{\jointFrameDef{#1}}{#1}
}


% #1 Spatial frame or body frame, #2 joint index.
\newcommand{\linkVel}[2]{\twistVel^{#1}_{L#2}}%
\newcommand{\linkVelBody}[2]{
  \textcolor{bodyVelColor}{
    \twistVel^{#1}_{L#2}
  }
}%
\newcommand{\linkAcc}[2]{\dot{\twistVel}^{#1}_{L#2}}%

% The force transmits across a joint.
% #1 Reference frame, #2 joint index.
\newcommand{\body}[1]{#1^{b}}
\newcommand{\spatial}[1]{#1^{s}}

% Joint positions are unified via \jangles[overhead][superscript][type][postscript]:
%   \jangles             -> \bs{q}              (default, joint position)
%   \jangles[dot]        -> \dot{\bs{q}}
%   \jangles[ddot]       -> \ddot{\bs{q}}
%   \jangles[][][actuated]    -> \bs{\theta}    (actuated joint angles)
%   \jangles[][][generalized] -> \bs{\mathcal{X}} (floating-base generalized coordinates)
%   \jangles[][][fb]     -> \bs{q}_{\text{fb}}  (floating-base joint position)
% Simple symbol aliases for use with subscripts in equations:
\newcommand{\jointPosition}{\bs{\theta}}
\newcommand{\jointPositiond}{\dot{\jointPosition}}
\newcommand{\jointPositiondd}{\ddot{\jointPosition}}

% ################## Coordinate frames

% The frame defined by Christian Ott:
\newcommand{\jointAngle}{\theta}
\newcommand{\jointAngled}{\dot{\jointAngle}}
% \jangles[][][actuated] is now unified as \jangles[][][actuated]

%%%%%%%%%%%%%%%%%%%%%%%%%%%%%%%%%%%%%%%%%%%%%%%%%%%%%%%%%%%%%%%%%%%%%%%%%%%%%%%%
%% 12. JACOBIANS
%%%%%%%%%%%%%%%%%%%%%%%%%%%%%%%%%%%%%%%%%%%%%%%%%%%%%%%%%%%%%%%%%%%%%%%%%%%%%%%%

\newcommand{\applyJacobianSubscript}[1]{%
    \IfStrEqCase{#1}{%
    {contact}{
    \contactPoint
    }%
    }[#1] % "default" string
}

\newcommand{\jacobianSymbol}{\text{J}}
\NewDocumentCommand{\jacobian}{O{}O{}}{%
{    \ensuremath{%
        % Check the first optional argument for dot or double dot
        \applyOverheadSymbols{#1}{\jacobianSymbol}
        % 'contact',
        _{\applyJacobianSubscript{#2}}%
    }
    }
}



\newcommand{\jacobianVector}{\bs{\jacobian}}

\newcommand{\jac}[1]{\jacobian_{#1}} % Jacobian of a special quantity
\newcommand{\jacd}[1]{\dot{\jacobian}_{#1}} % Jacobian of a special quantity
\newcommand{\jacdd}[1]{\ddot{\jacobian}_{#1}} % Jacobian of a special quantity

\newcommand{\jacobianMotion}{\jacobian_{\text{m}}}%


%%%%%%%%%%%%%%%% The articulated-body inertia
% Matrix spans a subspace
\newcommand{\rVelJac}[2]{{
  \textcolor{bodyVelColor}{
    \jacobian^{'\bodyFrame}_{#1#2}
  }
}}

% The relative velocity Jacobian: vel = vel(induced by cmm vel) + relative vel
% #1 inertial frame #2 body frame
\newcommand{\rVel}[2]{{
  \textcolor{bodyVelColor}{
    \vc{V}^{'\bodyFrame}_{#1#2}
  }
}}
% #1 inertial frame #2 body frame
\newcommand{\rVelJacDef}[2]{{
    (\bodyJacobian{#1}{#2} - \twistTransform{\cmmFrame}{#2}\bodyJacobian{#1}{\cmmFrame})
}}



\newcommand{\appBodyJacobian}[2]{
  \textcolor{bodyVelColor}{
    {\tilde{{\jacobian}}^{\bodyFrame\text{t}}_{#1 #2}}
  }
}
\newcommand{\tBodyJacobian}[2]{
  \textcolor{bodyVelColor}{
    {\jacobian^{\bodyFrame\text{t}}_{#1 #2}}
  }
}
\newcommand{\rBodyJacobian}[2]{
  \textcolor{bodyVelColor}{
    {\jacobian^{\bodyFrame\text{R}}_{#1 #2}}
  }
}

%  Define a macro with an optional argument for dot, or more.
\newcommand{\bodyJacobian}[3][]{
  \textcolor{bodyVelColor}{
      #1{\jacobian}^{\bodyFrame}_{#2 #3}
  }
}

\newcommand{\bodyJacobiand}[2]{
  \textcolor{bodyVelColor}{
    \dot{\jacobian}^{\bodyFrame}_{#1 #2}
  }
}


%%%%%%%%%%%%%%%%%%%%%%%%%% QP Control law


%%%%%%%%%%%%%%%%%%%%%%%%%%%%%%%%%%%%%%%%%%%%%%%%%%%%%%%%%%%%%%%%%%%%%%%%%%%%%%%%
%% 13. INERTIA, MASS, AND MOMENTUM
%%%%%%%%%%%%%%%%%%%%%%%%%%%%%%%%%%%%%%%%%%%%%%%%%%%%%%%%%%%%%%%%%%%%%%%%%%%%%%%%

\newcommand{\smm}{
  \textcolor{bodyVelColor}{
    \bs{A}
  }
}

\newcommand{\systemMomenta}{
  \textcolor{bodyVelColor}{
    \bs{h}
  }
}
\newcommand{\momenta}{
  \bs{h}
}

% The momentum transform looks the same as the wrench transform
%%% #1:current frame, #2:target frame
\newcommand{\geometricMT}[2]{
  \geometricFT{#1}{#2}
}
\newcommand{\geometricMTTwo}[2]{
  \geometricFTTwo{#1}{#2}
}
\newcommand{\systemVel}{
  \textcolor{bodyVelColor}{
    \bs{V}
  }
}
\newcommand{\systemJacobian}{
  \textcolor{bodyVelColor}{
   \jacobian
  }
}
\newcommand{\systemJacobiand}{
  \textcolor{bodyVelColor}{
   \jacobian[dot]
  }
}

\newcommand{\momentum}{\bs{h}}
\newcommand{\momentumd}{\dot{\bs{h}}}

%%% #1 reference frame, #2 link index
\newcommand{\momentumFrame}[2]{\bs{h}_{#1#2}}

%%% The canonical momenta:
\newcommand{\canoMomenta}{{\bs{h}_{J}}}

%%% The system momentum per link
\newcommand{\sMomentum}{{\bs{h}_{s}}}

%%% The contact momentum per contact
\newcommand{\cMomentum}{{\bs{h}_{c}}}

% #1:reference frame, #2:contact index
\newcommand{\cMomentumDef}[2]{
  \equivalentMass_{#2} \spatialVel{#1}{#2}
}

\newcommand{\cMomentumDefDef}[2]{
  {\equivalentMassDef{#1}{#2}} \spatialVel{#1}{#2}
}


\newcommand{\cmmMomenta}{
  \textcolor{bodyVelColor}{
  \bs{h}_{\cmmFrame}
 }
}
\newcommand{\cmmInertia}{
  \textcolor{bodyVelColor}{
    \bs{I}_{\cmmFrame}
  }
}

\newcommand{\cmmInertiaDef}{
  \agg{i}{\dof}{\spatialGInertia{\cmmFrame}{i}}
}
\newcommand{\cmmRelativeInertia}{
  \textcolor{bodyVelColor}{
    \bs{I}'_{\cmmFrame}
  }
}
\newcommand{\cmmRelativeInertiaDef}{
  \textcolor{bodyVelColor}{
    \quadraticObj{\rVelJac{}{}}{\systemInertia}
  }
}

\newcommand{\cmmInertiaDefTwo}{
  \agg{i}{\dof}{\spatialGInertia{\cmmFrame}{i}}
}

\newcommand{\cmmInertiaDetail}{
  \textcolor{bodyVelColor}{
    \bs{I}_{\cmmFrame}
  }
}

\newcommand{\linearMomentumSymbol}{P}
\NewDocumentCommand{\linearMomentum}{O{}O{}O{}}{%
{    \ensuremath{%
        % Check the first optional argument for dot or double dot
        \applyOverheadSymbols{#1}{\linearMomentumSymbol}
        % Check the second optional argument for 'ref'
        ^{\applySuperscript{#2}}
        \applyPostscript{#3}
    }
    }
}
%  The angular Momentum
\newcommand{\angularMomentumSymbol}{\textcolor{textColor}{\mathcal{L}}}
\NewDocumentCommand{\angularMomentum}{O{}O{}O{}O{}}{%
{    \ensuremath{%
        % Check the first optional argument for dot or double dot
        \applyOverheadSymbols{#1}{\angularMomentumSymbol}
        % Check the second optional argument for 'ref'
        _{#2}
        ^{\applySuperscript{#3}}
        \applyPostscript{#4}
    }
    }
}



%  The linear Momentum
\newcommand{\aMomentum}{\mathcal{L}}

\newcommand{\cmmSymbol}{
\textcolor{bodyVelColor}{
A
}
}




\newcommand{\applyCmmSuperscript}[1]{%
\textcolor{bodyVelColor}{
    \IfStrEqCase{#1}{%
    {l}{
    \text{t}
    }%
    {r}{
    \text{r}
    }%
    }[] % "default" string
    }
}


% %%%%%%%%%%%%%%%%%%%%%%%%% EQUATION VARIABLES:
\NewDocumentCommand{\cmm}{O{}O{default}O{default}O{default}}{%
{    \ensuremath{%
        % Check the first optional argument for dot or double dot
        \applyOverheadSymbols{#1}
        {\cmmSymbol}_{\textcolor{bodyVelColor}{\cmmFrame}}
        % Check the second optional argument for linear or angular
        ^{\applyCmmSuperscript{#2}}%
         \applyPostscript{#3}
    }
    }
}


\newcommand{\cmmt}{
  \textcolor{bodyVelColor}{
    A^{\text{t}}_{\cmmFrame}(\jangles)
  }
}


\newcommand{\cmmr}{
  \textcolor{bodyVelColor}{
    A^{\text{r}}_{\cmmFrame}(\jangles)
  }
}

\newcommand{\cmmLinear}{A_{\bs{v}\cmmFrame}}

\newcommand{\cmmDot}{
  \textcolor{bodyVelColor}{
  \dot{A}_{\cmmFrame}(\jangles[][][actuated])
}
}

\newcommand{\cmmPlane}{
  \textcolor{bodyVelColor}{
    A_{\cmmFrame_{x,y}}
  }
}

\newcommand{\aii}{\bs{\Phi}}
\newcommand{\aiiDef}[2]{#1 \inverse{(\transpose{#1} #2 #1)}\transpose{#1}}
% The joint acc of a rigid body constrianed by a joint.
% #1 the joint jacobian, #2 the inertia
\newcommand{\aInertia}{{I^A}}%

%%%%%%%%%%%%%%%% Motion constraint:

% The implicit constraint: \Phi(\dot{q}) = 0
% #1 Joint position
\newcommand{\fbVel}{
  \bs{V}_{\fbFrame}
}
\newcommand{\fbAcc}{
  \dot{\bs{V}}_{\fbFrame}
}

\newcommand{\cmmVel}{
  \textcolor{bodyVelColor}{
    \bs{V}_{\cmmFrame}
  }
}

\newcommand{\cmmVelDef}{
  \bodyVel{\inertialFrame}{\cmmFrame}
}
\newcommand{\cmmVelDefDef}{
  \inverse{\cmmInertia} \cmmMomenta
}
%%%%%%%%%%%% Exponential Map given unit skew matrix and a real number
%%% #1:input matrix
\newcommand{\crbVel}{
  {^{crb}\bs{V}}
}

%%% Spatial CRB inertia
% #1 inertial frame #2 body frame
\newcommand{\spatialCrbInertia}[2]{
  {^{crb}I^s_{#1#2}}
}

%%% CRB inertia
\newcommand{\crbGInertia}{
  {^{crb}I}
}

%%% System inertia
\newcommand{\systemInertia}{
  \textcolor{bodyVelColor}{
    I
  }
}


%%%%%%%%%%%

\newcommand{\averageVel}{
  \bar{\systemVel}
}
\newcommand{\relativeVel}{{\systemVel'}}
% The relative velocity Jacobian: vel = vel(induced by cmm vel) + relative vel
% #1 inertial frame #2 body frame
\newcommand{\equivalentMass}{{M_{eq}}}
% The 6-by-6 equivalent inertensor
\newcommand{\eInertiaTensor}{{I_{eq}}}

\newcommand{\osInertia}{{\Lambda}}
% The operational space bias force
\newcommand{\osBias}{{\bs{p}}}

% #1:reference frame #2:contact index
\newcommand{\equivalentMassDef}[2]{
  \inverse{\left(
      \spatialJacobian{#1}{#2} \inverse{\inertiaMatrix}_{#2}
     \transpose{{\spatialJacobian{#1}{#2}}}
   \right)}
}

\newcommand{\geometricFTDef}[2]{
  \textcolor{geometricColorVariable}{
    \adgInvTransDef{#2}{#1}
  }
}

% From frame #1 to frame #2
\newcommand{\geometricFTTwo}[2]{
  \textcolor{geometricColorVariable}{
    \adgTrans{#1}{#2}
  }
}
\newcommand{\geometricFTTwoDef}[2]{
  \textcolor{geometricColorVariable}{
    \adgTransDef{#1}{#2}
  }
}

% From frame #1 to frame #2
\newcommand{\geometricFT}[2]{
  \textcolor{geometricColorVariable}{
    \adgInvTrans{#2}{#1}
  }
}

% From frame #1 to frame #2
\newcommand{\geometricFTd}[2]{
  \textcolor{geometricColorVariable}{
    \derivative{\adgInvTrans{#2}{#1}}
  }
}

% From frame #1 to frame #2
\newcommand{\geometricFTdDef}[2]{
  \textcolor{geometricColorVariable}{
    -\transpose{\cross{\spatialVel{#2}{#1}}{}}\geometricFT{#1}{#2}
  }
}
\newcommand{\geometricFTdDefTwo}[2]{
  \textcolor{geometricColorVariable}{
    \dualCross{\spatialVel{#2}{#1}}{}\geometricFT{#1}{#2}
  }
}
\newcommand{\geometricFTdDefThree}[2]{
  \textcolor{geometricColorVariable}{
        \transpose{(\bodyVel{#1}{#2}\times)} \transpose{\twistTransform{#2}{#1}}
  }
}

\newcommand{\bodyGInertia}[2]{
\textcolor{bodyVelColor}{
  \body{\gInertia}_{#1 #2}
  }
}


\newcommand{\mechanicalConnection}{
  \bs{A}
}




\newcommand{\bodyGInertiaDef}[2]{
  \textcolor{bodyVelColor}{
    \matrixTwo{\mass_{#2} \identityMatrix }{\zeroMatrix}{\zeroMatrix}{\bodyMetaInertia{#1}{#2}}
  }
}

\newcommand{\spatialGInertia}[2]{
  \textcolor{spatialVelColor}{
    \spatial{\gInertia}_{#1#2}
  }
}

\newcommand{\spatialGInertiad}[2]{
  \textcolor{spatialVelColor}{
    \spatial{\dot{\gInertia}}_{#1#2}
  }
}

\newcommand{\spatialGInertiadDef}[2]{
  \textcolor{spatialVelColor}{
    % -\transpose{(\cross{\spatialVel{#1}{#2}}{})} \spatialGInertia{#1}{#2} - \spatialGInertia{#1}{#2}\cross{\spatialVel{#1}{#2}}{}
    \dualCrossDef{\spatialVel{#1}{#2}}{\spatialGInertia{#1}{#2}}
    - \spatialGInertia{#1}{#2}\cross{\spatialVel{#1}{#2}}{}
  }
}

\newcommand{\spatialGInertiadDefTwo}[2]{
  \textcolor{spatialVelColor}{
    \dualCross{\spatialVel{#1}{#2}}{\spatialGInertia{#1}{#2}}
    - \spatialGInertia{#1}{#2}\cross{\spatialVel{#1}{#2}}{}
  }
}

% #1 Current reference frame #2 New Reference frame
\newcommand{\momentaTransform}[2]{
  \transpose{\twistTransform{#2}{#1}}
}
\newcommand{\momentaTransformDef}[2]{
  \transpose{\twistTransformDef{#2}{#1}}
}

% #1 Current reference frame #2 New Reference frame
\newcommand{\spatialGInertiaTransform}[3]{
  \transpose{\twistTransform{#3}{#1}}\spatialGInertia{#1}{#2}\twistTransform{#3}{#1}
}

% #1 and #2 defines the spatial generalized inertia to be transformed.
% #1 Reference frame #2 body frame
% #3 defines the new reference frame
\newcommand{\spatialCrbInertiaTransform}[3]{
  \transpose{\twistTransform{#3}{#1}}\spatialCrbInertia{#1}{#2}\twistTransform{#3}{#1}
}

\newcommand{\gInertiaTransform}[3]{
  \transpose{\twistTransform{#1}{#2}}
  #3
  \twistTransform{#1}{#2}
}

\newcommand{\spatialGInertiaDef}[2]{
  \transpose{\twistTransform{#1}{#2}}
  % \bodyGInertia{#1}{#2}
  \spatialGInertia{#2}{\com_{#2}}
  \twistTransform{#1}{#2}
}

\newcommand{\spatialGInertiaDefTwo}[2]{
  \transpose{\twistTransformTwo{#1}{#2}}
  % \bodyGInertia{#1}{#2}
  \spatialGInertia{#2}{\com_{#2}}
  \twistTransformTwo{#1}{#2}
}
% #1 Reference frame #2 body frame
\newcommand{\spatialGInertiaDefDef}[2]{
  % \transpose{{\adgInvDef{#1}{#2}}} \bodyGInertiaDef{#2}\adgInvDef{#1}{#2}
  \transpose{{\twistTransformDef{#1}{#2}}}
  \spatialGInertia{#2}{\com_{#2}}
  % \bodyGInertiaDef{#1}{#2}
  \twistTransformDef{#1}{#2}
}
\newcommand{\spatialGInertiaDefDefTwo}[2]{
  % \transpose{{\adgInvDef{#1}{#2}}} \bodyGInertiaDef{#2}\adgInvDef{#1}{#2}
  \transpose{{\twistTransformTwoDef{#1}{#2}}}
  % \bodyGInertiaDef{#1}{#2}
  \spatialGInertia{#2}{\com_{#2}}
  \twistTransformTwoDef{#1}{#2}
}

% #1 Reference frame #2 body frame
\newcommand{\spatialGInertiaDetail}[2]{
  \textcolor{spatialVelColor}{
\matrixTwo{
  \mass_{#2} \identityMatrix
}{
  \mass_{#2}\transpose{\translationSkew{#1}{#2}}
}{
  \mass\translationSkew{#1}{#2}
}{
  \mass \translationSkew{#1}{#2} \transpose{\translationSkew{#1}{#2}} + \rotation{#1}{#2}\bodyMetaInertia{#1}{#2}\rotationInv{#1}{#2}
}
}
}
% #1 Reference frame #2 body frame
\newcommand{\spatialGInertiaDiagonalDetail}[2]{
  \textcolor{spatialVelColor}{
\matrixTwo{
  \mass_{#2} \identityMatrix
}{
  \zeroMatrix
}{
  \zeroMatrix
}{
  \rotation{#1}{#2}\bodyMetaInertia{#1}{#2}\rotationInv{#1}{#2}
}
}
}

% #1 Reference frame #2 body frame
\newcommand{\metaInertia}{\mathcal{I}}



% #1: the reference frame and #2 the link frame
\newcommand{\spatialMetaInertia}[2]
{
\textcolor{spatialVelColor}{
  \spatial{\metaInertia}_{{#1}{#2}}
}
}
% #1: the reference frame and #2 the link frame
\newcommand{\bodyMetaInertia}[2]
{
  \textcolor{bodyVelColor}{
      \body{\metaInertia}_{{#1 #2}}
    }
}
\newcommand{\spatialInertia}{\mathcal{I}^{\prime}}

\newcommand{\gInertia}{\mathcal{M}}

%%%%%%%%%%%%%%%%%%%%%%%%%%%%%%%%%%%%%%%%%%%%%%%%%%%%%%%%%%%%%%%%%%%%%%%%%%%%%%%%
%% 14. EQUATIONS OF MOTION
%%%%%%%%%%%%%%%%%%%%%%%%%%%%%%%%%%%%%%%%%%%%%%%%%%%%%%%%%%%%%%%%%%%%%%%%%%%%%%%%

\newcommand{\numArticulatedJoints}{n}
\newcommand{\nAJ}{n}
\newcommand{\ajDim}{\text{n}}
\newcommand{\actuationMatrix}{\text{B}}

\newcommand{\numLinks}{{n_{L}}}
\newcommand{\fbDim}{6}
\newcommand{\jsDim}{(\numArticulatedJoints+\fbDim)}
\newcommand{\selectionMatrix}{
  \textcolor{textColor}{
  \text{S}
  }
}

\newcommand{\subspace}{S}

% The apparent inverse inertia. #1 subspace #2 inertia
\newcommand{\jaccConstrained}[2]{
  \inverse{(\transpose{#1} #2 #1)}\transpose{#1}(\jaWrench{}{} - #2\dot{#1}\jvelocities - \biasForce )
}

% #1 Spatial frame, #2 joint index.
\newcommand{\cwBasis}{T}%
% The orthonormal basis of the Actuation wrench. T_a \in \RRm{6}{1}
\newcommand{\awBasis}{{T_a}}%
% The motion subspace:
\newcommand{\mBasis}{{S}}%


% The projection matrix that describes the connections between different links in a multi-body system
% The connectivity matrix
\newcommand{\pMatrix}{
  \textcolor{spatialVelColor}{
    P
  }
}%

\newcommand{\pMatrixBody}{
  \textcolor{bodyVelColor}{
    P
  }
}%


% The articulated-body inertia
\newcommand{\ic}[1]{\Phi(#1)}


% The constraint kernel: K \dot{q} = 0
\newcommand{\ick}{K}
\newcommand{\ickd}{\dot{\ick}}
%%% #1 Joint position
\newcommand{\ickDef}[1]{\pd{\ic{#1}}{#1}}

% The explicit constraint:  q = \gamma(x)
% #1 End-effector position
\newcommand{\ec}[1]{\gamma(#1)}

% The explicit constraint kernel:  \dot{q} = \eck\dot{x}
% #1 End-effector position
\newcommand{\eck}{G}
\newcommand{\eckDef}[1]{\pd{\ec{#1}}{#1}}

\newcommand{\eckd}{\dot{\eck}}


\newcommand{\cForce}{\bs{\lambda}}

\newcommand{\cTorque}{\jtorques_c}

% Dimension of the constriants
\newcommand{\cDim}{m}
\newcommand{\cRank}{m_{ic}}


%%%%%%%%%%%%%%%%


% Contact Vel:
\newcommand{\gravitySymbol}{\mathbf{g}}
\newcommand{\gravity}[1][default]{
    \IfStrEqCase{#1}{%
        {value}{9.81}%
        {vec}{\vectorThree{0}{0}{-1g}}%
    }[\gravitySymbol] % "default" string
}

% Operational space graviational forces
\newcommand{\sDim}{\text{n}}
% input dimension
\newcommand{\iDim}{\text{m}}


% First-order partial derivative #1 quantity, #2 variable
\newcommand{\gForce}{{\text{\bf g} }}
\newcommand{\gForceDef}{\vectorThree{0}{0}{-1}g}
\newcommand{\gValue}{\text{9.81}}
\newcommand{\gSymbol}{\text{g}}




\newcommand{\dof}{n}
\newcommand{\jsim}{\inertiaMatrix}
\newcommand{\inertiaMatrix}{\text{M}}
\newcommand{\inertiaMatrixFull}{\text{M}(\jangles)}
\newcommand{\gravityandcoriolis}{\mathbf{N}}
\newcommand{\gravityandcoriolisFull}{\gravityandcoriolis(\jangles, \jangles[dot])}
\newcommand{\coriolis}{\mathbf{C}}

\newcommand{\inertiaMatrixd}{\dot{\inertiaMatrix}}

%%% #1:row index, #2:column index
\newcommand{\inertiaMatrixComponentDef}[2]{
\sum^\dof_{l = max(#1,#2)}{
  \transpose{
    \twistC{S}{#1}(t_0)
  }\transpose{
    \twistTransform{#1}{l}
  }
  \spatialGInertia{S}{\com_l}(t_0)
  \twistTransform{#2}{l}\twistC{S}{#2}(t_0)
}
}

\newcommand{\inertiaMatrixComponentDefTwo}[2]{
\sum^\dof_{l = max(#1,#2)}{
  \transpose{
    \twistC{S}{#1}(t_0)
  }\transpose{
    \adg{l}{#1}
  }
  \spatialGInertia{S}{\com_l}(t_0)
  \adg{l}{#2}
  \twistC{S}{#2}(t_0)
}
}

%%%%%%%%%%%%%%%%%% inertiaMatrix partial derivative:
%%% #1:row index, #2:column index, #3 theta index
\newcommand{\inertiaMatrixdComponentDef}[3]{
  \sum^\dof_{l = max(#1,#2)}{\left(
      \transpose{
        \liebra{\adg{#3}{#1}\twistC{S}{#1}}{\twistC{S}{#3}}
      }
      \transpose{
        \adg{l}{#3}
      }
      \spatialGInertia{S}{\com_l}(t_0)
      \adg{l}{#2}
      \twistC{S}{#2}(t_0)
      +
      \transpose{
        \twistC{S}{#1}(t_0)
      }\transpose{
        \adg{l}{#1}
      }
      \spatialGInertia{S}{\com_l}(t_0)
      \adg{l}{#3}\liebra{\adg{#3}{#2}\twistC{S}{#2}}{\twistC{S}{#3}}.
      % \adg{l}{#2}
      % \twistC{S}{#2}(t_0)
    \right)
  }
}



\newcommand{\cMatrix}{C(\jointPosition, \jointPositiond)}
\newcommand{\nVector}{N}
\newcommand{\rpTransform}{\mathcal{T}}
\newcommand{\rpTransformd}[1]{\mathcal{#1}}

\newcommand{\coordinateTransformed}[2]{
\transpose{{\inverse{#2}}} #1 \inverse{#2}
}



\newcommand{\cristoffelSymbol}[3]{\Gamma_{#1#2#3}}

%%% #1 Reference frame: #2 notation
\newcommand{\cMatrixComponent}[3]{
  \Gamma_{#1#2#3}
}
%%%%%%%%%%% #1:row index, #2:column index, #3 index (between 1 and dof)
\newcommand{\cMatrixComponentDef}[3]{
  \frac{1}{2}\left(
    \pd{\inertiaMatrix_{#1#2}}{\jointPosition_{#3}}
    +\pd{\inertiaMatrix_{#1#3}}{\jointPosition_{#2}}
    -\pd{\inertiaMatrix_{#3#2}}{\jointPosition_{#1}}
  \right)
}



% #1:reference frame, #2:notation of the twist
\newcommand{\biasForce}{\bs{p}}

\newcommand{\kineticEnergy}{T(\jointPosition, \jointPositiond)}
\newcommand{\potentialEnergy}{V(\jointPosition)}
\newcommand{\totalEnergy}{H}

\newcommand{\lagR}{\mathcal{L}(\jointPosition, \jointPositiond)}

\newcommand{\lagRDef}{
  \quadratic{\jointPositiond}{\inertiaMatrix} + \potentialEnergy
}

%%% #1: dof number
\newcommand{\lagRDefDef}[1]{
  \frac{1}{2}\sum^{#1}_{i,j=1}\inertiaMatrix_{ij}(\jointPosition)\jointAngled_i\jointAngled_j - \potentialEnergy
}

% The Lagrange multipliers
\newcommand{\lm}{\bs{\lambda}}

%%% #1: link index
\newcommand{\lagREquations}[1]{
  \derivative{\pd{\lagR}{\jointAngled_{#1}}} - \pd{\lagR}{\jointPosition_{#1}} = \torque_{#1}
}

% The Lagrange's equation of a link
%%% #1: link index
\newcommand{\lagREquationsDef}[1]{
  \agg{j}{\dof}{
  (\inertiaMatrixd_{#1j}\jointPositiond_j  + \inertiaMatrix_{#1j}\jointPositiondd_j)
}
- \frac{1}{2}\agg{#1}{\dof}
{
  \pd{\inertiaMatrix_{#1j}}{\jointPosition_#1} \jointPositiond_#1\jointPositiond_j
  - \pd{\potentialEnergy}{\jointPosition_#1}
}
  = \torque_{#1}
}


\newcommand{\lagREquationsDefDef}[1]{
  \agg{j}{\dof}{
    \inertiaMatrix_{#1j}\jointPositiondd_j
  }
  + \frac{1}{2}\agg{j,k}{\dof}{(
    \pd{\inertiaMatrix_{#1j}}{\jointPosition_k}
    +\pd{\inertiaMatrix_{#1k}}{\jointPosition_j}
    - \pd{\inertiaMatrix_{kj}}{\jointPosition_#1}
    % \christoffelSymbolDef{#1}{j}{k}
  )\jointPositiond_k\jointPositiond_j}
+ \pd{\potentialEnergy}{\jointPosition_#1}
  = \torque_{#1}
}
%%%% #1 link index #2,#3 iteration index
\newcommand{\christoffelSymbol}[3]{
  \Gamma_{#1#2#3}
}
\newcommand{\christoffelSymbolDef}[3]{
  \frac{1}{2}(
    \pd{\inertiaMatrix_{#1#2}}{\jointPosition_#3}
    +\pd{\inertiaMatrix_{#1#3}}{\jointPosition_#2}
    - \pd{\inertiaMatrix_{#3#2}}{\jointPosition_#1})
}


% The argument are: force, torque

%%%%%%%%%%%%%%%%%%%%%%%%%%%%%%%%%%%%%%%%%%%%%%%%%%%%%%%%%%%%%%%%%%%%%%%%%%%%%%%%
%% 15. WRENCHES, FORCES, AND TORQUES
%%%%%%%%%%%%%%%%%%%%%%%%%%%%%%%%%%%%%%%%%%%%%%%%%%%%%%%%%%%%%%%%%%%%%%%%%%%%%%%%

\newcommand{\jointForce}[2]{\force^{#1}_{J#2}}%
\newcommand{\jointWrench}[2]{\wrench^{#1}_{J#2}}%

% #1 Reference frame, #2 joint index.
\newcommand{\linkWrench}[2]{\wrench^{#1}_{L#2}}%

% The actuation force of a joint
% #1 Reference frame, #2 joint index.
\newcommand{\jaForce}[2]{\force^{#1}_{a#2}}%
\newcommand{\jaWrench}[2]{\wrench^{#1}_{a#2}}%
% #1 joint index.
\newcommand{\jaWrenchDef}[1]{\wrench^\jointFrameDef{#1}_{a\jointFrameDef{#1}}}%


% The constraint force implements the motion constriant.
% #1 Reference frame, #2 joint index.
\newcommand{\jcForce}[2]{\force^{#1}_{c#2}}%
\newcommand{\jcWrench}[2]{\wrench^{#1}_{c#2}}%
% #1 joint index.
\newcommand{\jcWrenchDef}[1]{\wrench^{\jointFrameDef{#1}}_{c\jointFrameDef{#1}}}%

% The orthonormal basis of the constraint wrench. For example, T \in \RRm{6}{5} for a joint. and T \in \RRm{6}{1} for a point contact.
\newcommand{\inertialWrench}{{\wrench_{\inertialFrame}}}
\newcommand{\inertialForce}{{\force_{\inertialFrame}}}
\newcommand{\inertialTorque}{{\torque_{\inertialFrame}}}

\newcommand{\comForce}{{\force_{\com}}}
\newcommand{\comTorque}{{\torque_{\com}}}


% Z direction of the inertial frame
\newcommand{\cwc}{\bs{C}}

\newcommand{\wrenchC}[2]{
\begin{bmatrix}
#1\\
#2
\end{bmatrix}
}

% The argument are: force, torque
\newcommand{\wrenchR}[2]{
[#1^\top, #2^\top]^\top
}


%%%%%%%%%%%%%%%%%%%%%%%%%%%%%%%%%%%%%%%%%%%%%%%%%%%%%%%%%%%%%%%%%%%%%%%%%%%%%%%%
%% 16. CONTACT, FRICTION, AND IMPACT DYNAMICS
%%%%%%%%%%%%%%%%%%%%%%%%%%%%%%%%%%%%%%%%%%%%%%%%%%%%%%%%%%%%%%%%%%%%%%%%%%%%%%%%

\newcommand{\impactDuration}{\delta t}



\newcommand{\nextState}[1]{{#1}_{t_{k+1}}}
\newcommand{\nextStateB}[2]{{#1}_{#2,t_{k+1}}}
\newcommand{\nextStateTwo}[1]{{#1(t_{k+1})}}

\newcommand{\jump}{\Delta}
\newcommand{\impactduration}{\delta t}
\newcommand{\pointer}[2]{\overrightarrow{{#1}{#2}}}
\newcommand{\pointerDef}[2]{{\bs{#2}  - \bs{#1}}}

\newcommand{\pointerPerp}[2]{\overrightarrow{#1#2}_{\perp}}


\newcommand{\work}[1]{w_{#1}} % work done by #1
\newcommand{\impulse}{\bs{\iota}}
\newcommand{\impulseScalar}{\iota}

\newcommand{\quantityJump}{\jump \quantity}
\newcommand{\jtorquesJump}{\jump \jtorques}
\newcommand{\eevelocity}{\bs{\dot{x}}}
\newcommand{\eevelocityJump}{\jump \eevelocity}
\newcommand{\forceJump}{\jump \force}

\newcommand{\initial}[1]{{#1}_{0}}
\newcommand{\viable}[1]{{#1}^{\viableSymbol}}

\newcommand{\current}[1]{{#1}_{t_{k}}}
\newcommand{\currentTwo}[1]{#1(t_{k})}
\newcommand{\nextTwo}[1]{{#1(t_{k+1})}}
\newcommand{\previous}[1]{{#1}_{t_{k-1}}}

\newcommand{\impactNormal}{\bs{n}}

\newcommand{\surfaceNormal}{\bs{n}}
\newcommand{\projection}{P}

\newcommand{\projectionDef}[1]{{#1}{#1}^{\top}}

%%%%%%%%%%%%%%%%%%%%%%%%%%% C++ code

\newcommand{\lcpMass}{\text{M}_{\text{lcp}}}
\newcommand{\lcpd}{\text{d}_{\text{lcp}}}
% separation velocity:
\newcommand{\seVelocity}{\dot{\bs{\xi}}}
\newcommand{\seDistance}{\bs{\xi}}

%%%%%%%%%%%%%%%%%%%%%%%%%% Impact dynamics:

%%%%% Impact duration and structural dynamics

% The critical impact velocity that leads to visible plastic indentation.
\newcommand{\icVel}{\contactVel_{cr}}

% The yielding impact velocity that leads to plastic deformations that are not visible.
\newcommand{\iyVel}{\contactVel_{Y}}

%%%%%

\newcommand{\contactVel}{\bs{v}}
\newcommand{\nContactVel}{v_n}
\newcommand{\nContactVelR}{\contactVel_{nr}}
\newcommand{\nContactVelS}{\contactVel_{ns}}
\newcommand{\nContactVelL}{\contactVel_{nl}}
\newcommand{\nContactVelF}{\contactVel_{n,f}}
% The point where the tangential contactvel and its derivative are collinear.
\newcommand{\tVelL}{\contactVel_{gl}}
\newcommand{\tVelZero}{{\tp{\contactVel}}_{0}}

\newcommand{\cln}{\bs{s}}

% notation of the impacting body
\newcommand{\ib}{1}

% notation of the object
\newcommand{\ob}{2}

% virtual spring length
\newcommand{\vSpringL}{x}
\newcommand{\vSpringLd}{\dot{x}}
% spring coefficient
\newcommand{\vSpringC}{k}
\newcommand{\vDamperC}{c}

% The potential energy stored in the virtual spring
\newcommand{\pEnergy}{{{E}_p}}
\newcommand{\dEnergy}{{{E}_d}}
\newcommand{\dEnergyd}{{\dot{{E}}_d}}

% The kinetic energy
\newcommand{\kEnergy}{{E_k}}
\newcommand{\kEnergyd}{{\dot{E}_k}}

% #1 denotes the normal impulse
\newcommand{\pEnergyIntOne}[1]{\Phi_1(#1)}
% #1 denotes the normal impulse
\newcommand{\pEnergyIntTwo}[1]{\Phi_2(#1)}

\newcommand{\pEnergyd}{\dot{E}}

% Impulse at the end of the impact process
\newcommand{\impulseEnd}{\impulse_r}

% Derivative w.r.t. the normal impulse
\newcommand{\dni}[1]{\frac{d #1}{d \nImpulse} }
\newcommand{\dzi}[1]{\frac{d #1}{d \zImpulse} }
% Inverse inertia matrix
\newcommand{\iim}{W}
% The tangential plane iim
\newcommand{\iimtp}{B}
\newcommand{\iimtpDef}{\vectorTwo{\transpose{\basisVec{x}}}{\transpose{\basisVec{y}}} \iim \columnTwo{\basisVec{x}}{\basisVec{y}}}
% The z-axis iim
\newcommand{\iimcross}{\bs{d}}
\newcommand{\iimcrossDef}{\vectorTwo{\transpose{\basisVec{x}}}{\transpose{\basisVec{y}}} \iim \basisVec{z}}

\NewDocumentCommand{\coef}{m}{%
\IfStrEqCase{#1}
{
{f}{\mu}% friction
{s}{\bar{\mu}}% Coefficient of Stick
{r}{c_{\text{r}}}% Restitution coefficient
{er}{e_{\text{r}}}% Energetic restitution coefficient
}[] % "default" string}
}

\newcommand{\postImpact}[1]{{#1}^{+}}
\newcommand{\preImpact}[1]{#1^{-}}

% The impulse when the restitution phase ends
\newcommand{\impulseR}{\impulse_{r}}
\newcommand{\contactVelR}{\contactVel_{r}}
\newcommand{\contactVelC}{\contactVel_{c}}
% The impulse when the collinear condition applies
\newcommand{\impulseL}{\impulse_{l}}

% tangential direction unit vector
\newcommand{\tangentialU}{\hat{\bs{t}}}

% normal direction unit vector
\newcommand{\normalU}{\hat{\bs{n}}}

% The  unit vector of a coordinate system:
\newcommand{\basisVec}[1]{\widehat{\bs{\bs{#1}}}}


% The tangential impulse
\newcommand{\tImpulse}{\impulseScalar_{t}}
\newcommand{\tpImpulse}{\tp{\impulse}}
\newcommand{\tImpulsed}{\dot{\impulseScalar}_{t}}
% The tangential force
\newcommand{\tForce}{\forceScalar_{t}}

% Denote the tangential plane:
\newcommand{\tp}[1]{{#1}_{\perp}}

\newcommand{\tpForce}{\tp{\force}}

% The tangential linear Vel
\newcommand{\tVel}{\linearVelScalar_{t}}
\newcommand{\tVels}{\bs{\linearVelScalar}_{t}}
\newcommand{\hodog}{\dni{\tVels}}

% Curvature of the hodograph
\newcommand{\hodogC}{\kappa}
\newcommand{\hodogStepSize}{h}

% The normal linear Vel
\newcommand{\nVel}{\linearVelScalar_{n}}

% The normal impulse
\newcommand{\nImpulse}{{\impulseScalar_{n}}}
\newcommand{\nImpulseIni}{{\impulseScalar_{n,0}}}
\newcommand{\nImpulseF}{{\impulseScalar_{n,f}}}
\newcommand{\zImpulse}{{\impulseScalar_{z}}}
\newcommand{\nImpulsed}{\dot{\impulseScalar}_{n}}
\newcommand{\zImpulsed}{\dot{\impulseScalar}_{z}}
\newcommand{\impulseCone}{{^{\impulse}K}}
\newcommand{\fConeNum}{N_{\coef{f}}}

\newcommand{\applyConeSubscript}[1]{%
    %% \IfStrEqCase{#1}{%
    %% {}{
    %% \coef[f]
    %% }%
    %% {contact}{
    %% \text{sc}
    %% }%
    %% }[] % "default" string
    #1
}
\newcommand{\applyConeSuperscript}[1]{%
    \IfStrEqCase{#1}{%
    {f}{ % Friction cone
    \coef{f}
    }%
    {cwc}{ % Contact wrench cone
    \text{cwc}
    }%
    }[] % "default" string
}

\newcommand{\coneSymbol}{\text{K}}
\NewDocumentCommand{\cone}{O{}O{}}{
{
\coneSymbol^{\applyConeSuperscript{#1}}_{\applyConeSubscript{#2}}}
}

\newcommand{\generator}[1]{{\bs{\lambda}_{#1}}}
\newcommand{\generatorScalar}[1]{{\lambda_{#1}}}

\newcommand{\frictionCone}{{^{\coef{f}}K}}
\newcommand{\normalCone}{{^{\surfaceNormal}K}}
\newcommand{\impulseGenerator}{\bs{\lambda}}


% The normal force
\newcommand{\nForce}{{\forceScalar_{n}}}


% The normal impulse when the compression phase ends
\newcommand{\impulseC}{\impulse_{c}}
\newcommand{\nImpulseC}{\impulseScalar_{nc}}
\newcommand{\zImpulseC}{\impulseScalar_{zc}}
\newcommand{\zImpulseLCE}{\zeta_{zc}}
\newcommand{\czImpulseR}{\zeta_{r}} % candidate of the \zImpulseR
\newcommand{\nImpulseCEini}{\tilde{\impulseScalar}_{nc}}

% The normal impulse when the restitution phase ends
\newcommand{\nImpulseR}{\impulseScalar_{nr}}
\newcommand{\zImpulseR}{\impulseScalar_{zr}}

\newcommand{\zImpulseL}{\impulseScalar_{zl}}

% The smallest value of normal impulse at which v_t = 0.
\newcommand{\nImpulseS}{\impulseScalar_{ns}}
\newcommand{\nImpulseL}{\impulseScalar_{nl}}
\newcommand{\zImpulseS}{\impulseScalar_{zs}}
\newcommand{\impulseS}{\impulse_{s}}

% The derivative of impulse w.r.t. normal impulse before I_ns
\newcommand{\iDniB}{\bs{\delta}}
% The derivative of impulse w.r.t. normal impulse after the collinear condition applied.
\newcommand{\iDniL}{\bs{\delta}}

% The derivative of impulse w.r.t. normal impulse after I_ns
\newcommand{\iDniA}{\bs{\sigma}}

% The potential energy when the collinear condition applies
\newcommand{\pEnergyL}{\pEnergy_{l}}

% The potential energy when compression phase ends
\newcommand{\pEnergyC}{\pEnergy_{c}}

% The potential energy when v_t = 0 first holds
\newcommand{\pEnergyT}{\pEnergy_{s}}


%%%%%%%%%%%%%%%%%%%%%%%%%% state space models

%%%%%%%%%%%%%%%%%%%%%%%%%%%%%%%%%%%%%%%%%%%%%%%%%%%%%%%%%%%%%%%%%%%%%%%%%%%%%%%%
%% 17. LOCOMOTION AND BALANCE
%%%%%%%%%%%%%%%%%%%%%%%%%%%%%%%%%%%%%%%%%%%%%%%%%%%%%%%%%%%%%%%%%%%%%%%%%%%%%%%%

\newcommand{\ankleTorque}{{^a\bs{\tau}}}

% The arguments are: origin(or a point), surface normal
\newcommand{\plane}[2]{\Pi(#1, #2)}

% It rotates from frame #1 to #2.

%%%%%%%%%%%%%%%%%%%%%%%%% Polynomials
%% #1-#3 Coeffecients #4 variable name
\newcommand{\cop}[1][default]{
        \IfStrEq{#1}{text}{\text{CoP}}{
        {\text{cop}}
        }
}


\newcommand{\cmp}{{\textcolor{textColor}{\text{\bf c}}_{\text{mp}}}} % The centroidal moment pivot


\newcommand{\height}[1]{{d_{#1}}}

% The divergent component of motion (DCM) or the instantaneous capture point (ICP)
\newcommand{\dcmDef}{{\com + \frac{\com[dot]}{\pendulumC}}}
\newcommand{\dcmxDef}{{\com_x + \frac{\com[dot]_x}{\pendulumC}}}

% The center of mass (COM)
\newcommand{\comText}{{\text{CoM}}\xspace}


% stance foot location
\newcommand{\stanceFL}{{\bs{p}_s}}
% capture foot location
\newcommand{\captureFL}{{\bs{p}_c}}

% The arguments are: (1) starting point (2) the ending point. For instance: \vec{a} - \vec{b} = \pointer{b}{a}.
\newcommand{\areaSymbol}{\mathcal{S}}

\newcommand{\applyAreaSubscript}[1]{%
    \IfStrEqCase{#1}{%
    {zmp}{
     \zmp
    }%
    {com}{
     \com
    }%
    {comd}{
    \com[dot]
    }%
    {dcm}{
     \dcm
    }%
    }[] % "default" string
}

\NewDocumentCommand{\area}{O{}}{%
{    \ensuremath{%
        {\areaSymbol}_{\applyAreaSubscript{#1}}
    }
    }
}

\newcommand{\ssRegion}{{\set{\text{ssr}}}}

\newcommand{\controlArea}{{\mathcal{S}_{\com\zmp}}}
\newcommand{\contactArea}{{\mathcal{S}_{\contactPoint}}}
\newcommand{\footTask}{p}

% #1 joint number
\newcommand{\stepIndicator}{{\bf \text{O}}}
\newcommand{\stepIndicatorSum}{\stepIndicator_{\mpcHorizon}}

% #1 foot index
\newcommand{\footLocation}[1]{\contactPoint^{#1}}

\newcommand{\fL}{p}

\newcommand{\timeStep}{N}
% footStep
\newcommand{\footStep}{M}


% footStep
\newcommand{\step}[2]{
  {#1^{#2}}
}

%%%%%%%%%%%%%%%%%%%%%%%%%% Ellipsoid:

%%%%%%%%%%%%%%%%%%%%%%%%%%%%%%%%%%%%%%%%%%%%%%%%%%%%%%%%%%%%%%%%%%%%%%%%%%%%%%%%
%% 18. CONTROL
%%%%%%%%%%%%%%%%%%%%%%%%%%%%%%%%%%%%%%%%%%%%%%%%%%%%%%%%%%%%%%%%%%%%%%%%%%%%%%%%

\newcommand{\costFunctional}{J}
\newcommand{\stateWeight}{Q}

\newcommand{\inputWeight}{R}

\newcommand{\discountF}{\gamma}

% Set
\newcommand{\est}[1]{{\hat{#1}}}

% The damping ratio
\newcommand{\dRatio}{\xi}

% The critical frequency

\newcommand{\gainConstant}{\text{k}}
\newcommand{\stiffness}{\text{k}_{\text{s}}}

\newcommand{\pGain}{\gainConstant_{\text{p}}}
\newcommand{\dGain}{\gainConstant_{\text{d}}}
\newcommand{\iGain}{\gainConstant_{\text{i}}}

% #1:index of the pole
\newcommand{\pole}[1]{p_{#1}}

% The system momentum matrix
\newcommand{\samplingPeriod}{\Delta t}
\newcommand{\sampletime}{\Delta t}
\newcommand{\Kinv}{\bs{K}^{-1}}
\newcommand{\JEstimate}{\bs{K}_{\jvelocitiesJump}}
\newcommand{\JpinvEstimate}{\Kinv_{\jvelocitiesJump}}
\newcommand{\LambdaINVEstimate}{\bs{K}_{\impulse}}
\newcommand{\LambdaEstimate}{\Kinv_{\impulse}}

\newcommand{\weight}{w}

\newcommand{\cFreq}{{\omega_0}}
% The damped natural frequency
\newcommand{\dcFreq}{{\omega_d}}

\newcommand{\costate}{\lambda}%
\newcommand{\hamiltonian}{H}%

\newcommand{\initialTime}{t_0}%
\newcommand{\finalTime}{t_f}%

% #1 denotes the input weight: R; #2 denotes the state
\newcommand{\lqrControl}[2]{
  - \inverse{(#1 + \quadraticObj{B}{P_{n+1}})}\transpose{B}P_{n+1}A#2
}%

\newcommand{\manifold}{\mathcal{S}}%

% %%%%%%%%%%%%%%%%%%%%%%%%% LCP  VARIABLES:
\newcommand{\ssDim}{{n_s}}

\newcommand{\ssA}{A}
\newcommand{\ssB}{B}
\newcommand{\ssC}{C}
\newcommand{\ssD}{D}

\newcommand{\ssAd}{{A_d}}
\newcommand{\ssBd}{{B_d}}
\newcommand{\ssCd}{{C_d}}
\newcommand{\ssDd}{{D_d}}


\newcommand{\stateSymbol}{\text{\bf x}}
% #1 derivatives, #2 ref/min/, #3 x/y/xy, #4 current/next/previous
\NewDocumentCommand{\state}{O{}O{}O{}O{}}{%
  {    \ensuremath{%
      % Check the first optional argument for dot or double dot
      %% \applyOverheadSymbols{#1}\text{\bf z}
      \applyOverheadSymbols{#1}{\stateSymbol}
      % Check the second optional argument for 'ref'
      ^{\applySuperscript{#2}}%
      _{\applySubscript{#3}}
      \applyPostscript{#4}
    }
  }
}

% #1 current/next/previous
\NewDocumentCommand{\randomNoise}{O{}}{%
  {    \ensuremath{%
      \epsilon^{\applyPostscript{#1}}
    }
  }
}


\newcommand{\ssInput}{\text{\bf u}}

\newcommand{\x}[1]{{#1}_{x}}
\newcommand{\y}[1]{{#1}_{y}}
\newcommand{\z}[1]{{#1}_{z}}
\newcommand{\xy}[1]{{#1}_{x,y}}

\newcommand{\fwAngle}{\theta} % Fly wheel angle


%%%%%%%%%%%%%%%%%%%%%%%%%% QP commands:
\newcommand{\eqA}[1]{{\text{A}_{#1}}}
\newcommand{\eqb}[1]{{\text{b}_{#1}}}
\newcommand{\ieqC}[1]{{\text{C}_{#1}}}
\newcommand{\ieqd}[1]{{\text{d}_{#1}}}
\newcommand{\cpFrame}{\text{c}}

% Task-space elastic wrench:
\newcommand{\eWrench}{\wrenchFrame{\ee}{\Delta}}


% Object-level impedance control:
\newcommand{\ipdMass}{\text{K}_{\text{M}}}
\newcommand{\ipdDamper}{\text{K}_{\text{D}}}


% Operational space inertia matrix
\newcommand{\osim}[2][default]{
  \IfStrEq{#1}{def}
  {
    % Definition
    \inverse{(\quadraticObj{\jacobian_{#2}}
      {\inverse{\jsim}})}
  }{
    % Default notation
    \Lambda_{{#2}}
  }
}


% The desired frame
\newcommand{\oscc}[2][default]{%
  \IfStrEq{#1}{def}{
    % Gives the definition
    \transpose[inv]{\jacobian} \coriolis \inverse{\jacobian} - \osim{#2} \jacobian[dot]\inverse{\jacobian}
  }{
    \IfStrEq{#1}{dot}{
      % Defines the derivative
    }
    {
      % The default behavior
      \Gamma_{#2}}
  }
}


\newcommand{\osg}[2][default]{
  \IfStrEq{#1}{def}
  {
    % Definition
    \transpose{\jacobian_{#2}}
    \gravity
  }{
    \eta_{{#2}}
  }
}


%%%%%%%%%%%%%%%%%%%%%%%%%% MPC commands:
\newcommand{\olGain}{{\text{k}}}
\newcommand{\clGain}{{\text{K}}}



\newcommand{\totalCost}{{\text{J}_{0}}}


% Argument of the Minimum as a function of the perturbations around the i-th (x, u) pair
\newcommand{\argMinF}{{\text{Q}}}

% The argument denotes the step number
\newcommand{\costToGo}[1]{{\text{J}_{#1}}}
\newcommand{\costToGoDef}[1]{{\text{J}_{#1}}(\state, \controlSequence_{#1})}

% Value: the cost-to-go given a minimizing control sequence (not necessarily the optimal one), i.e., an instantiation of the cost-to-go.
\newcommand{\ctgIns}[1]{{\text{V}}(\state, #1)}

\newcommand{\ctgInsSimple}[1]{{\text{V}}^{#1}}

% #2 specifies the state
\newcommand{\ctgInsTwo}[2]{{\text{V}}(#2, #1)}

\newcommand{\ctgInsDef}[1]{\min_{\controlSequence_{#1}}\costToGoDef{#1}}

% #2 specifies the state
\newcommand{\ctgInsDefRecursion}[2]{\min_{\controlInput{#1}}
  (\runningCostDef{#1} + \ctgInsTwo{#1 + 1}{#2})
}



\newcommand{\totalCostDef}{\totalCost(\state, \controlSequence)}

\newcommand{\runningCost}{l}
\newcommand{\runningCostDef}[1]{\runningCost(\controlInput{#1}, \state_{#1})}

\newcommand{\finalCost}{{\runningCost_{\text{f}}}}
\newcommand{\finalCostDef}{\runningCost_{\text{f}}(\state_{\timeHL})}

\newcommand{\controlSequence}{{\bf{\text{U}}}}

% time horizon length
\newcommand{\timeHL}{N}

% time horizon
\newcommand{\timeH}{\text{T}}

% state dimension
\newcommand{\controlSequenceDef}[1]{\{\controlInput{#1}, \controlInput{#1 + 1},\ldots, \controlInput{\timeHL-1}\}}
\newcommand{\controlSequenceDefTwo}{\{\controlInput{0}, \controlInput{1},\ldots, \controlInput{\timeHL-1}\}}

% System Dynamics
\newcommand{\system}{{\text{f}}}

% Perturbed state
\newcommand{\perturb}[1]{{\delta #1}}
\newcommand{\perturbedState}[1]{{#1 + \delta #1}}

\newcommand{\objF}{J}

\newcommand{\mpcHorizon}{N}

\newcommand{\mpcfy}[1]{\bar{\bf #1}}



\newcommand{\mpcInput}{\mpcfy{\bf \controlInput{}}}
\newcommand{\mpcState}{\mpcfy{\state}}
\newcommand{\mpcStateWeight}{\mpcfy{Q}}
\newcommand{\mpcInputWeight}{\mpcfy{R}}

\newcommand{\mpcA}{\mpcfy{A}}
\newcommand{\mpcB}{\mpcfy{B}}
\newcommand{\mpcC}{\mpcfy{C}}
\newcommand{\mpcD}{\mpcfy{D}}

% Foot Step Location:
\newcommand{\decisionVariable}{\boldsymbol{\nu}}
\newcommand{\decisionVar}{\boldsymbol{\nu}}

% #1 velocity #2 mass
\newcommand{\desired}[1]{#1^{\text{des}}}
\newcommand{\optimized}[1]{#1^{*}}
\newcommand{\optimal}[1]{#1^{*}}
\newcommand{\reff}[1]{#1^{\text{ref}}}
\newcommand{\measured}[1]{#1^{\circ}}
\newcommand{\disturb}[1]{\delta{#1}}


% integration error
\newcommand{\iError}[1]{\int{\bs{e}_{#1}}}
\newcommand{\iErrorDef}[1]{\int{\errorDef{#1}}dt}

% Defines the function to be minimized in a QP
\newcommand{\function}[1]{\text{f}_{#1}}

\newcommand{\error}[1]{\bs{e}_{#1}}
\newcommand{\errord}[1]{\dot{\bs{e}}_{#1}}
\newcommand{\errordd}[1]{\ddot{\bs{e}}_{#1}}

\newcommand{\errorDef}[1]{
  #1 - \reff{#1}
}

% The constant term of the error dynamics
\newcommand{\cTerm}[1]{
  \text{C}_{#1}
}

% The quadratic term of the error dynamics
\newcommand{\qTerm}[1]{
  \text{Q}_{#1}
}

% The penalty term of the error dynamics
\newcommand{\pTerm}[1]{
  \text{P}_{#1}
}


\newcommand{\errorDefFull}[1]{
 \selectionMatrix #1 - \selectionMatrix \reff{#1}
}

% #1:index of the input
\newcommand{\controlInput}[1]{{\text{\bf{u}}_{#1}}}
\newcommand{\control}[1]{{\text{\bf{u}}_{#1}}}

%% #1 real part, #2 imaginery part

%%%%%%%%%%%%%%%%%%%%%%%%%%%%%%%%%%%%%%%%%%%%%%%%%%%%%%%%%%%%%%%%%%%%%%%%%%%%%%%%
%% 19. STATE ESTIMATION AND OBSERVERS
%%%%%%%%%%%%%%%%%%%%%%%%%%%%%%%%%%%%%%%%%%%%%%%%%%%%%%%%%%%%%%%%%%%%%%%%%%%%%%%%

\newcommand{\oGain}{K_O} % The diagonal observer gain matrix.

\newcommand{\reminder}{\bs{r}_{\text{M}}}
\newcommand{\reminderd}{\dot{\bs{r}}_{\text{M}}}
\newcommand{\reminderdDef}{\oGain(\momentum - \est{\momentumd})}

\newcommand{\ext}[1]{{#1}_{\text{ext}}}

%%%%%%%%%%%%%%%%%%%%%%%%%% special terms
\newcommand{\bNoise}[1]{\text{b}_{#1}}

% Gaussian noise
\newcommand{\gNoise}[1]{w_{#1}}


% covariance matrix
\newcommand{\cov}[1]{\text{Q}_{#1}}

% Leg kinematics
% Distance from the CoM to a contact point i: p_i - com
\newcommand{\disDef}[1]{\pointer{\com}{\contactPoint_{#1}}}
\newcommand{\dis}[1]{\text{s}_{#1}}



% Forward kinematics
\newcommand{\fk}[1]{\text{lkin}_{#1}(\jangles)}

\newcommand{\imu}{\text{m}}



% subequations

%%%%%%%%%%%%%%%%%%%%%%%%%% Impedance control commands:

% The compliant frame

%%%%%%%%%%%%%%%%%%%%%%%%%%%%%%%%%%%%%%%%%%%%%%%%%%%%%%%%%%%%%%%%%%%%%%%%%%%%%%%%
%% 20. COLLISION AVOIDANCE
%%%%%%%%%%%%%%%%%%%%%%%%%%%%%%%%%%%%%%%%%%%%%%%%%%%%%%%%%%%%%%%%%%%%%%%%%%%%%%%%

\newcommand{\link}[1]{\text{l}_{#1}}


%  Define a macro with an optional argument for dotting
\newcommand{\cDist}[3][]{#1{\text{dis}}_{#2,#3}}
% The secure distance that two links should never reach
\newcommand{\sDist}[1]{\sigma_{\text{s}, #1}}

% The threshold distance below which, the collision avoidance constraint will be activated
\newcommand{\tDist}[1]{\sigma_{\text{t}, #1}}

\newcommand{\dConstant}[1]{\text{c}_{#1}}

%%%%%%%%%%%%%%%%%%%%%%%%%% State estimation commands:

% bias (drift) noise

%%%%%%%%%%%%%%%%%%%%%%%%%%%%%%%%%%%%%%%%%%%%%%%%%%%%%%%%%%%%%%%%%%%%%%%%%%%%%%%%
%% 21. MISCELLANEOUS
%%%%%%%%%%%%%%%%%%%%%%%%%%%%%%%%%%%%%%%%%%%%%%%%%%%%%%%%%%%%%%%%%%%%%%%%%%%%%%%%

\newcommand{\uncertain}[1]{\widetilde{#1}}

\newcommand{\eePos}{\bs{x}} % End-effector position
\newcommand{\eeVel}{\dot{\eePos}}
\newcommand{\eeAcc}{\ddot{\eePos}}
\newcommand{\optVars}{\bs{u}}

% The center of pressure (COP) point
\newcommand{\specialJac}{\mathcal{J}}
\newcommand{\specialJacTwo}{\mathcal{C}}

%%%%%%%%%%%%%  Different support areas
\newcommand{\stateVar}{\bs{x}}

% #1 variable with different steps, #2 step number
\newcommand{\semiAxis}{
a
}

\newcommand{\density}{
\rho
}



%%%%%%%%%%%%%%%%%%%%%%%%%% Eigenvalue:


\newcommand{\eigV}{
\sigma
}

\newcommand{\eigVec}{
\bs{e}
}



%%%%%%%%%%%%%%%%%%%%%%%%%% Rigid body motion


\newcommand{\guassD}{\mathcal{D}}

\newcommand{\quantity}{\bs{d}}
\newcommand{\quantityd}{\dot{\quantity}}
\newcommand{\quantitydd}{\ddot{\quantity}}

%%%%%%%%%%%% Axis:
\newcommand{\pre}[1]{
  \setSymbol^p(#1)
}


%%% all the links between [root and predecessors] while(link.hasPredecessor()){set+=predecessor(link); link = link.predecessor}):
\newcommand{\recursivePre}[1]{
  \setSymbol^{rp}(#1)
}

%%% successor:
\newcommand{\suc}[1]{
  \setSymbol^s(#1)
}

%%% all the links between [root and predecessors] while(link.hasSuccessor()){set+=successor(link); link = link.successor}):
\newcommand{\recursiveSuc}[1]{
  \setSymbol^{rs}(#1)
}


%%% CRB vel


\newtheorem{theorem}{\oBlock{Theorem}}
\newtheorem{Lemma}{\oBlock{Lemma}}
\newtheorem{remark}{\oBlock{Remark}}[section]
\newtheorem{assumption}{assumption}
\newtheorem{hypothesis}{hypothesis}
\newtheorem{proposition}[theorem]{proposition}
\newtheorem{property}[section]{\oBlock{property}}
\newtheorem{corollary}[theorem]{corollary}
\newtheorem{Example}[theorem]{Example}
\newtheorem{problem}{Problem}
\newtheorem{Solution}{Solution}
\newtheorem{Definition}{Definition}


\newtheorem{thm}{Theorem}[section]
\newtheorem{cor}[thm]{Corollary}
\newtheorem{lem}[thm]{Lemma}
\newtheorem{prop}[thm]{Proposition}

\newtheorem{defn}[thm]{Definition}

\newtheorem{rem}[thm]{Remark}

\usepackage{yuquanTitle}


%%%%%%%%%%%%%%%%%%%%%%%%%%% -------------------------------------------


% Report-local rolling-contact notation. Each definition builds on the shared
% notation and is introduced only because the shared package has no wheel-
% specific equivalent.
\newcommand{\generalizedConfiguration}{\jangles[][][generalized]}
% Keep the velocity symbol script-free so shared sample operators such as
% \nextState can append their own time index without a double subscript.
\newcommand{\generalizedVelocity}{%
  \applyOverheadSymbols{dot}{\applyJangles{generalized}}}
\newcommand{\generalizedAcceleration}{\jangles[ddot][][generalized]}
\newcommand{\rollingDirection}{\bs{t}}
\newcommand{\lateralDirection}{\bs{s}}
\newcommand{\wheelAngle}{\phi}
\newcommand{\wheelRadius}{r}
\newcommand{\wheelSelector}{\bs{S}^{\wheelAngle}}
\newcommand{\tangentProjector}{\text{P}}
\newcommand{\slipVelocity}[1]{\contactVel_{\text{slip},#1}}
\newcommand{\rollingMatrix}{\text{A}}
\newcommand{\rollingTarget}{\bs{b}}
\newcommand{\rollingResidual}{\bs{e}^{\text{roll}}}
\newcommand{\slackVariable}{\bs{\sigma}}

\project{\tinytf Rolling contacts in whole-body QP control}
\author{
  Yuquan WANG \\
  Tencent Robotics X \\
  Shenzhen, P.R. China\\
  \vspace{5pt}
  \texttt{hasbragwang@tencent.com}\vspace{30pt} \\
}
\summary{
  This tutorial extends an acceleration-level whole-body quadratic program
  with rolling, sticking, and sliding wheel contacts. It derives QP-ready
  kinematic constraints, connects them to friction and actuation limits, and
  provides a mode-aware implementation and validation procedure for floating-
  base wheeled and leg-wheeled robots.
}

\rhead{Robotics-X \colorbox{teal}{\textcolor{white}{\quad \bf Rolling-contact QP control}}}
%%%%%%%%%%%%%%%%%%%%%%%%%%% -------------------------------------------


\begin{document}

\maketitle

\section{Introduction}
\label{sec:rolling-introduction}

Optimization-based whole-body control computes joint accelerations,
actuation torques, and contact forces while respecting joint, contact,
collision, and balance constraints. The baseline optimization-based control
report treats a sustained contact as a point
whose Cartesian acceleration is regulated to zero. This model is suitable for
a fixed foot or hand, but it prevents the intentional tangential motion of a
wheel. Conversely, removing the contact constraint entirely loses
no-penetration and lateral no-slip conditions.

The rolling-contact notes~\cite{rollingContactNotion,completeRollingContactsNotion}
identify the required components: a local wheel-contact frame, longitudinal
rolling and lateral constraints, slip detection, coupled multi-wheel
kinematics, friction and torque feasibility, floating-base integration, and
tests on ramps, non-coplanar terrain, and mixed leg-wheel support. We turn
those components into a consistent acceleration-level QP extension. In
particular, we use one generalized velocity throughout, restore the missing
equal signs in the slip definitions, and use the dimensionally correct torque
relation $\norm[abs]{\jtorques_{\wheelAngle_i}} \geq
\wheelRadius_i\norm[abs]{\innerP{\rollingDirection_i}{\force_i}}$ only as a
quasi-static interpretation of the complete equations of motion.

We define the rolling-contact geometry and slip residual in
\secRef{sec:rolling-model}. We differentiate the ideal velocity relation to
obtain acceleration-level constraints in \secRef{sec:rolling-constraints}, and
connect them to friction, actuation, and balance in
\secRef{sec:rolling-dynamics}. We then formulate rolling objectives, contact
mode transitions, and conflicting contacts in \secRef{sec:rolling-tasks}.
A simplified ideal-rolling QP and its stepwise extension to the complete
mode-aware controller appear in \secRef{sec:rolling-qp}, followed by
implementation guidance in \secRef{sec:rolling-implementation} and validation
cases in \secRef{sec:rolling-cases}.

\section{Rolling-contact model}
\label{sec:rolling-model}

We first choose a convention that separates the velocity of the wheel carrier
from the surface velocity generated by wheel spin. This distinction removes
an ambiguity in the source notes and determines whether the rolling term
$\wheelRadius_i\dot{\wheelAngle}_i$ belongs inside or outside the contact
Jacobian.

\subsection{Generalized state and local contact frame}

Let $\generalizedConfiguration$ denote the floating-base generalized
configuration, and let $\generalizedVelocity$ and
$\generalizedAcceleration$ denote its velocity and acceleration. The vector
$\generalizedVelocity$ contains the floating-base twist and every actuated
joint velocity, including the wheel rate $\dot{\wheelAngle}_i$.

For the \ith wheel, we define a unit surface normal $\surfaceNormal_i$, a unit
rolling direction $\rollingDirection_i$, and a unit lateral direction
$\lateralDirection_i$:
\quickEq{eq:rolling-local-frame}{
  \lateralDirection_i = \cross{\surfaceNormal_i}{\rollingDirection_i},
  \qquad
  \rotation{\inertialFrame}{c_i}
  = \begin{bmatrix}
      \rollingDirection_i & \lateralDirection_i & \surfaceNormal_i
    \end{bmatrix}.
}
The three directions must be orthonormal. They are recomputed from the current
wheel orientation and local terrain geometry at every control cycle, rather
than copied from a common horizontal plane.

Attach a carrier frame $\cframe[c_i]$ to the \ith wheel. Its origin
$\point[][][\inertialFrame,c_i]$ is the wheel rolling center: the center of the
undeformed circular wheel, located at the hub or axle intersection. It is not
the instantaneous terrain contact point. For a rigid wheel of radius
$\wheelRadius_i$ on the locally planar surface, the nominal contact point is
\quickEq{eq:wheel-center-contact-point}{
  \point[][][\inertialFrame,p_i]
  = \point[][][\inertialFrame,c_i]
  - \wheelRadius_i\surfaceNormal_i.
}
The wheel plane is spanned by $\rollingDirection_i$ and $\surfaceNormal_i$,
while $\lateralDirection_i$ is parallel to the axle. For a steering wheel,
the rolling and lateral directions rotate with its steering angle, but the
origin $\point[][][\inertialFrame,c_i]$ remains the carrier center.

Let $\jacobian[][c_i](\generalizedConfiguration)$ be the translational
Jacobian of this explicitly defined rolling center. Its velocity is
\quickEq{eq:rolling-carrier-velocity}{
  \contactVel_{c_i}
  = \jacobian[][c_i](\generalizedConfiguration)\generalizedVelocity.
}
The rolling, lateral, and normal components are therefore
\begin{align}
  \contactVel_{c_i}^{\text{roll}}
  &= \innerP{\rollingDirection_i}{\contactVel_{c_i}},
  \label{eq:rolling-component}\\
  \contactVel_{c_i}^{\text{lat}}
  &= \innerP{\lateralDirection_i}{\contactVel_{c_i}},
  \label{eq:lateral-component}\\
  \contactVel_{c_i}^{\text{nor}}
  &= \innerP{\surfaceNormal_i}{\contactVel_{c_i}}.
  \label{eq:normal-component}
\end{align}

\shadowRemark{
  \label{re:rolling-jacobian-convention}
  Equation~\eqref{eq:rolling-carrier-velocity} uses a carrier-point Jacobian
  whose longitudinal velocity excludes the rim velocity generated by
  $\dot{\wheelAngle}_i$. If an implementation instead constructs the Jacobian
  of the instantaneous wheel material point and includes wheel spin, pure
  rolling is simply zero material-point velocity. The two conventions are
  equivalent, but their terms must never be mixed in one constraint.
}

\subsection{Ideal rolling, tangential velocity, and slip residual}

For pure rolling, the carrier velocity along $\rollingDirection_i$ equals the
rim speed, while lateral and normal velocities vanish:
\begin{equation}
  \label{eq:ideal-rolling-velocity}
  \begin{aligned}
    \innerP{\rollingDirection_i}{\jacobian[][c_i]\generalizedVelocity}
      &= \wheelRadius_i\dot{\wheelAngle}_i,\\
    \innerP{\lateralDirection_i}{\jacobian[][c_i]\generalizedVelocity}
      &= 0,\\
    \innerP{\surfaceNormal_i}{\jacobian[][c_i]\generalizedVelocity}
      &= 0.
  \end{aligned}
\end{equation}
The third row is the no-penetration constraint. The second row prohibits
lateral skidding, and the first row prohibits longitudinal slip.

Define the tangent-plane projector as
\quickEq{eq:tangent-projector}{
  \tangentProjector_i
  = \bs{I}_3
  - \surfaceNormal_i\transpose{\surfaceNormal_i}
  = \rollingDirection_i\transpose{\rollingDirection_i}
  + \lateralDirection_i\transpose{\lateralDirection_i}.
}
Here, $\bs{I}_3\in\RRm{3}{3}$ is the identity matrix
on the ambient Cartesian space, so
$\tangentProjector_i\in\RRm{3}{3}$ projects onto the two-dimensional
tangent plane at contact $i$.
The carrier point's tangential velocity can be computed directly as
\quickEq{eq:tangential-carrier-velocity}{
  \contactVel_{c_i}^{\text{tan}}
  = \tangentProjector_i\jacobian[][c_i]\generalizedVelocity.
}
This quantity contains both the longitudinal carrier motion and any lateral
motion, but it does not by itself measure wheel--terrain slip. Under the
carrier-point Jacobian convention of Remark~\ref{re:rolling-jacobian-convention},
the relative tangential velocity at the wheel--terrain interface is obtained
by subtracting the rim velocity:
\quickEq{eq:slip-velocity}{
  \slipVelocity{i}
  = \contactVel_{c_i}^{\text{tan}}
  - \wheelRadius_i\dot{\wheelAngle}_i\rollingDirection_i
  = \left(
      \innerP{\rollingDirection_i}
        {\jacobian[][c_i]\generalizedVelocity}
      - \wheelRadius_i\dot{\wheelAngle}_i
    \right)\rollingDirection_i
  + \innerP{\lateralDirection_i}
      {\jacobian[][c_i]\generalizedVelocity}\lateralDirection_i.
}
Unlike a heuristic direction based only on the commanded wheel rate,
\eqref{eq:slip-velocity} measures the residual between the actual carrier
motion and the wheel surface speed. It therefore detects longitudinal slip,
lateral slip, and the residual motion produced by incompatible contacts.
Equations~\eqref{eq:tangential-carrier-velocity} and
\eqref{eq:slip-velocity} are definitions used for estimation, diagnostics,
and mode selection; they are not additional QP constraints. The ideal
velocity constraint is Equation~\eqref{eq:ideal-rolling-velocity}, whose
acceleration-level form is derived in Section~\ref{sec:rolling-constraints}.

We use hysteresis rather than a single threshold:
\begin{equation}
  \label{eq:slip-hysteresis}
  \begin{aligned}
    \norm{\measured{\slipVelocity{i}}}
      &> \varepsilon_{\text{enter}}
      &&\Longrightarrow \text{enter sliding mode},\\
    \norm{\measured{\slipVelocity{i}}}
      &< \varepsilon_{\text{exit}}
      &&\Longrightarrow \text{leave sliding mode},\\
    \varepsilon_{\text{exit}}
      &< \varepsilon_{\text{enter}}.
  \end{aligned}
\end{equation}
The measured residual should be filtered and checked for persistence over
several cycles. A learned detector may augment this residual, but it should not
replace the deterministic kinematic and force-feasibility checks.

\subsection{Stacked multi-contact relation}

For $\numContacts$ rolling candidates, define
\begin{equation}
  \label{eq:stacked-rolling-velocity}
  \rollingMatrix(\generalizedConfiguration)\generalizedVelocity
  = \rollingTarget(\dot{\wheelAngle}),
\end{equation}
where the \ith three-row block is
\quickEq{eq:rolling-block}{
  \rollingMatrix_i
  = \begin{bmatrix}
      \transpose{\rollingDirection_i}\jacobian[][c_i]\\
      \transpose{\lateralDirection_i}\jacobian[][c_i]\\
      \transpose{\surfaceNormal_i}\jacobian[][c_i]
    \end{bmatrix},
  \qquad
  \rollingTarget_i
  = \vectorThree{
      \wheelRadius_i\dot{\wheelAngle}_i
    }{0}{0}.
}
Equation~\eqref{eq:stacked-rolling-velocity} exposes coupling between wheels
through their common floating-base coordinates. It must be assembled with each
wheel's own local frame, especially for ramps and non-coplanar contacts.

If wheel spin is part of $\generalizedVelocity$, let $\wheelSelector_i$ select
$\dot{\wheelAngle}_i$ from it. The homogeneous version of the \ith block is
\quickEq{eq:homogeneous-rolling-block}{
  \underbrace{
  \begin{bmatrix}
    \transpose{\rollingDirection_i}\jacobian[][c_i]
      - \wheelRadius_i\wheelSelector_i\\
    \transpose{\lateralDirection_i}\jacobian[][c_i]\\
    \transpose{\surfaceNormal_i}\jacobian[][c_i]
  \end{bmatrix}}_{\text{G}_i(\generalizedConfiguration)}
  \generalizedVelocity = \zeroVector.
}
This form avoids treating the wheel rate both as an independent command and as
an element of the QP state.

\section{Rolling-contact constraints}
\label{sec:rolling-constraints}

The QP in the baseline report optimizes acceleration-level quantities. The
velocity constraint~\eqref{eq:stacked-rolling-velocity} must therefore be
differentiated once, just as a fixed Cartesian contact is converted into
$\jacobian\jaccelerations+\jacobian[dot]\jvelocities=0$.

\subsection{Acceleration-level equality}

Define the rolling residual
\quickEq{eq:rolling-residual}{
  \rollingResidual
  = \rollingMatrix\generalizedVelocity - \rollingTarget.
}
Differentiating it and adding proportional stabilization yields
\quickEq{eq:rolling-acceleration-constraint}{
  \rollingMatrix\generalizedAcceleration
  = \dot{\rollingTarget}
  - \dot{\rollingMatrix}\generalizedVelocity
  - \pGain\rollingResidual.
}
At a control cycle, $\rollingMatrix$, $\dot{\rollingMatrix}$,
$\generalizedVelocity$, $\rollingTarget$, and $\dot{\rollingTarget}$ are known.
Consequently, \eqref{eq:rolling-acceleration-constraint} is affine in the QP
decision variables. The derivative $\dot{\rollingMatrix}$ includes the time
variation of the Jacobian, steering direction, lateral direction, and terrain
normal. Dropping the direction derivatives is valid only for a fixed local
frame.

For the homogeneous relation~\eqref{eq:homogeneous-rolling-block}, the same
constraint becomes
\quickEq{eq:homogeneous-rolling-acceleration}{
  \text{G}_i\generalizedAcceleration
  = -\dot{\text{G}}_i\generalizedVelocity
  - \pGain\text{G}_i\generalizedVelocity.
}
This is the preferred implementation when wheel accelerations are elements of
$\generalizedAcceleration$.

\subsection{Hard and soft rows}

The three rows of \eqref{eq:rolling-acceleration-constraint} do not always
deserve the same priority:

\begin{itemize}
  \item \hlBlock{Normal row}: keep it hard while contact is active, because it
    enforces no penetration and contact maintenance.
  \item \hlBlock{Lateral row}: keep it hard for an ideal conventional wheel;
    soften or remove it for omni-wheels, casters with a different model, or
    detected lateral sliding.
  \item \hlBlock{Rolling row}: keep it hard only when pure rolling is feasible
    and required. Otherwise introduce a slack or a quadratic rolling task so
    fixed contacts and safety constraints remain feasible.
\end{itemize}

Let $\text{a}_{t,i}$ and $\text{d}_{t,i}$ denote the longitudinal row and
right-hand side of \eqref{eq:rolling-acceleration-constraint}. A soft rolling
equality writes
\quickEq{eq:soft-rolling-equality}{
  \text{a}_{t,i}\generalizedAcceleration
  - \text{d}_{t,i} = \slackVariable_i,
}
and the QP minimizes $\weight_{\sigma,i}\norm[two]{\slackVariable_i}$.
The slack is not merely a numerical device: after integration, it contributes
to the predicted rolling mismatch and provides a useful diagnostic for slip or
an infeasible wheel command.

\shadowRemark{
  A hard rolling equality can make the QP infeasible before a physical wheel
  begins to slide. Use hard constraints for geometry and safety, and represent
  a commanded wheel rate as a high-priority task unless feasibility has been
  established. The resulting residual is preferable to silently violating a
  fixed-contact or joint-limit constraint.
}

\section{Contact force, actuation, and balance}
\label{sec:rolling-dynamics}

Rolling kinematics alone cannot guarantee pure rolling. The optimized contact
force must lie inside the friction cone, and the actuator must supply the
torque required by the coupled robot-wheel dynamics. These checks reuse the
force generators and torque bounds already present in the baseline QP.

\subsection{Friction cone in the rolling frame}

Resolve the \ith contact force along the local directions:
\begin{align}
  \force_{t,i} &= \innerP{\rollingDirection_i}{\force_i},
  \label{eq:rolling-force-component}\\
  \force_{s,i} &= \innerP{\lateralDirection_i}{\force_i},
  \label{eq:lateral-force-component}\\
  \force_{n,i} &= \innerP{\surfaceNormal_i}{\force_i}.
  \label{eq:normal-force-component}
\end{align}
The Coulomb condition is
\quickEq{eq:rolling-friction-cone}{
  \norm{\vectorTwo{\force_{t,i}}{\force_{s,i}}}
  \leq \coef{f}_i\force_{n,i},
  \qquad
  \force_{n,i} \geq 0.
}
As in the baseline formulation, we approximate this cone with generators:
\quickEq{eq:rolling-force-generators}{
  \force_i = \cone[f][i]\generator{\force_i},
  \qquad
  \generator{\force_i} \geq \zeroVector.
}
The basis $\cone[f][i]$ must be constructed in the local rolling frame
$[\rollingDirection_i,\lateralDirection_i,\surfaceNormal_i]$ and updated when
the wheel steers or the terrain normal changes. A planar wheel contact with a
finite patch may additionally use the contact-wrench cone from the baseline
controller.

When sliding is established, the friction force opposes the measured slip
velocity, not necessarily the commanded rolling direction. Holding the
measured direction fixed for one control cycle gives the affine approximation
\quickEq{eq:sliding-friction-direction}{
  \tangentProjector_i\force_i
  = -\coef{f}_i\force_{n,i}
  \frac{\measured{\slipVelocity{i}}}
       {\norm{\measured{\slipVelocity{i}}}+\varepsilon}.
}
Equation~\eqref{eq:sliding-friction-direction} is used only for sustained
sliding. During sticking or pure rolling, tangential force remains an interior
decision of the friction cone and must not be forced to its boundary.

\subsection{Equations of motion and wheel torque}

With contact-force generators, the floating-base equations of motion write
\quickEq{eq:rolling-equations-of-motion}{
  \jsim\generalizedAcceleration + \gravityandcoriolis
  = \actuationMatrix\jtorques
  + \agg{i}{\numContacts}{
      \transpose{\jacobian[][p_i]}\cone[f][i]\generator{\force_i}
    }.
}
The dimensions and unactuated floating-base rows are handled by the actuation
matrix $\actuationMatrix$. Equation~\eqref{eq:rolling-equations-of-motion}
couples wheel acceleration, wheel torque, floating-base acceleration, and
contact forces; it must be imposed only once.

If $\wheelSelector_i$ selects the \ith wheel actuator row, its torque is
\quickEq{eq:wheel-torque-row}{
  \jtorques_{\wheelAngle_i}
  = \wheelSelector_i\left(
      \jsim\generalizedAcceleration + \gravityandcoriolis
      - \agg{j}{\numContacts}{
          \transpose{\jacobian[][p_j]}\force_j
        }
    \right).
}
The ordinary actuation bounds then apply:
\quickEq{eq:wheel-torque-bound}{
  \wheelSelector_i\lowerBound{\jtorques}
  \leq \jtorques_{\wheelAngle_i}
  \leq \wheelSelector_i\upperBound{\jtorques}.
}
For an isolated wheel with negligible rotational acceleration and losses,
\eqref{eq:wheel-torque-row} reduces to the intuitive necessary condition
\quickEq{eq:quasistatic-wheel-force}{
  \norm[abs]{\innerP{\rollingDirection_i}{\force_i}}
  \leq
  \frac{\norm[abs]{\jtorques_{\wheelAngle_i}}}{\wheelRadius_i}.
}
Do not impose \eqref{eq:quasistatic-wheel-force} in addition to a complete
wheel row of \eqref{eq:rolling-equations-of-motion}; doing so may duplicate or
overconstrain the same dynamics.

\subsection{Floating-base balance}

Rolling contacts contribute to the net centroidal wrench exactly like other
sustained contacts. Therefore, the CoM, ZMP, DCM, and angular-momentum
constraints of the baseline QP remain valid after the fixed-contact set is
replaced by a mode-aware union of sticking, rolling, and sliding contacts. The
resultant contact wrench must still satisfy the centroidal dynamics, while
each local force satisfies \eqref{eq:rolling-force-generators}.

\shadowRemark{
  Kinematic rolling feasibility and dynamic force feasibility are different
  tests. A wheel-rate command may satisfy
  \eqref{eq:ideal-rolling-velocity} while requiring a force outside
  \eqref{eq:rolling-friction-cone} or a torque outside
  \eqref{eq:wheel-torque-bound}. In that case the controller must predict or
  detect sliding and change contact mode; it must not retain the pure-rolling
  equality and merely clip the torque.
}

\section{Rolling-contact task formulation and contact modes}
\label{sec:rolling-tasks}

We now combine the geometric residual, force feasibility, and contact state
into QP tasks. Let $\set{f}$, $\set{r}$, and $\set{s}$ denote fixed/sticking,
rolling, and sliding contact sets, respectively.

\subsection{Geometric carrier--base coupling}
\label{sec:carrier-base-coupling}

For a differential-drive or four-steering-wheel chassis on a plane, the QP
must explicitly couple every carrier-center velocity to the same
floating-base twist. Let $\Pi$ be the common terrain plane, either horizontal
or a ramp, with unit normal $\surfaceNormal_{\Pi}$. Choose orthonormal tangent
axes ${}^{\inertialFrame}\bs{e}_{\Pi,x}$ and
${}^{\inertialFrame}\bs{e}_{\Pi,y}$ in the inertial frame, and define
${}^{\inertialFrame}\text{E}_{\Pi}
=\begin{bmatrix}
  {}^{\inertialFrame}\bs{e}_{\Pi,x}&
  {}^{\inertialFrame}\bs{e}_{\Pi,y}
\end{bmatrix}$.
The origin $\point[][][\inertialFrame,\fbFrame]$ is the center of the
floating-base task frame projected onto $\Pi$. The planar coordinates of the
\ith rolling center relative to this origin are
\quickEq{eq:wheel-center-planar-offset}{
  \bs{\rho}_i
  = \vectorTwo{x_i}{y_i}
  = \transpose{{}^{\inertialFrame}\text{E}_{\Pi}}
    \left(
      \point[][][\inertialFrame,c_i]
      - \point[][][\inertialFrame,\fbFrame]
    \right).
}
Thus, $\point[][][\inertialFrame,\fbFrame]$ and
$\point[][][\inertialFrame,c_i]$ are distinct, explicitly defined origins.

Write the floating-base velocity as the body twist
\quickEq{eq:base-body-velocity}{
  \bodyVel{\inertialFrame}{\fbFrame}
  = \vectorTwo{
      \bodyTV{\inertialFrame}{\fbFrame}
    }{
      \bodyRV{\inertialFrame}{\fbFrame}
    }.
}
For the body-twist equations, rotate the plane geometry into $\cframe[fb]$:
\quickEq{eq:terrain-basis-body-frame}{
  {}^{\fbFrame}\text{E}_{\Pi}
  = \rotationInv{\inertialFrame}{\fbFrame}
    {}^{\inertialFrame}\text{E}_{\Pi},
  \qquad
  {}^{\fbFrame}\surfaceNormal_{\Pi}
  = \rotationInv{\inertialFrame}{\fbFrame}
    {}^{\inertialFrame}\surfaceNormal_{\Pi}.
}
The plane-compatible part of the body twist is
\begin{equation}
  \label{eq:base-plane-projection}
  \prescript{\Pi}{}{\bodyVel{\inertialFrame}{\fbFrame}}
  =
  \underbrace{
    \begin{bmatrix}
      {}^{\fbFrame}\text{E}_{\Pi}
        \transpose{{}^{\fbFrame}\text{E}_{\Pi}} & \zeroMatrix\\
      \zeroMatrix &
        {}^{\fbFrame}\surfaceNormal_{\Pi}
        \transpose{{}^{\fbFrame}\surfaceNormal_{\Pi}}
    \end{bmatrix}
  }_{\text{P}^{\fbFrame}_{\Pi}}
  \bodyVel{\inertialFrame}{\fbFrame}.
\end{equation}
It retains translation tangent to the plane and rotation about its normal.
The corresponding three-dimensional planar twist is
\quickEq{eq:base-planar-twist}{
  \bs{\xi}^{\Pi}_{\fbFrame}
  = \vectorThree{
      \innerP{{}^{\fbFrame}\bs{e}_{\Pi,x}}
        {\bodyTV{\inertialFrame}{\fbFrame}}
    }{
      \innerP{{}^{\fbFrame}\bs{e}_{\Pi,y}}
        {\bodyTV{\inertialFrame}{\fbFrame}}
    }{
      \innerP{{}^{\fbFrame}\surfaceNormal_{\Pi}}
        {\bodyRV{\inertialFrame}{\fbFrame}}
    }.
}

If I understand correclty, \eqref{eq:base-plane-projection} might be wrong. If it projects the $\bodyVel{\inertialFrame}{\fbFrame}$ into the terrain plane $\Pi$, we can directly use the twist transform: 
\quickEq{ni}{
  \bodyVel{\Pi}{\fbFrame}   
   = \bodyVelTransform{\inertialFrame}{\fbFrame}{\Pi}
}. As the $bodyVel{\inertialFrame}{\fbFrame}$ is zero, so actually $\bodyVel{\Pi}{\fbFrame}$ is equal to $\bodyVel{\inertialFrame}{\fbFrame}$




The shared body-velocity transform gives the velocity of the carrier frame:
\quickEq{eq:wheel-vel-calc}{
  \bodyVel{\inertialFrame}{c_i}
  = \bodyVelTransform{\fbFrame}{c_i}{\inertialFrame}.
}
The right-hand side of \eqref{eq:wheel-vel-calc} expands to
$\twistTransform{\fbFrame}{c_i}\bodyVel{\inertialFrame}{\fbFrame}
+\bodyVel{\fbFrame}{c_i}$. Here $\cframe[c_i]$ is the chassis-aligned carrier
frame located at the wheel center; the steering angle is encoded by
$\rollingDirection_i$ rather than by rotating this carrier frame. For a rigid
carrier mount,
$\bodyVel{\fbFrame}{c_i}=\zeroVector$, so the projected relation becomes
\quickEq{eq:wheel-vel-def}{
  \prescript{\Pi}{}{\bodyVel{\inertialFrame}{c_i}}
  = \twistTransform{\fbFrame}{c_i}
    \text{P}^{\fbFrame}_{\Pi}
    \bodyVel{\inertialFrame}{\fbFrame}.
}
The translational part of this twist equation has the general planar form
\quickEq{eq:planar-carrier-map}{
  \contactVel_{c_i}^{\Pi}
  = \underbrace{
      \begin{bmatrix}
        1 & 0 & -y_i\\
        0 & 1 &  x_i
      \end{bmatrix}
    }_{\text{H}_i}
    \bs{\xi}^{\Pi}_{\fbFrame}.
}
This is the common geometric map for differential and steering wheels.

Let $\rollingDirection_i^{\Pi}$ and $\lateralDirection_i^{\Pi}$ be the
two-dimensional coordinates of the rolling and lateral directions in the
basis $[\bs{e}_{\Pi,x},\bs{e}_{\Pi,y}]$. Each wheel must satisfy
\quickEq{eq:geometric-wheel-velocity-task}{
  \begin{bmatrix}
    \transpose{\rollingDirection_i^{\Pi}}\\
    \transpose{\lateralDirection_i^{\Pi}}
  \end{bmatrix}
  \text{H}_i\bs{\xi}^{\Pi}_{\fbFrame}
  = \vectorTwo{\wheelRadius_i\dot{\wheelAngle}_i}{0}.
}
Equation~\eqref{eq:geometric-wheel-velocity-task} is the explicit geometric
task that makes every wheel rate agree with one floating-base velocity. If
$\text{S}_{\Pi}\text{S}_{\fbFrame}\generalizedVelocity
=\bs{\xi}^{\Pi}_{\fbFrame}$ and $\wheelSelector_i$ extracts the wheel rate,
its homogeneous QP row is
\quickEq{eq:geometric-wheel-homogeneous-task}{
  \underbrace{\left(
    \begin{bmatrix}
      \transpose{\rollingDirection_i^{\Pi}}\\
      \transpose{\lateralDirection_i^{\Pi}}
    \end{bmatrix}
    \text{H}_i\text{S}_{\Pi}\text{S}_{\fbFrame}
    - \vectorTwo{\wheelRadius_i}{0}\wheelSelector_i
  \right)}_{\text{G}^{\text{geo}}_i}
  \generalizedVelocity = \zeroVector.
}
The acceleration-level equality inserted into the QP is consequently
\quickEq{eq:geometric-wheel-acceleration-task}{
  \text{G}^{\text{geo}}_i\generalizedAcceleration
  = -\dot{\text{G}}^{\text{geo}}_i\generalizedVelocity
    - \pGain\text{G}^{\text{geo}}_i\generalizedVelocity.
}
For these two chassis types, use the rows of
\eqref{eq:geometric-wheel-acceleration-task} in place of the equivalent
longitudinal and lateral rows of
\eqref{eq:homogeneous-rolling-acceleration}; do not impose both copies. The
normal row remains the contact-maintenance constraint.

\subsection{Rolling task}

The longitudinal acceleration residual of wheel $i$ is
\quickEq{eq:rolling-task}{
  \function{\text{roll},i}(\generalizedAcceleration)
  = \text{a}_{t,i}\generalizedAcceleration - \text{d}_{t,i}.
}
Minimizing $\weight_{\text{roll},i}
\norm[two]{\function{\text{roll},i}}$ makes the wheel track its commanded
rolling rate without sacrificing higher-priority feasibility. The normal and
lateral rows remain hard for a wheel in $\set{r}$, while all three zero-motion
rows remain hard for a fixed contact in $\set{f}$.

For a sliding wheel in $\set{s}$, only the normal contact row remains hard.
The tangential kinematic equalities are removed, and the contact force obeys
the cone plus the one-cycle friction-direction approximation in
\eqref{eq:sliding-friction-direction}. This separation prevents an impossible
combination of nonzero slip velocity and a hard no-slip constraint.

\subsection{Conflicting fixed and rolling contacts}

Consider the question from the notes: three coplanar contacts are fixed while
the fourth wheel is commanded to roll. Let
$\text{A}_{\text{fix}}\generalizedAcceleration=\bs{d}_{\text{fix}}$ collect
the differentiated fixed-contact constraints. The best feasible rolling
command is obtained from
\begin{equation}
  \label{eq:fixed-rolling-conflict}
  \begin{aligned}
    \min_{\generalizedAcceleration,\slackVariable}
      \quad & \weight_{\sigma}\norm{\slackVariable}_2^2\\
    \mbox{s.t.}\quad
      & \text{A}_{\text{fix}}\generalizedAcceleration
        = \bs{d}_{\text{fix}},\\
      & \text{a}_{t,\text{roll}}\generalizedAcceleration
        - \text{d}_{t,\text{roll}} = \slackVariable.
  \end{aligned}
\end{equation}
If the fixed constraints eliminate the commanded motion, the optimizer
preserves them and exposes the incompatible rolling command through
$\slackVariable$. After integration, the sliding direction is the tangential
residual in \eqref{eq:slip-velocity}; it is not an arbitrary direction chosen
from the wheel command alone.

In practice, the controller should either release one or more fixed contacts,
accept wheel slip, or reject the command. Which response is correct depends on
the contact schedule and safety requirements, not only on a numerical task
weight.

\subsection{Mode transitions}

Abruptly adding or removing contact rows can produce discontinuous QP
solutions. Reuse the baseline task-activation strategy by assigning each
contact a continuous activation $\alpha_i(t)\in\interval{0}{1}$. Keep a fixed
maximum number of rows and vary the relevant task weight, slack penalty, or
constraint relaxation continuously during a transition.

The mode logic should combine:
\begin{itemize}
  \item \hlBlock{Kinematic evidence}: the filtered residual
    $\norm{\measured{\slipVelocity{i}}}$ and the slack
    $\norm{\slackVariable_i}$;
  \item \hlBlock{Force evidence}: friction-cone margin and normal force;
  \item \hlBlock{Actuation evidence}: distance to the wheel-torque bound;
  \item \hlBlock{Contact schedule}: whether the wheel is expected to stick,
    roll, detach, or slide;
  \item \hlBlock{Hysteresis}: separate entry and exit thresholds with a
    minimum dwell time.
\end{itemize}

\section{Rolling-contact QP: from ideal rolling to contact modes}
\label{sec:rolling-qp}

\subsection{Simplified ideal-rolling QP}

We first consider the minimal extension of the baseline whole-body QP. Assume
that every active wheel rolls ideally, that no wheel is sliding
($\set{s}=\emptyset$), and that the prescribed wheel motion is compatible
with the other hard constraints. No rolling slack or contact-mode variable is
needed. The decision vector is
\quickEq{eq:ideal-rolling-decision-vector}{
  \decisionVar_{\text{ideal}}
  = \begin{bmatrix}
      \generalizedAcceleration\\
      \jtorques\\
      \generator{\force_1}\\
      \vdots\\
      \generator{\force_{\numContacts}}
    \end{bmatrix}.
}
Over one sampling period, use the affine velocity prediction
\quickEq{eq:ideal-rolling-velocity-prediction}{
  \nextState{\generalizedVelocity}
  = \generalizedVelocity
  + \samplingPeriod\generalizedAcceleration.
}
Because $\wheelSelector_i\nextState{\generalizedVelocity}$ is the predicted
wheel rate, the ideal tangential-velocity residual is
\quickEq{eq:ideal-tangential-velocity-residual}{
  \function{\text{tan},i}(\decisionVar_{\text{ideal}})
  = \left(
      \tangentProjector_i\jacobian[][c_i]
      - \wheelRadius_i\rollingDirection_i\wheelSelector_i
    \right)
    \nextState{\generalizedVelocity}.
}
Its component along $\rollingDirection_i$ equates carrier speed and rim
speed, while its component along $\lateralDirection_i$ suppresses lateral
skidding. Thus, $\function{\text{tan},i}=\zeroVector$ is precisely the
tangential part of the ideal relation~\eqref{eq:ideal-rolling-velocity}; it
does not introduce a separate slip-velocity state.

Let the first objective sum contain the baseline posture, end-effector,
force, and balance tasks that are required in the application. The simplified
ideal-rolling controller is
\begin{equation}
  \label{eq:ideal-rolling-whole-body-qp}
  \begin{aligned}
    \min_{\decisionVar_{\text{ideal}}}\quad &
      \agg{j}{\num[obj]}{
        \weight_j\norm{\function{j}(\decisionVar_{\text{ideal}})}_2^2
      }
      \\
    \mbox{s.t.}\quad
      &\highlightblue{\text{ Equations of motion and joint bounds:}}\\
      &\quad \text{Equation~\eqref{eq:rolling-equations-of-motion} and the
        position, velocity, acceleration, and torque bounds},\\
      &\highlightgreen{\text{ Ideal rolling contacts }i\in\set{r}:}\\
      &\quad \function{\text{tan},i}(\decisionVar_{\text{ideal}})
        = \zeroVector,\\
      &\quad \transpose{\surfaceNormal_i}
        \left(
          \jacobian[][c_i]\generalizedAcceleration
          + \jacobian[dot][c_i]\generalizedVelocity
        \right)=0,\\
      &\highlightgreen{\text{ Contact-force feasibility:}}\\
      &\quad \generator{\force_i}\geq\zeroVector,
        \qquad \force_i=\cone[f][i]\generator{\force_i}.
  \end{aligned}
\end{equation}
The tangential equality in
\eqref{eq:ideal-rolling-whole-body-qp} contains both longitudinal rolling and
lateral no-skid conditions; the following scalar row maintains normal
contact. The QP neither estimates a sliding direction nor switches contact
models. This is the appropriate starting point when pure rolling is a valid
assumption and the wheel commands have already been checked for geometric
compatibility.

\subsection{Explicit planar chassis QPs}

For differential and four-steering chassis, the generic ideal formulation can
be written with explicit wheel variables. These are reduced-coordinate views
of \eqref{eq:rolling-equations-of-motion}, not independent dynamics models.
The matrices below are obtained by selecting the plane-compatible
floating-base and wheel rows from the full dynamics. Additional arm, torso,
or suspension variables can be appended to the common block without changing
the number of variables introduced by each wheel.

\subsubsection{Differential-drive QP}

For $i\in\{L,R\}$, wheel $i$ introduces exactly the rolling acceleration and
drive torque
\quickEq{eq:differential-wheel-decision-block}{
  \decisionVar^{\text{diff}}_i
  = \vectorTwo{\ddot{\wheelAngle}_i}{\jtorques_{\wheelAngle_i}}
  \in\RRv{2}.
}
Let $\bs{\xi}^{\Pi}_{\fbFrame}=\vectorThree{v_x}{v_y}{\omega_{\Pi}}$
and define
\quickEq{eq:differential-drive-decision-vector}{
  \decisionVar_{\text{diff}}
  = \begin{bmatrix}
      \dot{\bs{\xi}}^{\Pi}_{\fbFrame}\\
      \generator{\force_L}\\
      \generator{\force_R}\\
      \decisionVar^{\text{diff}}_L\\
      \decisionVar^{\text{diff}}_R
    \end{bmatrix},
  \qquad
  \bs{a}_{\text{diff}}
  = \begin{bmatrix}
      \dot{\bs{\xi}}^{\Pi}_{\fbFrame}\\
      \ddot{\wheelAngle}_L\\
      \ddot{\wheelAngle}_R
    \end{bmatrix}.
}
Here the force generators belong to the common contact model; the two
two-dimensional wheel blocks contribute four actuator variables in total.
\begin{samepage}
For constant track width and wheel radii, an explicit ideal differential-drive
QP is
\begin{equation}
  \label{eq:differential-drive-qp}
  \begin{aligned}
    \min_{\decisionVar_{\text{diff}}}\quad &
      \weight_{\xi}
      \norm{
        \bs{\xi}^{\Pi}_{\fbFrame}
        + \samplingPeriod\dot{\bs{\xi}}^{\Pi}_{\fbFrame}
        - \bs{\xi}^{\Pi,\text{ref}}_{\fbFrame}
      }_2^2\\
      &+ \sum_{i\in\{L,R\}}
      \left(
        \weight_{\wheelAngle,i}
        \squared{\ddot{\wheelAngle}_i-\reff{\ddot{\wheelAngle}}_i}
        + \weight_{\jtorques,i}\squared{\jtorques_{\wheelAngle_i}}
      \right)
      \\
    \mbox{s.t.}\quad
      &\text{M}_{\text{diff}}\bs{a}_{\text{diff}}
        + \bs{h}_{\text{diff}}
        = \text{B}_{\text{diff}}
          \vectorTwo{\jtorques_{\wheelAngle_L}}{\jtorques_{\wheelAngle_R}}
        + \sum_{i\in\{L,R\}}
          \transpose{\text{J}_{\text{diff},i}}
          \cone[f][i]\generator{\force_i},\\
      &\dot v_x-\frac{b}{2}\dot\omega_{\Pi}
        = \wheelRadius_L\ddot{\wheelAngle}_L,\\
      &\dot v_x+\frac{b}{2}\dot\omega_{\Pi}
        = \wheelRadius_R\ddot{\wheelAngle}_R,\\
      &\dot v_y=0,\\
      &\generator{\force_i}\geq\zeroVector,
        \quad \force_i=\cone[f][i]\generator{\force_i},
        \quad i\in\{L,R\},\\
      &\ddot{\wheelAngle}_{i,\text{min}}
        \leq\ddot{\wheelAngle}_i
        \leq\ddot{\wheelAngle}_{i,\text{max}},
        \quad
        \jtorques_{\wheelAngle_i,\text{min}}
        \leq\jtorques_{\wheelAngle_i}
        \leq\jtorques_{\wheelAngle_i,\text{max}}.
  \end{aligned}
\end{equation}
\end{samepage}
The three kinematic rows are the differentiated form of
\eqref{eq:differential-geometric-constraints} under the ideal zero-residual
assumption. If the measured residual is nonzero, add the proportional
stabilization term from \eqref{eq:geometric-wheel-acceleration-task} to the
right-hand side.

\subsubsection{Four-steering-wheel QP}

For $i\in\{1,\ldots,4\}$, let $\delta_i$ be the steering angle. Each wheel
introduces exactly four variables: rolling and steering accelerations, and
their two actuator torques,
\quickEq{eq:steering-wheel-decision-block}{
  \decisionVar^{\text{4s}}_i
  = \begin{bmatrix}
      \ddot{\wheelAngle}_i\\
      \ddot\delta_i\\
      \jtorques_{\wheelAngle_i}\\
      \jtorques_{\delta_i}
    \end{bmatrix}
  \in\RRv{4}.
}
The complete reduced decision and acceleration vectors are
\quickEq{eq:four-steering-decision-vector}{
  \decisionVar_{\text{4s}}
  = \begin{bmatrix}
      \dot{\bs{\xi}}^{\Pi}_{\fbFrame}\\
      \generator{\force_1}\\
      \vdots\\
      \generator{\force_4}\\
      \decisionVar^{\text{4s}}_1\\
      \vdots\\
      \decisionVar^{\text{4s}}_4
    \end{bmatrix},
  \qquad
  \bs{a}_{\text{4s}}
  = \begin{bmatrix}
      \dot{\bs{\xi}}^{\Pi}_{\fbFrame}\\
      \ddot{\wheelAngle}_1\\
      \ddot\delta_1\\
      \vdots\\
      \ddot{\wheelAngle}_4\\
      \ddot\delta_4
    \end{bmatrix}.
}
Thus, the four wheels add sixteen actuator variables. Define the actuation
vector
\quickEq{eq:four-steering-torque-vector}{
  \jtorques_{\text{4s}}
  = \begin{bmatrix}
      \jtorques_{\wheelAngle_1} & \jtorques_{\delta_1} & \cdots &
      \jtorques_{\wheelAngle_4} & \jtorques_{\delta_4}
    \end{bmatrix}^{\top}.
}
With the planar directions from \eqref{eq:steering-directions}, their known
time derivatives at the current state are
\quickEq{eq:steering-direction-derivatives}{
  \dot{\rollingDirection}^{\Pi}_i
  = \dot\delta_i\lateralDirection^{\Pi}_i,
  \qquad
  \dot{\lateralDirection}^{\Pi}_i
  = -\dot\delta_i\rollingDirection^{\Pi}_i.
}
An explicit ideal four-steering-wheel QP is
\begin{equation}
  \label{eq:four-steering-wheel-qp}
  \begin{aligned}
    \min_{\decisionVar_{\text{4s}}}\quad &
      \weight_{\xi}
      \norm{
        \bs{\xi}^{\Pi}_{\fbFrame}
        + \samplingPeriod\dot{\bs{\xi}}^{\Pi}_{\fbFrame}
        - \bs{\xi}^{\Pi,\text{ref}}_{\fbFrame}
      }_2^2\\
      &+\sum_{i=1}^{4}\left(
        \weight_{\wheelAngle,i}
        \squared{\ddot{\wheelAngle}_i-\reff{\ddot{\wheelAngle}}_i}
        +\weight_{\delta,i}\squared{\ddot\delta_i-\reff{\ddot\delta}_i}
        +\weight_{\jtorques_{\wheelAngle},i}
          \squared{\jtorques_{\wheelAngle_i}}
        +\weight_{\jtorques_{\delta},i}\squared{\jtorques_{\delta_i}}
      \right)
      \\
    \mbox{s.t.}\quad
      &\text{M}_{\text{4s}}\bs{a}_{\text{4s}}
        +\bs{h}_{\text{4s}}
        =\text{B}_{\text{4s}}\jtorques_{\text{4s}}
        +\sum_{i=1}^{4}
          \transpose{\text{J}_{\text{4s},i}}
          \cone[f][i]\generator{\force_i},\\
      &\left(\rollingDirection^{\Pi}_i\right)^\top
          \text{H}_i\dot{\bs{\xi}}^{\Pi}_{\fbFrame}
        +\left(\dot{\rollingDirection}^{\Pi}_i\right)^\top
          \text{H}_i\bs{\xi}^{\Pi}_{\fbFrame}
        =\wheelRadius_i\ddot{\wheelAngle}_i,\\
      &\left(\lateralDirection^{\Pi}_i\right)^\top
          \text{H}_i\dot{\bs{\xi}}^{\Pi}_{\fbFrame}
        +\left(\dot{\lateralDirection}^{\Pi}_i\right)^\top
          \text{H}_i\bs{\xi}^{\Pi}_{\fbFrame}
        =0,
        \qquad i=1,\ldots,4,\\
      &\generator{\force_i}\geq\zeroVector,
        \quad \force_i=\cone[f][i]\generator{\force_i},
        \qquad i=1,\ldots,4,\\
      &\ddot{\wheelAngle}_{i,\text{min}}
        \leq\ddot{\wheelAngle}_i
        \leq\ddot{\wheelAngle}_{i,\text{max}},
        \quad
        \ddot{\delta}_{i,\text{min}}
        \leq\ddot\delta_i
        \leq\ddot{\delta}_{i,\text{max}},\\
      &\jtorques_{\wheelAngle_i,\text{min}}
        \leq\jtorques_{\wheelAngle_i}
        \leq\jtorques_{\wheelAngle_i,\text{max}},
        \quad
        \jtorques_{\delta_i,\text{min}}
        \leq\jtorques_{\delta_i}
        \leq\jtorques_{\delta_i,\text{max}}.
  \end{aligned}
\end{equation}
The instantaneous rolling rows depend on the measured steering rate through
\eqref{eq:steering-direction-derivatives}; the steering acceleration
$\ddot\delta_i$ enters the dynamics, bounds, and steering objective, and
updates the direction used at the next cycle. Keeping the direction fixed
inside one solve preserves the QP structure. Normal contact rows from the
generic model must be appended to both explicit planar QPs on a flat plane or
a ramp.

\subsection{Stepwise extensions}

The simplified controller can be extended in explicit stages rather than all
at once:
\begin{enumerate}
  \item \hlBlock{Use the exact acceleration-level rows.}\ Replace the
    one-step tangential equality by
    \eqref{eq:rolling-acceleration-constraint} or its homogeneous form
    \eqref{eq:homogeneous-rolling-acceleration}. For a differential or
    four-steering chassis, use the explicit coupled rows in
    \eqref{eq:geometric-wheel-acceleration-task}; never impose both generic
    and geometric copies.
  \item \hlBlock{Protect feasibility.}\ If the longitudinal row can conflict
    with fixed contacts or wheel commands, replace the hard row by the soft
    equality~\eqref{eq:soft-rolling-equality} and penalize
    $\norm{\slackVariable_i}_2^2$. Equivalently, use the direct quadratic
    task~\eqref{eq:rolling-task} without adding the same residual a second
    time.
  \item \hlBlock{Introduce contact sets and activation.}\ Partition contacts
    into fixed, rolling, and sliding sets $\set{f}$, $\set{r}$, and
    $\set{s}$, and vary row activation continuously during transitions.
    Fixed contacts retain three zero-motion rows; rolling contacts retain
    hard normal and lateral rows plus a hard or softened longitudinal row.
  \item \hlBlock{Add sliding only when detected.}\ For $i\in\set{s}$,
    remove the tangential no-slip rows, retain the normal row, and impose the
    measured one-cycle friction direction
    \eqref{eq:sliding-friction-direction}. The residual
    \eqref{eq:slip-velocity}, force margin, and hysteresis determine this
    mode outside the QP.
  \item \hlBlock{Complete safeguards.}\ Add collision,
    CoM, ZMP, DCM, angular-momentum, transition-smoothing, and any
    application-specific constraints from the baseline controller.
\end{enumerate}

\subsection{Complete mode-aware whole-body QP}

After softening the longitudinal rows, extend the ideal decision vector with
the rolling slacks while retaining joint accelerations, torques, and
friction-cone generators:
\quickEq{eq:rolling-decision-vector}{
  \decisionVar
  = \begin{bmatrix}
      \generalizedAcceleration\\
      \jtorques\\
      \generator{\force_1}\\
      \vdots\\
      \generator{\force_{\numContacts}}\\
      \slackVariable
    \end{bmatrix}.
}
\newpage
At one control cycle, all Jacobians, direction vectors, mode activations, and
measured residuals are fixed data. Choosing the soft-equality alternative for
the longitudinal rolling rows gives the complete mode-aware whole-body QP
\begin{equation}
  \label{eq:rolling-whole-body-qp}
  \begin{aligned}
    \min_{\decisionVar}\quad &
      \agg{j}{\num[obj]}{
        \weight_j\norm{\function{j}(\decisionVar)}_2^2
      }
      + \sum_{i\in\set{r}}
        \weight_{\sigma,i}\norm{\slackVariable_i}_2^2
      + \weight_{\Delta}
        \norm{\current{\decisionVar}-\previous{\decisionVar}}_2^2
      \\
    \mbox{s.t.}\quad
      &\highlightblue{\text{ Equations of motion and joint constraints:}}\\
      &\quad \text{Equation~\eqref{eq:rolling-equations-of-motion}, joint
        position, velocity, acceleration, and torque bounds},\\
      &\highlightgreen{\text{ Fixed/sticking contacts }i\in\set{f}:}\\
      &\quad \text{three hard zero-motion acceleration rows},\\
      &\highlightgreen{\text{ Rolling contacts }i\in\set{r}:}\\
      &\quad \text{hard normal and lateral rows from
        \eqref{eq:rolling-acceleration-constraint}},\\
      &\quad \function{\text{roll},i}(\decisionVar)
        = \slackVariable_i
        \quad \text{as in \eqref{eq:soft-rolling-equality}},\\
      &\quad \text{for differential or four-steering chassis, replace the
        equivalent tangential rows by
        \eqref{eq:geometric-wheel-acceleration-task}},\\
      &\highlightgreen{\text{ Sliding contacts }i\in\set{s}:}\\
      &\quad \text{hard normal row and, for sustained sliding,
        \eqref{eq:sliding-friction-direction}},\\
      &\highlightgreen{\text{ Contact-force feasibility:}}\\
      &\quad \generator{\force_i}\geq\zeroVector,
        \quad \force_i=\cone[f][i]\generator{\force_i},\\
      &\quad \text{contact-wrench cone where a finite patch is modeled},\\
      &\highlightblue{\text{ Whole-body safety and balance:}}\\
      &\quad \text{collision, CoM, ZMP, DCM, and angular-momentum constraints}.
  \end{aligned}
\end{equation}
The complete formulation reduces to
\eqref{eq:ideal-rolling-whole-body-qp} when $\set{s}=\emptyset$, every
rolling row is imposed exactly, the slack variables and transition penalty
are removed, and the additional safeguards are omitted. Conversely, each
step above adds one identifiable modeling feature without requiring a
separate optimizer or a second contact-force solution.

\subsection{Convexity and numerical structure}

The formulation remains a QP when geometry is evaluated at the measured state
and held constant during one solve. In particular:

\begin{itemize}
  \item $\rollingMatrix$, $\dot{\rollingMatrix}$, local directions, and
    activations are QP data rather than decision-dependent expressions;
  \item the friction cone is polyhedralized through
    \eqref{eq:rolling-force-generators};
  \item the sliding direction in \eqref{eq:sliding-friction-direction} is
    measured before the solve and protected by $\varepsilon$;
  \item mode decisions are made outside the QP, while slacks preserve
    feasibility inside it;
  \item unused contact rows and force variables stay in the fixed-size problem
    with zero activation or bounded zero force, which preserves sparsity and
    warm-start structure.
\end{itemize}

If exact Coulomb complementarity or the contact mode itself becomes a decision
variable, the problem is no longer the same convex QP and should be treated as
a separate mixed-integer, complementarity, or nonlinear formulation.

\section{Solver and implementation procedure}
\label{sec:rolling-implementation}

The extension should be implemented as an additional contact type within the
existing whole-body controller rather than as a parallel kinematics solver.
This keeps floating-base dynamics, torque limits, force distribution, and
contact tasks consistent.

\subsection{Rolling-contact data}

For every wheel, store:
\begin{itemize}
  \item rolling-center origin $\point[][][\inertialFrame,c_i]$, nominal
    contact point $\point[][][\inertialFrame,p_i]$, and wheel radius
    $\wheelRadius_i$;
  \item projected base origin $\point[][][\inertialFrame,\fbFrame]$, terrain
    basis ${}^{\inertialFrame}\text{E}_{\Pi}$, its body-frame representation
    ${}^{\fbFrame}\text{E}_{\Pi}$, plane normal, and carrier offset
    $\bs{\rho}_i$;
  \item rolling direction $\rollingDirection_i$ and lateral direction
    $\lateralDirection_i$;
  \item wheel angle, rate, acceleration, and steering angle when applicable;
  \item carrier-center Jacobian $\jacobian[][c_i]$, contact-force Jacobian
    $\jacobian[][p_i]$, their derivatives, and the differentiated rolling
    blocks;
  \item friction coefficient, force generators, and torque limits;
  \item requested and estimated contact modes, activations, hysteresis state,
    and dwell time;
  \item measured slip, QP rolling slack, friction margin, and torque margin.
\end{itemize}

Sign conventions must be tested explicitly: positive wheel spin must produce
positive rim speed along $\rollingDirection_i$, and the drive-torque sign must
match the generalized-force convention in
\eqref{eq:rolling-equations-of-motion}.

\subsection{Control-cycle sequence}

At every cycle:
\begin{enumerate}
  \item Update the terrain plane, projected base origin, rolling-center
    origins, and orthonormal local frames.
  \item Project the measured floating-base body twist with
    \eqref{eq:base-plane-projection}; compute $\text{H}_i$ and the explicit
    geometric rows \eqref{eq:geometric-wheel-homogeneous-task} for a
    differential or four-steering chassis.
  \item Compute carrier/contact Jacobians, their derivatives, and the generic
    rolling matrices in \eqref{eq:rolling-block} or
    \eqref{eq:homogeneous-rolling-block}. Use either the explicit geometric
    tangential rows or their generic Jacobian equivalents, not both.
  \item Estimate \eqref{eq:slip-velocity} from the measured generalized
    velocity and wheel rate; filter it without introducing excessive delay.
  \item Update the contact mode using the requested schedule, hysteresis,
    normal force, friction margin, torque margin, and previous mode.
  \item Assemble the ideal QP~\eqref{eq:ideal-rolling-whole-body-qp} when the
    no-slip assumptions have been certified. For a planar differential or
    four-steering chassis, instantiate \eqref{eq:differential-drive-qp} or
    \eqref{eq:four-steering-wheel-qp}, respectively. Otherwise assemble the
    mode-aware QP~\eqref{eq:rolling-whole-body-qp} with the required
    activations and slacks.
  \item Warm start and solve the QP. Reject non-finite or infeasible results and
    apply the controller's existing safe fallback.
  \item Integrate the solved generalized acceleration, fuse it with the state
    estimator, and recompute the predicted slip and feasibility margins.
\end{enumerate}

\subsection{Floating-base state update}

In a no-slip planar differential-drive case, closed-form wheel odometry can
provide a useful measurement or cross-check. It must not become a second state
update that competes with the full generalized solution. On uneven terrain,
with steering wheels, or whenever slip is detected, update the floating-base
state from the complete constrained dynamics and state estimator. Ideal wheel
integration alone is then inconsistent with the observed contact motion.

For real-time performance, keep the number and ordering of decision variables
and constraints constant, preallocate all matrix blocks, exploit sparsity, and
warm start from the previous solution. Scale rolling residuals so that meters
per second, radians per second, force, torque, and acceleration objectives have
comparable numerical influence before tuning their semantic priorities.

\section{Usual rolling-contact cases}
\label{sec:rolling-cases}

The validation scope is restricted to differential-drive and
four-steering-wheel chassis on either flat terrain or a constant-slope ramp.
Each test verifies the explicit carrier--base coupling in
\eqref{eq:geometric-wheel-velocity-task}, QP feasibility, force and torque
margins, and floating-base state updates.
Under certified ideal rolling, use \eqref{eq:differential-drive-qp} for the
two-wheel chassis and \eqref{eq:four-steering-wheel-qp} for the four-wheel
chassis before enabling the mode-aware extensions.

\begin{table}[h]
  \centering
  \caption{Rolling-contact validation cases and required outcomes}
  \label{tab:rolling-validation}
  \begin{tabular}{|C{0.22\columnwidth}|C{0.30\columnwidth}|C{0.37\columnwidth}|}
    \hline
    Case & Setup & Required outcome \\
    \hline
    Differential drive---flat & Two fixed carrier offsets and independent
      wheel rates on a horizontal plane &
      The two longitudinal rows recover
      \eqref{eq:differential-linear-velocity}--
      \eqref{eq:differential-angular-velocity}; the independent lateral row
      gives $v_y=0$ \\
    \hline
    Differential drive---ramp & Same chassis on a constant-slope plane;
      express all geometry in its tangent basis &
      The flat-terrain equations hold in the ramp basis, the normal rows prevent
      penetration, and force/torque margins reflect gravity \\
    \hline
    Four steering wheels---flat & Four offsets, steering angles, and wheel
      rates on a horizontal plane &
      The stacked rows of \eqref{eq:steering-expanded-constraints} recover one
      consistent chassis twist or expose the incompatible command as rolling
      slack \\
    \hline
    Four steering wheels---ramp & Same steering geometry in a constant-slope
      tangent basis &
      All four carrier velocities match the transformed projected base twist;
      rolling, lateral, and normal residuals remain within tolerance \\
    \hline
  \end{tabular}
\end{table}

For every test, assert the algebraic residuals, not only the rendered motion.
At minimum, record $\norm{\rollingResidual_i}$,
$\norm{\slipVelocity{i}}$, $\norm{\slackVariable_i}$, normal force, friction
margin, wheel-torque margin, solver status, and solve time. Compare the
measured and predicted residuals to distinguish model error from true sliding.

\section{Future works}

The QP extension leaves several important developments open:
\begin{itemize}
  \item Add a terrain/contact estimator for uncertain normals, wheel radius,
    deformation, and moving support surfaces.
  \item Extend the point model to finite tire patches, anisotropic friction,
    rolling resistance, and compliant tire dynamics.
  \item Compare model-based slip detection with a learned detector, while
    retaining deterministic residual and feasibility monitors as safeguards.
  \item Investigate complementarity or mixed-integer contact-mode selection
    when an external mode manager is insufficient.
  \item Automate rolling-task and slack-weight tuning using normalized residuals
    and explicit priority guarantees.
  \item Expose batched rolling kinematics, slip penalties, and constraint
    residuals for IsaacLab without coupling simulation-specific code to the
    controller's core contact model.
\end{itemize}

% Bibliography
% -----------------------------------------------------------------
% This report cites only the two source notes. The pinned shared bibliography
% is intentionally not loaded because its legacy revision contains unrelated
% BibTeX syntax errors.
\bibliography{rolling-contact-qp}
\bibliographystyle{unsrtnat}

\appendix
\section{Planar wheel specializations}
\label{app:rolling-specializations}

Equations~\eqref{eq:wheel-vel-def}--
\eqref{eq:geometric-wheel-acceleration-task} are the implementation source of
truth. This appendix expands their longitudinal and lateral rows for the two
chassis types retained in Table~\ref{tab:rolling-validation}. The algebra is
unchanged between flat terrain and a ramp: only the tangent basis
${}^{\inertialFrame}\text{E}_{\Pi}$ and normal
${}^{\inertialFrame}\surfaceNormal_{\Pi}$ change.

\subsection{Differential drive}

Place $\point[][][\inertialFrame,\fbFrame]$ at the midpoint of the drive axle
projected onto $\Pi$. For track width $b$, the rolling-center offsets and
direction coordinates are
\quickEq{eq:differential-wheel-geometry}{
  \bs{\rho}_L=\vectorTwo{0}{b/2},
  \qquad
  \bs{\rho}_R=\vectorTwo{0}{-b/2},
  \qquad
  \rollingDirection_L^{\Pi}=\rollingDirection_R^{\Pi}=\vectorTwo{1}{0},
  \qquad
  \lateralDirection_L^{\Pi}=\lateralDirection_R^{\Pi}=\vectorTwo{0}{1}.
}
Substitution into \eqref{eq:planar-carrier-map} gives
\quickEq{eq:differential-carrier-velocities}{
  \contactVel_{c_L}^{\Pi}
  = \vectorTwo{v_x-\frac{b}{2}\omega_{\Pi}}{v_y},
  \qquad
  \contactVel_{c_R}^{\Pi}
  = \vectorTwo{v_x+\frac{b}{2}\omega_{\Pi}}{v_y}.
}
Hence, the independent geometric constraints imposed in the QP are
\begin{equation}
  \label{eq:differential-geometric-constraints}
  \begin{aligned}
    v_x-\frac{b}{2}\omega_{\Pi}
      &= \wheelRadius_L\dot{\wheelAngle}_L,\\
    v_x+\frac{b}{2}\omega_{\Pi}
      &= \wheelRadius_R\dot{\wheelAngle}_R,\\
    v_y &= 0.
  \end{aligned}
\end{equation}
Both wheels produce the same lateral equation, so only one independent copy
should be inserted. Keeping both copies is algebraically harmless for an exact
equality solver but introduces a redundant row and can worsen conditioning.

Solving the first two rows of
\eqref{eq:differential-geometric-constraints} yields the familiar diagnostic
relations
\begin{align}
  v_x
  &= \frac{
      \wheelRadius_L\dot{\wheelAngle}_L
      + \wheelRadius_R\dot{\wheelAngle}_R
    }{2},
  \label{eq:differential-linear-velocity}\\
  \omega_{\Pi}
  &= \frac{
      \wheelRadius_R\dot{\wheelAngle}_R
      - \wheelRadius_L\dot{\wheelAngle}_L
    }{b}.
  \label{eq:differential-angular-velocity}
\end{align}
The associated inertial-frame motion on either a flat plane or a ramp is
\quickEq{eq:differential-pose-rate}{
  \point[dot][][\inertialFrame,\fbFrame]
  = {}^{\inertialFrame}\text{E}_{\Pi}\vectorTwo{v_x}{v_y},
  \qquad
  {}^{\inertialFrame}\bs{\omega}_{\fbFrame}
  = \omega_{\Pi}{}^{\inertialFrame}\surfaceNormal_{\Pi}.
}
These closed-form equations are validation identities. The controller imposes
the differentiated rows of
\eqref{eq:differential-geometric-constraints} through
\eqref{eq:geometric-wheel-acceleration-task}, together with dynamics, normal
contact, friction, and torque limits.

\subsection{Four steering wheels}

For steering angle $\delta_i$ measured in the tangent basis of $\Pi$, define
\quickEq{eq:steering-directions}{
  \rollingDirection_i^{\Pi}
  = \vectorTwo{\cos\delta_i}{\sin\delta_i},
  \qquad
  \lateralDirection_i^{\Pi}
  = \vectorTwo{-\sin\delta_i}{\cos\delta_i}.
}
For $\bs{\rho}_i=\transpose{[x_i,\,y_i]}$, the transformed carrier velocity
from \eqref{eq:planar-carrier-map} is
\quickEq{eq:steering-carrier-velocity}{
  \contactVel_{c_i}^{\Pi}
  = \vectorTwo{
      v_x-\omega_{\Pi} y_i
    }{
      v_y+\omega_{\Pi} x_i
    }.
}
Equation~\eqref{eq:geometric-wheel-velocity-task} therefore expands to
\begin{equation}
  \label{eq:steering-expanded-constraints}
  \begin{aligned}
    \cos\delta_i\left(v_x-\omega_{\Pi}y_i\right)
    +\sin\delta_i\left(v_y+\omega_{\Pi}x_i\right)
      &= \wheelRadius_i\dot{\wheelAngle}_i,\\
    -\sin\delta_i\left(v_x-\omega_{\Pi}y_i\right)
    +\cos\delta_i\left(v_y+\omega_{\Pi}x_i\right)
      &=0.
  \end{aligned}
\end{equation}
Stacking \eqref{eq:steering-expanded-constraints} for all four wheels ties
their rates and steering angles to one projected floating-base twist. If the
commands are mutually incompatible, the rolling rows use the slack of
\eqref{eq:soft-rolling-equality}; the optimizer must not compute four
independent chassis velocities. When steering changes during motion,
$\dot{\text{G}}^{\text{geo}}_i$ in
\eqref{eq:geometric-wheel-acceleration-task} includes
$\dot{\delta}_i$. On a ramp, the same equations hold after expressing
$\bs{\rho}_i$, $\rollingDirection_i$, and $\lateralDirection_i$ in the ramp
tangent basis.

\section{Source provenance and corrected assumptions}
\label{app:rolling-sources}

The initial note~\cite{rollingContactNotion} supplied the design questions:
sliding detection, floating-base updates, coupled differential and steering
wheels, friction direction, actuation limits, conflicting coplanar contacts,
and tests on ramps and leg-wheeled systems. The completed
note~\cite{completeRollingContactsNotion} supplied the local-frame equations,
stacked velocity solve, friction inequalities, planar models, mode logic, and
implementation phases.

Before placing that material in the baseline QP, this report defines the
rolling-center origin explicitly, uses $\inertialFrame$ consistently for the
inertial frame, and distinguishes carrier- from material-point Jacobians. It
also restores the slip equalities, differentiates the velocity relation,
separates sticking from saturated sliding friction, and replaces the invalid
torque--radius relation with complete dynamics and
\eqref{eq:quasistatic-wheel-force}. Finally, the explicit carrier--base task
\eqref{eq:geometric-wheel-acceleration-task} couples every differential or
steering wheel to one projected floating-base twist; slacks and mode
transitions handle incompatible commands without a separate pseudoinverse.

\end{document}
